\documentclass[12pt, ukrainian]{article}
\usepackage[a4paper]{geometry}
\setlength\parindent{0mm}
%\setlength\parskip{4mm}
\geometry{top=1.3cm,bottom=1.3cm,left=1.5cm,right=1.5cm}
\pagestyle{empty}
\usepackage[T2A]{fontenc}
\usepackage[utf8]{inputenc}
\usepackage[ukrainian]{babel}
\begin{document}
%begin
{{\begin {scriptsize}\par
{ \center  Київський національний університет імені Тараса Шевченка \par}
{\parbox {5 cm}{Освітня програма \hfill}{\parbox {7 cm}{Інформатика\hfill}}\parbox {6 cm}{\hfill Семестр 1}}\par
{\parbox {5 cm} {Навчальний предмет \hfill}\parbox {7 cm}{Програмування \hfill}\parbox {6cm}{\hfill}\par}\vspace{-15pt} \end{scriptsize}}{\center Екзаменаційний білет \No \textbf { 1 }\par}}
{%beginitem
1. Компільовані та інтерпретовані мови. Переваги та недоліки.%enditem
\par
%beginitem
2. Написати програму, що запитує у користувача значення дійсних параметрів $x$ та $y$, обчислює для них значення виразу 
$$\frac{\ln x - 30}{(\arctg x + 0.9) (y + 77)}$$ 
та повiдомляє введенi данi та результат обчислення.

%enditem
\par
%beginitem
3. Написати функцію, що для інтервалу $[a, b)$ знаходить найближчу пару чисел з цього інтервалу, що кожне з них має парну кількість різних простих дільників. Якщо таких пар кілька, то цікавить пара з найменшим добутком.

По завданням 2 та 3 оцінюється правильність та якість коду, а також економність обчислень.

%enditem
}
{\smallskip\smallskip \begin {scriptsize} Затверджено на засіданні кафедри теоретичної кібернетики, протокол \No 2  від   07.10.2024 р.\par\smallskip\smallskip\smallskip\smallskip
{\parbox {6.5 cm} {Зав. кафедри \dotfill Ю.В.~Крак}}\hspace {0.7cm} {\hbox to 6.5 cm { Екзаменатор \dotfill Т.О.~Карнаух}} \end{scriptsize}}
%end
\vspace{0.9cm}
%begin
{{\begin {scriptsize}\par
{ \center  Київський національний університет імені Тараса Шевченка \par}
{\parbox {5 cm}{Освітня програма \hfill}{\parbox {7 cm}{Інформатика\hfill}}\parbox {6 cm}{\hfill Семестр 1}}\par
{\parbox {5 cm} {Навчальний предмет \hfill}\parbox {7 cm}{Програмування \hfill}\parbox {6cm}{\hfill}\par}\vspace{-15pt} \end{scriptsize}}{\center Екзаменаційний білет \No \textbf { 2 }\par}}
{%beginitem
1. Математичний підхід до визначення складності обчислень.%enditem
\par
%beginitem
2. Написати програму, що запитує у користувача значення дійсних параметрів $x$ та $y$, обчислює для них значення виразу 
$$\frac{\pi x + 98}{(\sin x - 0.2) (y - 21)}$$ 
та повiдомляє введенi данi та результат обчислення.

%enditem
\par
%beginitem
3. Написати функцію, що для послідовності цілих (задається масивом) знаходить найдовший відрізок, в якому всі сусідні числа відрізняються якнайменш на 3. Якщо таких відрізків кілька, то повернути перший (як індекс початку та довжину). Використовувати додаткові структури даних (у тому числі рядки) заборонено.

По завданням 2 та 3 оцінюється правильність та якість коду, а також економність обчислень.

%enditem
}
{\smallskip\smallskip \begin {scriptsize} Затверджено на засіданні кафедри теоретичної кібернетики, протокол \No 2  від   07.10.2024 р.\par\smallskip\smallskip\smallskip\smallskip
{\parbox {6.5 cm} {Зав. кафедри \dotfill Ю.В.~Крак}}\hspace {0.7cm} {\hbox to 6.5 cm { Екзаменатор \dotfill Т.О.~Карнаух}} \end{scriptsize}}
%end
\vspace{0.9cm}
%begin
{{\begin {scriptsize}\par
{ \center  Київський національний університет імені Тараса Шевченка \par}
{\parbox {5 cm}{Освітня програма \hfill}{\parbox {7 cm}{Інформатика\hfill}}\parbox {6 cm}{\hfill Семестр 1}}\par
{\parbox {5 cm} {Навчальний предмет \hfill}\parbox {7 cm}{Програмування \hfill}\parbox {6cm}{\hfill}\par}\vspace{-15pt} \end{scriptsize}}{\center Екзаменаційний білет \No \textbf { 3 }\par}}
{%beginitem
1. Оператори порівняння (relational), рівності (equality) та булеві оператори.%enditem
\par
%beginitem
2. Написати програму, що запитує у користувача значення дійсних параметрів $x$ та $y$, обчислює для них значення виразу 
$$\frac{\sqrt x + 92}{(\cos x + 0.1) (y + 58)}$$ 
та повiдомляє введенi данi та результат обчислення.

%enditem
\par
%beginitem
3. Написати функцію, що в інтервалі $[a, b)$ знаходить ціле число, що ділиться на найбільший степінь числа 3. Якщо таких чисел декілька, повернути найбільше.

По завданням 2 та 3 оцінюється правильність та якість коду, а також економність обчислень.

%enditem
}
{\smallskip\smallskip \begin {scriptsize} Затверджено на засіданні кафедри теоретичної кібернетики, протокол \No 2  від   07.10.2024 р.\par\smallskip\smallskip\smallskip\smallskip
{\parbox {6.5 cm} {Зав. кафедри \dotfill Ю.В.~Крак}}\hspace {0.7cm} {\hbox to 6.5 cm { Екзаменатор \dotfill Т.О.~Карнаух}} \end{scriptsize}}
%end
\newpage
%begin
{{\begin {scriptsize}\par
{ \center  Київський національний університет імені Тараса Шевченка \par}
{\parbox {5 cm}{Освітня програма \hfill}{\parbox {7 cm}{Інформатика\hfill}}\parbox {6 cm}{\hfill Семестр 1}}\par
{\parbox {5 cm} {Навчальний предмет \hfill}\parbox {7 cm}{Програмування \hfill}\parbox {6cm}{\hfill}\par}\vspace{-15pt} \end{scriptsize}}{\center Екзаменаційний білет \No \textbf { 4 }\par}}
{%beginitem
1. Семантика функції та її виклику, параметри та аргументи.%enditem
\par
%beginitem
2. Написати програму, що запитує у користувача значення дійсних параметрів $x$ та $y$, обчислює для них значення виразу 
$$\frac{\log_{10} x - 59}{(\ctg x - 0.6) (y - 39)}$$ 
та повiдомляє введенi данi та результат обчислення.

%enditem
\par
%beginitem
3. Написати функцію, що повертає число з послідовності цілих (задається масивом), що має найменшу кількість різних простих дільників. Якщо таких чисел кілька, то повернути найбільше. Використовувати додаткові структури даних (у тому числі рядки) заборонено.

По завданням 2 та 3 оцінюється правильність та якість коду, а також економність обчислень.

%enditem
}
{\smallskip\smallskip \begin {scriptsize} Затверджено на засіданні кафедри теоретичної кібернетики, протокол \No 2  від   07.10.2024 р.\par\smallskip\smallskip\smallskip\smallskip
{\parbox {6.5 cm} {Зав. кафедри \dotfill Ю.В.~Крак}}\hspace {0.7cm} {\hbox to 6.5 cm { Екзаменатор \dotfill Т.О.~Карнаух}} \end{scriptsize}}
%end
\vspace{0.9cm}
%begin
{{\begin {scriptsize}\par
{ \center  Київський національний університет імені Тараса Шевченка \par}
{\parbox {5 cm}{Освітня програма \hfill}{\parbox {7 cm}{Інформатика\hfill}}\parbox {6 cm}{\hfill Семестр 1}}\par
{\parbox {5 cm} {Навчальний предмет \hfill}\parbox {7 cm}{Програмування \hfill}\parbox {6cm}{\hfill}\par}\vspace{-15pt} \end{scriptsize}}{\center Екзаменаційний білет \No \textbf { 5 }\par}}
{%beginitem
1. Винятки та їх оброблення.%enditem
\par
%beginitem
2. Написати програму, що запитує у користувача значення дійсних параметрів $x$ та $y$, обчислює для них значення виразу 
$$\frac{\ x + 49}{(\tg x + 0.5) (y - 20)}$$ 
та повiдомляє введенi данi та результат обчислення.

%enditem
\par
%beginitem
3. Написати функцію, що для цілого числа знаходить простий дільник, який входить у його розклад у найбільшому степені. Якщо таких кілька, то повернути найменший. Наприклад, для числа $2^{9}{\cdot}3^5{\cdot}7^9{\cdot}11$ функція має повернути $2$.

По завданням 2 та 3 оцінюється правильність та якість коду, а також економність обчислень.

%enditem
}
{\smallskip\smallskip \begin {scriptsize} Затверджено на засіданні кафедри теоретичної кібернетики, протокол \No 2  від   07.10.2024 р.\par\smallskip\smallskip\smallskip\smallskip
{\parbox {6.5 cm} {Зав. кафедри \dotfill Ю.В.~Крак}}\hspace {0.7cm} {\hbox to 6.5 cm { Екзаменатор \dotfill Т.О.~Карнаух}} \end{scriptsize}}
%end
\vspace{0.9cm}
%begin
{{\begin {scriptsize}\par
{ \center  Київський національний університет імені Тараса Шевченка \par}
{\parbox {5 cm}{Освітня програма \hfill}{\parbox {7 cm}{Інформатика\hfill}}\parbox {6 cm}{\hfill Семестр 1}}\par
{\parbox {5 cm} {Навчальний предмет \hfill}\parbox {7 cm}{Програмування \hfill}\parbox {6cm}{\hfill}\par}\vspace{-15pt} \end{scriptsize}}{\center Екзаменаційний білет \No \textbf { 6 }\par}}
{%beginitem
1. Засоби керування порядком обчислень. Інструкція if.%enditem
\par
%beginitem
2. Написати програму, що запитує у користувача значення дійсних параметрів $x$ та $y$, обчислює для них значення виразу 
$$\frac{\log_{2} x - 71}{(\arcsin x - 0.7) (y + 69)}$$ 
та повiдомляє введенi данi та результат обчислення.

%enditem
\par
%beginitem
3. Написати функцію, що для інтервалу $[a, b)$ знаходить найближчу пару чисел з цього інтервалу, що кожне з них має непарну кількість різних простих дільників. Якщо таких пар кілька, то цікавить пара з найбільшою сумою.

По завданням 2 та 3 оцінюється правильність та якість коду, а також економність обчислень.

%enditem
}
{\smallskip\smallskip \begin {scriptsize} Затверджено на засіданні кафедри теоретичної кібернетики, протокол \No 2  від   07.10.2024 р.\par\smallskip\smallskip\smallskip\smallskip
{\parbox {6.5 cm} {Зав. кафедри \dotfill Ю.В.~Крак}}\hspace {0.7cm} {\hbox to 6.5 cm { Екзаменатор \dotfill Т.О.~Карнаух}} \end{scriptsize}}
%end
\newpage
%begin
{{\begin {scriptsize}\par
{ \center  Київський національний університет імені Тараса Шевченка \par}
{\parbox {5 cm}{Освітня програма \hfill}{\parbox {7 cm}{Інформатика\hfill}}\parbox {6 cm}{\hfill Семестр 1}}\par
{\parbox {5 cm} {Навчальний предмет \hfill}\parbox {7 cm}{Програмування \hfill}\parbox {6cm}{\hfill}\par}\vspace{-15pt} \end{scriptsize}}{\center Екзаменаційний білет \No \textbf { 7 }\par}}
{%beginitem
1. Цикли while та do-while, їх вигляд і семантика, умова продовження циклу.%enditem
\par
%beginitem
2. Написати програму, що запитує у користувача значення дійсних параметрів $x$ та $y$, обчислює для них значення виразу 
$$\frac{\log_{10} x + 57}{(\arccos x + 0.3) (y + 87)}$$ 
та повiдомляє введенi данi та результат обчислення.

%enditem
\par
%beginitem
3. Написати функцію, що для послідовності цілих (задається масивом) знаходить довжину найдовшого відрізку, що складається з чисел, у сімковому запису якого нема цифри 2. За відсутності має повертати $0$. Використовувати додаткові структури даних (у тому числі рядки) заборонено.

По завданням 2 та 3 оцінюється правильність та якість коду, а також економність обчислень.

%enditem
}
{\smallskip\smallskip \begin {scriptsize} Затверджено на засіданні кафедри теоретичної кібернетики, протокол \No 2  від   07.10.2024 р.\par\smallskip\smallskip\smallskip\smallskip
{\parbox {6.5 cm} {Зав. кафедри \dotfill Ю.В.~Крак}}\hspace {0.7cm} {\hbox to 6.5 cm { Екзаменатор \dotfill Т.О.~Карнаух}} \end{scriptsize}}
%end
\vspace{0.9cm}
%begin
{{\begin {scriptsize}\par
{ \center  Київський національний університет імені Тараса Шевченка \par}
{\parbox {5 cm}{Освітня програма \hfill}{\parbox {7 cm}{Інформатика\hfill}}\parbox {6 cm}{\hfill Семестр 1}}\par
{\parbox {5 cm} {Навчальний предмет \hfill}\parbox {7 cm}{Програмування \hfill}\parbox {6cm}{\hfill}\par}\vspace{-15pt} \end{scriptsize}}{\center Екзаменаційний білет \No \textbf { 8 }\par}}
{%beginitem
1. Глибина рекурсії та загальна кількість рекурсивних викликів, їх вплив на розміри програмного стеку та на час виконання програми.%enditem
\par
%beginitem
2. Написати програму, що запитує у користувача значення дійсних параметрів $x$ та $y$, обчислює для них значення виразу 
$$\frac{\log_{2} x - 12}{(\arctg x - 0.8) (y - 63)}$$ 
та повiдомляє введенi данi та результат обчислення.

%enditem
\par
%beginitem
3. Написати функцію, що повертає найбільше число з послідовності цілих (задається масивом), що має рівно три різних простих дільники. За відсутності має повертати None. Використовувати додаткові структури даних (у тому числі рядки) заборонено.

По завданням 2 та 3 оцінюється правильність та якість коду, а також економність обчислень.

%enditem
}
{\smallskip\smallskip \begin {scriptsize} Затверджено на засіданні кафедри теоретичної кібернетики, протокол \No 2  від   07.10.2024 р.\par\smallskip\smallskip\smallskip\smallskip
{\parbox {6.5 cm} {Зав. кафедри \dotfill Ю.В.~Крак}}\hspace {0.7cm} {\hbox to 6.5 cm { Екзаменатор \dotfill Т.О.~Карнаух}} \end{scriptsize}}
%end
\vspace{0.9cm}
%begin
{{\begin {scriptsize}\par
{ \center  Київський національний університет імені Тараса Шевченка \par}
{\parbox {5 cm}{Освітня програма \hfill}{\parbox {7 cm}{Інформатика\hfill}}\parbox {6 cm}{\hfill Семестр 1}}\par
{\parbox {5 cm} {Навчальний предмет \hfill}\parbox {7 cm}{Програмування \hfill}\parbox {6cm}{\hfill}\par}\vspace{-15pt} \end{scriptsize}}{\center Екзаменаційний білет \No \textbf { 9 }\par}}
{%beginitem
1. Поняття типу даних. Вбудовані арифметичні та цілі (intergral) типи мови С++.%enditem
\par
%beginitem
2. Написати програму, що запитує у користувача значення дійсних параметрів $x$ та $y$, обчислює для них значення виразу 
$$\frac{\pi x - 25}{(\sin x + 0.4) (y - 61)}$$ 
та повiдомляє введенi данi та результат обчислення.

%enditem
\par
%beginitem
3. Написати функцію, що повертає число з послідовності цілих (задається масивом), що має найбільшу кількість різних простих дільників. Якщо таких чисел кілька, то повернути найменше. Використовувати додаткові структури даних (у тому числі рядки) заборонено.

По завданням 2 та 3 оцінюється правильність та якість коду, а також економність обчислень.

%enditem
}
{\smallskip\smallskip \begin {scriptsize} Затверджено на засіданні кафедри теоретичної кібернетики, протокол \No 2  від   07.10.2024 р.\par\smallskip\smallskip\smallskip\smallskip
{\parbox {6.5 cm} {Зав. кафедри \dotfill Ю.В.~Крак}}\hspace {0.7cm} {\hbox to 6.5 cm { Екзаменатор \dotfill Т.О.~Карнаух}} \end{scriptsize}}
%end
\newpage
%begin
{{\begin {scriptsize}\par
{ \center  Київський національний університет імені Тараса Шевченка \par}
{\parbox {5 cm}{Освітня програма \hfill}{\parbox {7 cm}{Інформатика\hfill}}\parbox {6 cm}{\hfill Семестр 1}}\par
{\parbox {5 cm} {Навчальний предмет \hfill}\parbox {7 cm}{Програмування \hfill}\parbox {6cm}{\hfill}\par}\vspace{-15pt} \end{scriptsize}}{\center Екзаменаційний білет \No \textbf { 10 }\par}}
{%beginitem
1. Арифметичні операції. Вирази. Пріоритети операторів. Порядок обчислення виразів.%enditem
\par
%beginitem
2. Написати програму, що запитує у користувача значення дійсних параметрів $x$ та $y$, обчислює для них значення виразу 
$$\frac{\ln x + 90}{(\arcsin x - 0.9) (y + 84)}$$ 
та повiдомляє введенi данi та результат обчислення.

%enditem
\par
%beginitem
3. Написати функцію, що знаходить найбільше число з послідовності цілих (задається масивом), у десятковому запису якого нема цифри 3. Повертає індекс останнього входження цього числа або $-1$ (за відсутності). Використовувати додаткові структури даних (у тому числі рядки) заборонено.

По завданням 2 та 3 оцінюється правильність та якість коду, а також економність обчислень.

%enditem
}
{\smallskip\smallskip \begin {scriptsize} Затверджено на засіданні кафедри теоретичної кібернетики, протокол \No 2  від   07.10.2024 р.\par\smallskip\smallskip\smallskip\smallskip
{\parbox {6.5 cm} {Зав. кафедри \dotfill Ю.В.~Крак}}\hspace {0.7cm} {\hbox to 6.5 cm { Екзаменатор \dotfill Т.О.~Карнаух}} \end{scriptsize}}
%end
\vspace{0.9cm}
%begin
{{\begin {scriptsize}\par
{ \center  Київський національний університет імені Тараса Шевченка \par}
{\parbox {5 cm}{Освітня програма \hfill}{\parbox {7 cm}{Інформатика\hfill}}\parbox {6 cm}{\hfill Семестр 1}}\par
{\parbox {5 cm} {Навчальний предмет \hfill}\parbox {7 cm}{Програмування \hfill}\parbox {6cm}{\hfill}\par}\vspace{-15pt} \end{scriptsize}}{\center Екзаменаційний білет \No \textbf { 11 }\par}}
{%beginitem
1. Цикл for, вигляд і семантика інструкції циклу for.%enditem
\par
%beginitem
2. Написати програму, що запитує у користувача значення дійсних параметрів $x$ та $y$, обчислює для них значення виразу 
$$\frac{\ x + 60}{(\arccos x + 0.8) (y - 46)}$$ 
та повiдомляє введенi данi та результат обчислення.

%enditem
\par
%beginitem
3. Написати функцію, що повертає найбільше число з послідовності цілих (задається масивом), у сімковому запису якого нема цифри 2. За відсутності має повертати None. Використовувати додаткові структури даних (у тому числі рядки) заборонено.

По завданням 2 та 3 оцінюється правильність та якість коду, а також економність обчислень.

%enditem
}
{\smallskip\smallskip \begin {scriptsize} Затверджено на засіданні кафедри теоретичної кібернетики, протокол \No 2  від   07.10.2024 р.\par\smallskip\smallskip\smallskip\smallskip
{\parbox {6.5 cm} {Зав. кафедри \dotfill Ю.В.~Крак}}\hspace {0.7cm} {\hbox to 6.5 cm { Екзаменатор \dotfill Т.О.~Карнаух}} \end{scriptsize}}
%end
\vspace{0.9cm}
%begin
{{\begin {scriptsize}\par
{ \center  Київський національний університет імені Тараса Шевченка \par}
{\parbox {5 cm}{Освітня програма \hfill}{\parbox {7 cm}{Інформатика\hfill}}\parbox {6 cm}{\hfill Семестр 1}}\par
{\parbox {5 cm} {Навчальний предмет \hfill}\parbox {7 cm}{Програмування \hfill}\parbox {6cm}{\hfill}\par}\vspace{-15pt} \end{scriptsize}}{\center Екзаменаційний білет \No \textbf { 12 }\par}}
{%beginitem
1. Поняття часової та просторової складності обчислень. Складність у середньому. Асимптотичні оцінки складності.%enditem
\par
%beginitem
2. Написати програму, що запитує у користувача значення дійсних параметрів $x$ та $y$, обчислює для них значення виразу 
$$\frac{\sqrt x - 67}{(\tg x - 0.1) (y + 43)}$$ 
та повiдомляє введенi данi та результат обчислення.

%enditem
\par
%beginitem
3. Написати функцію, що знаходить найменше число з послідовності цілих (задається масивом), у десятковому запису якого нема цифри 5. Повертає індекс останнього входження цього числа або $-1$ (за відсутності). Використовувати додаткові структури даних (у тому числі рядки) заборонено.

По завданням 2 та 3 оцінюється правильність та якість коду, а також економність обчислень.

%enditem
}
{\smallskip\smallskip \begin {scriptsize} Затверджено на засіданні кафедри теоретичної кібернетики, протокол \No 2  від   07.10.2024 р.\par\smallskip\smallskip\smallskip\smallskip
{\parbox {6.5 cm} {Зав. кафедри \dotfill Ю.В.~Крак}}\hspace {0.7cm} {\hbox to 6.5 cm { Екзаменатор \dotfill Т.О.~Карнаух}} \end{scriptsize}}
%end
\newpage
%begin
{{\begin {scriptsize}\par
{ \center  Київський національний університет імені Тараса Шевченка \par}
{\parbox {5 cm}{Освітня програма \hfill}{\parbox {7 cm}{Інформатика\hfill}}\parbox {6 cm}{\hfill Семестр 1}}\par
{\parbox {5 cm} {Навчальний предмет \hfill}\parbox {7 cm}{Програмування \hfill}\parbox {6cm}{\hfill}\par}\vspace{-15pt} \end{scriptsize}}{\center Екзаменаційний білет \No \textbf { 13 }\par}}
{%beginitem
1. Рекурсивні означення, рекурсивні функції.%enditem
\par
%beginitem
2. Написати програму, що запитує у користувача значення дійсних параметрів $x$ та $y$, обчислює для них значення виразу 
$$\frac{\log_{10} x - 18}{(\ctg x - 0.7) (y + 56)}$$ 
та повiдомляє введенi данi та результат обчислення.

%enditem
\par
%beginitem
3. Написати функцію, що для інтервалу $[a, b)$ знаходить найближчу пару чисел з цього інтервалу, що кожне з них має парну кількість різних простих дільників. Якщо таких пар кілька, то цікавить пара з найбільшою сумою.

По завданням 2 та 3 оцінюється правильність та якість коду, а також економність обчислень.

%enditem
}
{\smallskip\smallskip \begin {scriptsize} Затверджено на засіданні кафедри теоретичної кібернетики, протокол \No 2  від   07.10.2024 р.\par\smallskip\smallskip\smallskip\smallskip
{\parbox {6.5 cm} {Зав. кафедри \dotfill Ю.В.~Крак}}\hspace {0.7cm} {\hbox to 6.5 cm { Екзаменатор \dotfill Т.О.~Карнаух}} \end{scriptsize}}
%end
\vspace{0.9cm}
%begin
{{\begin {scriptsize}\par
{ \center  Київський національний університет імені Тараса Шевченка \par}
{\parbox {5 cm}{Освітня програма \hfill}{\parbox {7 cm}{Інформатика\hfill}}\parbox {6 cm}{\hfill Семестр 1}}\par
{\parbox {5 cm} {Навчальний предмет \hfill}\parbox {7 cm}{Програмування \hfill}\parbox {6cm}{\hfill}\par}\vspace{-15pt} \end{scriptsize}}{\center Екзаменаційний білет \No \textbf { 14 }\par}}
{%beginitem
1. Алгоритм двійкового пошуку та його складність.%enditem
\par
%beginitem
2. Написати програму, що запитує у користувача значення дійсних параметрів $x$ та $y$, обчислює для них значення виразу 
$$\frac{\log_{2} x + 26}{(\cos x + 0.3) (y - 93)}$$ 
та повiдомляє введенi данi та результат обчислення.

%enditem
\par
%beginitem
3. Написати функцію, що для інтервалу $[a, b)$ знаходить кількість чисел, квадрати яких діляться на суму квадратів їх цифр. (Розглядаємо десятковий запис.)

По завданням 2 та 3 оцінюється правильність та якість коду, а також економність обчислень.

%enditem
}
{\smallskip\smallskip \begin {scriptsize} Затверджено на засіданні кафедри теоретичної кібернетики, протокол \No 2  від   07.10.2024 р.\par\smallskip\smallskip\smallskip\smallskip
{\parbox {6.5 cm} {Зав. кафедри \dotfill Ю.В.~Крак}}\hspace {0.7cm} {\hbox to 6.5 cm { Екзаменатор \dotfill Т.О.~Карнаух}} \end{scriptsize}}
%end
\vspace{0.9cm}
%begin
{{\begin {scriptsize}\par
{ \center  Київський національний університет імені Тараса Шевченка \par}
{\parbox {5 cm}{Освітня програма \hfill}{\parbox {7 cm}{Інформатика\hfill}}\parbox {6 cm}{\hfill Семестр 1}}\par
{\parbox {5 cm} {Навчальний предмет \hfill}\parbox {7 cm}{Програмування \hfill}\parbox {6cm}{\hfill}\par}\vspace{-15pt} \end{scriptsize}}{\center Екзаменаційний білет \No \textbf { 15 }\par}}
{%beginitem
1. Змінні, їх властивості (у тому числі час існування, область дії оголошення імені) та операції над ними.%enditem
\par
%beginitem
2. Написати програму, що запитує у користувача значення дійсних параметрів $x$ та $y$, обчислює для них значення виразу 
$$\frac{\ln x + 73}{(\ctg x - 0.4) (y + 41)}$$ 
та повiдомляє введенi данi та результат обчислення.

%enditem
\par
%beginitem
3. Написати функцію, що повертає найменше число з послідовності цілих (задається масивом), що має не більше трьох різних простих дільників. За відсутності має повертати None. Використовувати додаткові структури даних (у тому числі рядки) заборонено.

По завданням 2 та 3 оцінюється правильність та якість коду, а також економність обчислень.

%enditem
}
{\smallskip\smallskip \begin {scriptsize} Затверджено на засіданні кафедри теоретичної кібернетики, протокол \No 2  від   07.10.2024 р.\par\smallskip\smallskip\smallskip\smallskip
{\parbox {6.5 cm} {Зав. кафедри \dotfill Ю.В.~Крак}}\hspace {0.7cm} {\hbox to 6.5 cm { Екзаменатор \dotfill Т.О.~Карнаух}} \end{scriptsize}}
%end
\newpage
%begin
{{\begin {scriptsize}\par
{ \center  Київський національний університет імені Тараса Шевченка \par}
{\parbox {5 cm}{Освітня програма \hfill}{\parbox {7 cm}{Інформатика\hfill}}\parbox {6 cm}{\hfill Семестр 1}}\par
{\parbox {5 cm} {Навчальний предмет \hfill}\parbox {7 cm}{Програмування \hfill}\parbox {6cm}{\hfill}\par}\vspace{-15pt} \end{scriptsize}}{\center Екзаменаційний білет \No \textbf { 16 }\par}}
{%beginitem
1. Емпіричний підхід до визначення складності обчислень.%enditem
\par
%beginitem
2. Написати програму, що запитує у користувача значення дійсних параметрів $x$ та $y$, обчислює для них значення виразу 
$$\frac{\ x - 55}{(\tg x + 0.5) (y - 50)}$$ 
та повiдомляє введенi данi та результат обчислення.

%enditem
\par
%beginitem
3. Написати функцію, що знаходить найменше число з послідовності цілих (задається масивом), у десятковому запису якого нема цифри 7. Повертає індекс першого входження цього числа або $-1$ (за відсутності). Використовувати додаткові структури даних (у тому числі рядки) заборонено.

По завданням 2 та 3 оцінюється правильність та якість коду, а також економність обчислень.

%enditem
}
{\smallskip\smallskip \begin {scriptsize} Затверджено на засіданні кафедри теоретичної кібернетики, протокол \No 2  від   07.10.2024 р.\par\smallskip\smallskip\smallskip\smallskip
{\parbox {6.5 cm} {Зав. кафедри \dotfill Ю.В.~Крак}}\hspace {0.7cm} {\hbox to 6.5 cm { Екзаменатор \dotfill Т.О.~Карнаух}} \end{scriptsize}}
%end
\vspace{0.9cm}
%begin
{{\begin {scriptsize}\par
{ \center  Київський національний університет імені Тараса Шевченка \par}
{\parbox {5 cm}{Освітня програма \hfill}{\parbox {7 cm}{Інформатика\hfill}}\parbox {6 cm}{\hfill Семестр 1}}\par
{\parbox {5 cm} {Навчальний предмет \hfill}\parbox {7 cm}{Програмування \hfill}\parbox {6cm}{\hfill}\par}\vspace{-15pt} \end{scriptsize}}{\center Екзаменаційний білет \No \textbf { 17 }\par}}
{%beginitem
1. Поняття масиву. Моделювання послідовностей масивами. Параметри, що зображують масиви у функціях.%enditem
\par
%beginitem
2. Написати програму, що запитує у користувача значення дійсних параметрів $x$ та $y$, обчислює для них значення виразу 
$$\frac{\pi x - 82}{(\sin x + 0.2) (y + 79)}$$ 
та повiдомляє введенi данi та результат обчислення.

%enditem
\par
%beginitem
3. Написати функцію, що для цілого числа знаходить кількість його простих дільників, що входять у його розклад не більш ніж у 5му степеню. Наприклад, для числа $3^8{\cdot}7^9{\cdot}11{\cdot}23^{2}$ таких дільників два – $11$ та $23$, тому функція має повернути $2$.

По завданням 2 та 3 оцінюється правильність та якість коду, а також економність обчислень.

%enditem
}
{\smallskip\smallskip \begin {scriptsize} Затверджено на засіданні кафедри теоретичної кібернетики, протокол \No 2  від   07.10.2024 р.\par\smallskip\smallskip\smallskip\smallskip
{\parbox {6.5 cm} {Зав. кафедри \dotfill Ю.В.~Крак}}\hspace {0.7cm} {\hbox to 6.5 cm { Екзаменатор \dotfill Т.О.~Карнаух}} \end{scriptsize}}
%end
\vspace{0.9cm}
%begin
{{\begin {scriptsize}\par
{ \center  Київський національний університет імені Тараса Шевченка \par}
{\parbox {5 cm}{Освітня програма \hfill}{\parbox {7 cm}{Інформатика\hfill}}\parbox {6 cm}{\hfill Семестр 1}}\par
{\parbox {5 cm} {Навчальний предмет \hfill}\parbox {7 cm}{Програмування \hfill}\parbox {6cm}{\hfill}\par}\vspace{-15pt} \end{scriptsize}}{\center Екзаменаційний білет \No \textbf { 18 }\par}}
{%beginitem
1. Принципи зображення цілих та дійсних чисел у комп'ютері.%enditem
\par
%beginitem
2. Написати програму, що запитує у користувача значення дійсних параметрів $x$ та $y$, обчислює для них значення виразу 
$$\frac{\sqrt x + 96}{(\arctg x - 0.6) (y - 40)}$$ 
та повiдомляє введенi данi та результат обчислення.

%enditem
\par
%beginitem
3. Написати функцію, що для послідовності цілих (задається масивом) знаходить найдовший відрізок, в якому всі сусідні числа різні. Якщо таких відрізків кілька, то повернути перший (як індекс початку та довжину). Використовувати додаткові структури даних (у тому числі рядки) заборонено.

По завданням 2 та 3 оцінюється правильність та якість коду, а також економність обчислень.

%enditem
}
{\smallskip\smallskip \begin {scriptsize} Затверджено на засіданні кафедри теоретичної кібернетики, протокол \No 2  від   07.10.2024 р.\par\smallskip\smallskip\smallskip\smallskip
{\parbox {6.5 cm} {Зав. кафедри \dotfill Ю.В.~Крак}}\hspace {0.7cm} {\hbox to 6.5 cm { Екзаменатор \dotfill Т.О.~Карнаух}} \end{scriptsize}}
%end
\newpage
%begin
{{\begin {scriptsize}\par
{ \center  Київський національний університет імені Тараса Шевченка \par}
{\parbox {5 cm}{Освітня програма \hfill}{\parbox {7 cm}{Інформатика\hfill}}\parbox {6 cm}{\hfill Семестр 1}}\par
{\parbox {5 cm} {Навчальний предмет \hfill}\parbox {7 cm}{Програмування \hfill}\parbox {6cm}{\hfill}\par}\vspace{-15pt} \end{scriptsize}}{\center Екзаменаційний білет \No \textbf { 19 }\par}}
{%beginitem
1. Винятки та їх оброблення.%enditem
\par
%beginitem
2. Написати програму, що запитує у користувача значення дійсних параметрів $x$ та $y$, обчислює для них значення виразу 
$$\frac{\ x + 75}{(\arcsin x + 0.1) (y - 19)}$$ 
та повiдомляє введенi данi та результат обчислення.

%enditem
\par
%beginitem
3. Написати функцію, що для послідовності цілих (задається масивом) знаходить довжину найдовшого відрізку, в якому всі сусідні числа різні. Використовувати додаткові структури даних (у тому числі рядки) заборонено.

По завданням 2 та 3 оцінюється правильність та якість коду, а також економність обчислень.

%enditem
}
{\smallskip\smallskip \begin {scriptsize} Затверджено на засіданні кафедри теоретичної кібернетики, протокол \No 2  від   07.10.2024 р.\par\smallskip\smallskip\smallskip\smallskip
{\parbox {6.5 cm} {Зав. кафедри \dotfill Ю.В.~Крак}}\hspace {0.7cm} {\hbox to 6.5 cm { Екзаменатор \dotfill Т.О.~Карнаух}} \end{scriptsize}}
%end
\vspace{0.9cm}
%begin
{{\begin {scriptsize}\par
{ \center  Київський національний університет імені Тараса Шевченка \par}
{\parbox {5 cm}{Освітня програма \hfill}{\parbox {7 cm}{Інформатика\hfill}}\parbox {6 cm}{\hfill Семестр 1}}\par
{\parbox {5 cm} {Навчальний предмет \hfill}\parbox {7 cm}{Програмування \hfill}\parbox {6cm}{\hfill}\par}\vspace{-15pt} \end{scriptsize}}{\center Екзаменаційний білет \No \textbf { 20 }\par}}
{%beginitem
1. Рекурсивні означення, рекурсивні функції.%enditem
\par
%beginitem
2. Написати програму, що запитує у користувача значення дійсних параметрів $x$ та $y$, обчислює для них значення виразу 
$$\frac{\log_{2} x - 74}{(\cos x - 0.3) (y + 42)}$$ 
та повiдомляє введенi данi та результат обчислення.

%enditem
\par
%beginitem
3. Написати функцію, що для інтервалу $[a, b)$ знаходить знаходить найбільший степінь числа 3, на який ділиться хоча б одне ціле число з цього інтервалу.

По завданням 2 та 3 оцінюється правильність та якість коду, а також економність обчислень.

%enditem
}
{\smallskip\smallskip \begin {scriptsize} Затверджено на засіданні кафедри теоретичної кібернетики, протокол \No 2  від   07.10.2024 р.\par\smallskip\smallskip\smallskip\smallskip
{\parbox {6.5 cm} {Зав. кафедри \dotfill Ю.В.~Крак}}\hspace {0.7cm} {\hbox to 6.5 cm { Екзаменатор \dotfill Т.О.~Карнаух}} \end{scriptsize}}
%end
\vspace{0.9cm}
%begin
{{\begin {scriptsize}\par
{ \center  Київський національний університет імені Тараса Шевченка \par}
{\parbox {5 cm}{Освітня програма \hfill}{\parbox {7 cm}{Інформатика\hfill}}\parbox {6 cm}{\hfill Семестр 1}}\par
{\parbox {5 cm} {Навчальний предмет \hfill}\parbox {7 cm}{Програмування \hfill}\parbox {6cm}{\hfill}\par}\vspace{-15pt} \end{scriptsize}}{\center Екзаменаційний білет \No \textbf { 21 }\par}}
{%beginitem
1. Оператори порівняння (relational), рівності (equality) та булеві оператори.%enditem
\par
%beginitem
2. Написати програму, що запитує у користувача значення дійсних параметрів $x$ та $y$, обчислює для них значення виразу 
$$\frac{\log_{10} x - 31}{(\arccos x - 0.9) (y + 28)}$$ 
та повiдомляє введенi данi та результат обчислення.

%enditem
\par
%beginitem
3. Написати функцію, що для цілого числа знаходить суму його простих дільників, що входять у його розклад не більш ніж у 4му степеню. Наприклад, для числа $3^5{\cdot}7^9{\cdot}11{\cdot}23^{2}$ таких дільників два – $11$ та $23$, тому функція має повернути $34$.

По завданням 2 та 3 оцінюється правильність та якість коду, а також економність обчислень.

%enditem
}
{\smallskip\smallskip \begin {scriptsize} Затверджено на засіданні кафедри теоретичної кібернетики, протокол \No 2  від   07.10.2024 р.\par\smallskip\smallskip\smallskip\smallskip
{\parbox {6.5 cm} {Зав. кафедри \dotfill Ю.В.~Крак}}\hspace {0.7cm} {\hbox to 6.5 cm { Екзаменатор \dotfill Т.О.~Карнаух}} \end{scriptsize}}
%end
\newpage
%begin
{{\begin {scriptsize}\par
{ \center  Київський національний університет імені Тараса Шевченка \par}
{\parbox {5 cm}{Освітня програма \hfill}{\parbox {7 cm}{Інформатика\hfill}}\parbox {6 cm}{\hfill Семестр 1}}\par
{\parbox {5 cm} {Навчальний предмет \hfill}\parbox {7 cm}{Програмування \hfill}\parbox {6cm}{\hfill}\par}\vspace{-15pt} \end{scriptsize}}{\center Екзаменаційний білет \No \textbf { 22 }\par}}
{%beginitem
1. Поняття типу даних. Вбудовані арифметичні та цілі (intergral) типи мови С++.%enditem
\par
%beginitem
2. Написати програму, що запитує у користувача значення дійсних параметрів $x$ та $y$, обчислює для них значення виразу 
$$\frac{\pi x + 13}{(\sin x + 0.6) (y - 72)}$$ 
та повiдомляє введенi данi та результат обчислення.

%enditem
\par
%beginitem
3. Написати функцію, що для послідовності цілих (задається масивом) знаходить довжину найдовшого відрізку, в якому всі сусідні числа відрізняються якнайменш на 2. Використовувати додаткові структури даних (у тому числі рядки) заборонено.

По завданням 2 та 3 оцінюється правильність та якість коду, а також економність обчислень.

%enditem
}
{\smallskip\smallskip \begin {scriptsize} Затверджено на засіданні кафедри теоретичної кібернетики, протокол \No 2  від   07.10.2024 р.\par\smallskip\smallskip\smallskip\smallskip
{\parbox {6.5 cm} {Зав. кафедри \dotfill Ю.В.~Крак}}\hspace {0.7cm} {\hbox to 6.5 cm { Екзаменатор \dotfill Т.О.~Карнаух}} \end{scriptsize}}
%end
\vspace{0.9cm}
%begin
{{\begin {scriptsize}\par
{ \center  Київський національний університет імені Тараса Шевченка \par}
{\parbox {5 cm}{Освітня програма \hfill}{\parbox {7 cm}{Інформатика\hfill}}\parbox {6 cm}{\hfill Семестр 1}}\par
{\parbox {5 cm} {Навчальний предмет \hfill}\parbox {7 cm}{Програмування \hfill}\parbox {6cm}{\hfill}\par}\vspace{-15pt} \end{scriptsize}}{\center Екзаменаційний білет \No \textbf { 23 }\par}}
{%beginitem
1. Математичний підхід до визначення складності обчислень.%enditem
\par
%beginitem
2. Написати програму, що запитує у користувача значення дійсних параметрів $x$ та $y$, обчислює для них значення виразу 
$$\frac{\ln x - 54}{(\arcsin x - 0.5) (y - 83)}$$ 
та повiдомляє введенi данi та результат обчислення.

%enditem
\par
%beginitem
3. Написати функцію, що для інтервалу $[a, b)$ знаходить найближчу пару чисел з цього інтервалу, що мають від трьох до п'яти (включно) простих дільників.

По завданням 2 та 3 оцінюється правильність та якість коду, а також економність обчислень.

%enditem
}
{\smallskip\smallskip \begin {scriptsize} Затверджено на засіданні кафедри теоретичної кібернетики, протокол \No 2  від   07.10.2024 р.\par\smallskip\smallskip\smallskip\smallskip
{\parbox {6.5 cm} {Зав. кафедри \dotfill Ю.В.~Крак}}\hspace {0.7cm} {\hbox to 6.5 cm { Екзаменатор \dotfill Т.О.~Карнаух}} \end{scriptsize}}
%end
\vspace{0.9cm}
%begin
{{\begin {scriptsize}\par
{ \center  Київський національний університет імені Тараса Шевченка \par}
{\parbox {5 cm}{Освітня програма \hfill}{\parbox {7 cm}{Інформатика\hfill}}\parbox {6 cm}{\hfill Семестр 1}}\par
{\parbox {5 cm} {Навчальний предмет \hfill}\parbox {7 cm}{Програмування \hfill}\parbox {6cm}{\hfill}\par}\vspace{-15pt} \end{scriptsize}}{\center Екзаменаційний білет \No \textbf { 24 }\par}}
{%beginitem
1. Глибина рекурсії та загальна кількість рекурсивних викликів, їх вплив на розміри програмного стеку та на час виконання програми.%enditem
\par
%beginitem
2. Написати програму, що запитує у користувача значення дійсних параметрів $x$ та $y$, обчислює для них значення виразу 
$$\frac{\sqrt x + 48}{(\ctg x + 0.4) (y + 38)}$$ 
та повiдомляє введенi данi та результат обчислення.

%enditem
\par
%beginitem
3. Написати функцію, що для цілого числа знаходить кількість простих дільників, які входять у його розклад у найбільшому степені. Наприклад, для числа $3^5{\cdot}7^9{\cdot}11{\cdot}23^{9}$ таких дільників два – $23$ та $7$, тому функція має повернути $2$.

По завданням 2 та 3 оцінюється правильність та якість коду, а також економність обчислень.

%enditem
}
{\smallskip\smallskip \begin {scriptsize} Затверджено на засіданні кафедри теоретичної кібернетики, протокол \No 2  від   07.10.2024 р.\par\smallskip\smallskip\smallskip\smallskip
{\parbox {6.5 cm} {Зав. кафедри \dotfill Ю.В.~Крак}}\hspace {0.7cm} {\hbox to 6.5 cm { Екзаменатор \dotfill Т.О.~Карнаух}} \end{scriptsize}}
%end
\newpage
%begin
{{\begin {scriptsize}\par
{ \center  Київський національний університет імені Тараса Шевченка \par}
{\parbox {5 cm}{Освітня програма \hfill}{\parbox {7 cm}{Інформатика\hfill}}\parbox {6 cm}{\hfill Семестр 1}}\par
{\parbox {5 cm} {Навчальний предмет \hfill}\parbox {7 cm}{Програмування \hfill}\parbox {6cm}{\hfill}\par}\vspace{-15pt} \end{scriptsize}}{\center Екзаменаційний білет \No \textbf { 25 }\par}}
{%beginitem
1. Принципи зображення цілих та дійсних чисел у комп'ютері.%enditem
\par
%beginitem
2. Написати програму, що запитує у користувача значення дійсних параметрів $x$ та $y$, обчислює для них значення виразу 
$$\frac{\ln x + 37}{(\tg x + 0.8) (y + 14)}$$ 
та повiдомляє введенi данi та результат обчислення.

%enditem
\par
%beginitem
3. Написати функцію, що для інтервалу $[a, b)$ знаходить найближчу пару чисел з цього інтервалу, що мають не менше трьох простих дільників та мають цифру 3 у своєму десятковому запису.

По завданням 2 та 3 оцінюється правильність та якість коду, а також економність обчислень.

%enditem
}
{\smallskip\smallskip \begin {scriptsize} Затверджено на засіданні кафедри теоретичної кібернетики, протокол \No 2  від   07.10.2024 р.\par\smallskip\smallskip\smallskip\smallskip
{\parbox {6.5 cm} {Зав. кафедри \dotfill Ю.В.~Крак}}\hspace {0.7cm} {\hbox to 6.5 cm { Екзаменатор \dotfill Т.О.~Карнаух}} \end{scriptsize}}
%end
\vspace{0.9cm}
%begin
{{\begin {scriptsize}\par
{ \center  Київський національний університет імені Тараса Шевченка \par}
{\parbox {5 cm}{Освітня програма \hfill}{\parbox {7 cm}{Інформатика\hfill}}\parbox {6 cm}{\hfill Семестр 1}}\par
{\parbox {5 cm} {Навчальний предмет \hfill}\parbox {7 cm}{Програмування \hfill}\parbox {6cm}{\hfill}\par}\vspace{-15pt} \end{scriptsize}}{\center Екзаменаційний білет \No \textbf { 26 }\par}}
{%beginitem
1. Цикл for, вигляд і семантика інструкції циклу for.%enditem
\par
%beginitem
2. Написати програму, що запитує у користувача значення дійсних параметрів $x$ та $y$, обчислює для них значення виразу 
$$\frac{\pi x - 22}{(\arctg x - 0.7) (y - 94)}$$ 
та повiдомляє введенi данi та результат обчислення.

%enditem
\par
%beginitem
3. Написати функцію, що повертає число з послідовності цілих (задається масивом), що має парну кількість різних простих дільників. Якщо таких чисел кілька, то повернути найбільше. Використовувати додаткові структури даних (у тому числі рядки) заборонено.

По завданням 2 та 3 оцінюється правильність та якість коду, а також економність обчислень.

%enditem
}
{\smallskip\smallskip \begin {scriptsize} Затверджено на засіданні кафедри теоретичної кібернетики, протокол \No 2  від   07.10.2024 р.\par\smallskip\smallskip\smallskip\smallskip
{\parbox {6.5 cm} {Зав. кафедри \dotfill Ю.В.~Крак}}\hspace {0.7cm} {\hbox to 6.5 cm { Екзаменатор \dotfill Т.О.~Карнаух}} \end{scriptsize}}
%end
\vspace{0.9cm}
%begin
{{\begin {scriptsize}\par
{ \center  Київський національний університет імені Тараса Шевченка \par}
{\parbox {5 cm}{Освітня програма \hfill}{\parbox {7 cm}{Інформатика\hfill}}\parbox {6 cm}{\hfill Семестр 1}}\par
{\parbox {5 cm} {Навчальний предмет \hfill}\parbox {7 cm}{Програмування \hfill}\parbox {6cm}{\hfill}\par}\vspace{-15pt} \end{scriptsize}}{\center Екзаменаційний білет \No \textbf { 27 }\par}}
{%beginitem
1. Алгоритм двійкового пошуку та його складність.%enditem
\par
%beginitem
2. Написати програму, що запитує у користувача значення дійсних параметрів $x$ та $y$, обчислює для них значення виразу 
$$\frac{\ x - 15}{(\arccos x - 0.2) (y - 47)}$$ 
та повiдомляє введенi данi та результат обчислення.

%enditem
\par
%beginitem
3. Написати функцію, що для цілого числа знаходить кількість його простих дільників, що входять у його розклад рівно у другому степеню. Наприклад, для числа $3^2{\cdot}7^9{\cdot}11{\cdot}23^{2}$ таких дільників два – $3$ та $23$, тому функція має повернути $2$.

По завданням 2 та 3 оцінюється правильність та якість коду, а також економність обчислень.

%enditem
}
{\smallskip\smallskip \begin {scriptsize} Затверджено на засіданні кафедри теоретичної кібернетики, протокол \No 2  від   07.10.2024 р.\par\smallskip\smallskip\smallskip\smallskip
{\parbox {6.5 cm} {Зав. кафедри \dotfill Ю.В.~Крак}}\hspace {0.7cm} {\hbox to 6.5 cm { Екзаменатор \dotfill Т.О.~Карнаух}} \end{scriptsize}}
%end
\newpage
%begin
{{\begin {scriptsize}\par
{ \center  Київський національний університет імені Тараса Шевченка \par}
{\parbox {5 cm}{Освітня програма \hfill}{\parbox {7 cm}{Інформатика\hfill}}\parbox {6 cm}{\hfill Семестр 1}}\par
{\parbox {5 cm} {Навчальний предмет \hfill}\parbox {7 cm}{Програмування \hfill}\parbox {6cm}{\hfill}\par}\vspace{-15pt} \end{scriptsize}}{\center Екзаменаційний білет \No \textbf { 28 }\par}}
{%beginitem
1. Поняття часової та просторової складності обчислень. Складність у середньому. Асимптотичні оцінки складності.%enditem
\par
%beginitem
2. Написати програму, що запитує у користувача значення дійсних параметрів $x$ та $y$, обчислює для них значення виразу 
$$\frac{\log_{10} x + 86}{(\cos x + 0.7) (y + 68)}$$ 
та повiдомляє введенi данi та результат обчислення.

%enditem
\par
%beginitem
3. Написати функцію, що повертає найбільше число з послідовності цілих (задається масивом), що має не більше трьох різних простих дільників. За відсутності має повертати None. Використовувати додаткові структури даних (у тому числі рядки) заборонено.

По завданням 2 та 3 оцінюється правильність та якість коду, а також економність обчислень.

%enditem
}
{\smallskip\smallskip \begin {scriptsize} Затверджено на засіданні кафедри теоретичної кібернетики, протокол \No 2  від   07.10.2024 р.\par\smallskip\smallskip\smallskip\smallskip
{\parbox {6.5 cm} {Зав. кафедри \dotfill Ю.В.~Крак}}\hspace {0.7cm} {\hbox to 6.5 cm { Екзаменатор \dotfill Т.О.~Карнаух}} \end{scriptsize}}
%end
\vspace{0.9cm}
%begin
{{\begin {scriptsize}\par
{ \center  Київський національний університет імені Тараса Шевченка \par}
{\parbox {5 cm}{Освітня програма \hfill}{\parbox {7 cm}{Інформатика\hfill}}\parbox {6 cm}{\hfill Семестр 1}}\par
{\parbox {5 cm} {Навчальний предмет \hfill}\parbox {7 cm}{Програмування \hfill}\parbox {6cm}{\hfill}\par}\vspace{-15pt} \end{scriptsize}}{\center Екзаменаційний білет \No \textbf { 29 }\par}}
{%beginitem
1. Арифметичні операції. Вирази. Пріоритети операторів. Порядок обчислення виразів.%enditem
\par
%beginitem
2. Написати програму, що запитує у користувача значення дійсних параметрів $x$ та $y$, обчислює для них значення виразу 
$$\frac{\log_{2} x + 97}{(\arcsin x + 0.3) (y + 23)}$$ 
та повiдомляє введенi данi та результат обчислення.

%enditem
\par
%beginitem
3. Написати функцію, яка для цілого числа знаходить кількість його різних простих дільників, що мають у десятковому запису цифру 1.

По завданням 2 та 3 оцінюється правильність та якість коду, а також економність обчислень.

%enditem
}
{\smallskip\smallskip \begin {scriptsize} Затверджено на засіданні кафедри теоретичної кібернетики, протокол \No 2  від   07.10.2024 р.\par\smallskip\smallskip\smallskip\smallskip
{\parbox {6.5 cm} {Зав. кафедри \dotfill Ю.В.~Крак}}\hspace {0.7cm} {\hbox to 6.5 cm { Екзаменатор \dotfill Т.О.~Карнаух}} \end{scriptsize}}
%end
\vspace{0.9cm}
%begin
{{\begin {scriptsize}\par
{ \center  Київський національний університет імені Тараса Шевченка \par}
{\parbox {5 cm}{Освітня програма \hfill}{\parbox {7 cm}{Інформатика\hfill}}\parbox {6 cm}{\hfill Семестр 1}}\par
{\parbox {5 cm} {Навчальний предмет \hfill}\parbox {7 cm}{Програмування \hfill}\parbox {6cm}{\hfill}\par}\vspace{-15pt} \end{scriptsize}}{\center Екзаменаційний білет \No \textbf { 30 }\par}}
{%beginitem
1. Змінні, їх властивості (у тому числі час існування, область дії оголошення імені) та операції над ними.%enditem
\par
%beginitem
2. Написати програму, що запитує у користувача значення дійсних параметрів $x$ та $y$, обчислює для них значення виразу 
$$\frac{\sqrt x - 76}{(\arctg x - 0.6) (y - 64)}$$ 
та повiдомляє введенi данi та результат обчислення.

%enditem
\par
%beginitem
3. Написати функцію, що для цілого числа знаходить простий дільник, який входить у його розклад ц найбільшому степені. Якщо таких кілька, то повернути найбільший. Наприклад, для числа $2^{9}{\cdot}3^5{\cdot}7^9{\cdot}11$ функція має повернути $7$.

По завданням 2 та 3 оцінюється правильність та якість коду, а також економність обчислень.

%enditem
}
{\smallskip\smallskip \begin {scriptsize} Затверджено на засіданні кафедри теоретичної кібернетики, протокол \No 2  від   07.10.2024 р.\par\smallskip\smallskip\smallskip\smallskip
{\parbox {6.5 cm} {Зав. кафедри \dotfill Ю.В.~Крак}}\hspace {0.7cm} {\hbox to 6.5 cm { Екзаменатор \dotfill Т.О.~Карнаух}} \end{scriptsize}}
%end
\newpage
%begin
{{\begin {scriptsize}\par
{ \center  Київський національний університет імені Тараса Шевченка \par}
{\parbox {5 cm}{Освітня програма \hfill}{\parbox {7 cm}{Інформатика\hfill}}\parbox {6 cm}{\hfill Семестр 1}}\par
{\parbox {5 cm} {Навчальний предмет \hfill}\parbox {7 cm}{Програмування \hfill}\parbox {6cm}{\hfill}\par}\vspace{-15pt} \end{scriptsize}}{\center Екзаменаційний білет \No \textbf { 31 }\par}}
{%beginitem
1. Поняття масиву. Моделювання послідовностей масивами. Параметри, що зображують масиви у функціях.%enditem
\par
%beginitem
2. Написати програму, що запитує у користувача значення дійсних параметрів $x$ та $y$, обчислює для них значення виразу 
$$\frac{\ x + 45}{(\tg x + 0.9) (y - 91)}$$ 
та повiдомляє введенi данi та результат обчислення.

%enditem
\par
%beginitem
3. Написати функцію, що в інтервалі $[a, b)$ знаходить ціле число, що ділиться на найбільший степінь числа 3. Якщо таких чисел декілька, повернути найменше.

По завданням 2 та 3 оцінюється правильність та якість коду, а також економність обчислень.

%enditem
}
{\smallskip\smallskip \begin {scriptsize} Затверджено на засіданні кафедри теоретичної кібернетики, протокол \No 2  від   07.10.2024 р.\par\smallskip\smallskip\smallskip\smallskip
{\parbox {6.5 cm} {Зав. кафедри \dotfill Ю.В.~Крак}}\hspace {0.7cm} {\hbox to 6.5 cm { Екзаменатор \dotfill Т.О.~Карнаух}} \end{scriptsize}}
%end
\vspace{0.9cm}
%begin
{{\begin {scriptsize}\par
{ \center  Київський національний університет імені Тараса Шевченка \par}
{\parbox {5 cm}{Освітня програма \hfill}{\parbox {7 cm}{Інформатика\hfill}}\parbox {6 cm}{\hfill Семестр 1}}\par
{\parbox {5 cm} {Навчальний предмет \hfill}\parbox {7 cm}{Програмування \hfill}\parbox {6cm}{\hfill}\par}\vspace{-15pt} \end{scriptsize}}{\center Екзаменаційний білет \No \textbf { 32 }\par}}
{%beginitem
1. Засоби керування порядком обчислень. Інструкція if.%enditem
\par
%beginitem
2. Написати програму, що запитує у користувача значення дійсних параметрів $x$ та $y$, обчислює для них значення виразу 
$$\frac{\sqrt x - 81}{(\cos x - 0.8) (y + 70)}$$ 
та повiдомляє введенi данi та результат обчислення.

%enditem
\par
%beginitem
3. Написати функцію, що знаходить найбільше число з послідовності цілих (задається масивом), у десятковому запису якого нема цифри 4. Повертає індекс першого входження цього числа або $-1$ (за відсутності). Використовувати додаткові структури даних (у тому числі рядки) заборонено.

По завданням 2 та 3 оцінюється правильність та якість коду, а також економність обчислень.

%enditem
}
{\smallskip\smallskip \begin {scriptsize} Затверджено на засіданні кафедри теоретичної кібернетики, протокол \No 2  від   07.10.2024 р.\par\smallskip\smallskip\smallskip\smallskip
{\parbox {6.5 cm} {Зав. кафедри \dotfill Ю.В.~Крак}}\hspace {0.7cm} {\hbox to 6.5 cm { Екзаменатор \dotfill Т.О.~Карнаух}} \end{scriptsize}}
%end
\vspace{0.9cm}
%begin
{{\begin {scriptsize}\par
{ \center  Київський національний університет імені Тараса Шевченка \par}
{\parbox {5 cm}{Освітня програма \hfill}{\parbox {7 cm}{Інформатика\hfill}}\parbox {6 cm}{\hfill Семестр 1}}\par
{\parbox {5 cm} {Навчальний предмет \hfill}\parbox {7 cm}{Програмування \hfill}\parbox {6cm}{\hfill}\par}\vspace{-15pt} \end{scriptsize}}{\center Екзаменаційний білет \No \textbf { 33 }\par}}
{%beginitem
1. Емпіричний підхід до визначення складності обчислень.%enditem
\par
%beginitem
2. Написати програму, що запитує у користувача значення дійсних параметрів $x$ та $y$, обчислює для них значення виразу 
$$\frac{\log_{2} x + 95}{(\ctg x + 0.2) (y + 52)}$$ 
та повiдомляє введенi данi та результат обчислення.

%enditem
\par
%beginitem
3. Написати функцію, що для інтервалу $[a, b)$ знаходить найближчу пару чисел з цього інтервалу, що кожне з них має непарну кількість різних простих дільників. Якщо таких пар кілька, то цікавить пара з найменшим добутком.

По завданням 2 та 3 оцінюється правильність та якість коду, а також економність обчислень.

%enditem
}
{\smallskip\smallskip \begin {scriptsize} Затверджено на засіданні кафедри теоретичної кібернетики, протокол \No 2  від   07.10.2024 р.\par\smallskip\smallskip\smallskip\smallskip
{\parbox {6.5 cm} {Зав. кафедри \dotfill Ю.В.~Крак}}\hspace {0.7cm} {\hbox to 6.5 cm { Екзаменатор \dotfill Т.О.~Карнаух}} \end{scriptsize}}
%end
\newpage
%begin
{{\begin {scriptsize}\par
{ \center  Київський національний університет імені Тараса Шевченка \par}
{\parbox {5 cm}{Освітня програма \hfill}{\parbox {7 cm}{Інформатика\hfill}}\parbox {6 cm}{\hfill Семестр 1}}\par
{\parbox {5 cm} {Навчальний предмет \hfill}\parbox {7 cm}{Програмування \hfill}\parbox {6cm}{\hfill}\par}\vspace{-15pt} \end{scriptsize}}{\center Екзаменаційний білет \No \textbf { 34 }\par}}
{%beginitem
1. Компільовані та інтерпретовані мови. Переваги та недоліки.%enditem
\par
%beginitem
2. Написати програму, що запитує у користувача значення дійсних параметрів $x$ та $y$, обчислює для них значення виразу 
$$\frac{\ln x - 33}{(\sin x - 0.1) (y - 53)}$$ 
та повiдомляє введенi данi та результат обчислення.

%enditem
\par
%beginitem
3. Написати функцію, що для послідовності цілих (задається масивом) знаходить найдовший відрізок, в якому всі сусідні числа різні. Якщо таких відрізків кілька, то повернути останній (як індекс початку та довжину). Використовувати додаткові структури даних (у тому числі рядки) заборонено.

По завданням 2 та 3 оцінюється правильність та якість коду, а також економність обчислень.

%enditem
}
{\smallskip\smallskip \begin {scriptsize} Затверджено на засіданні кафедри теоретичної кібернетики, протокол \No 2  від   07.10.2024 р.\par\smallskip\smallskip\smallskip\smallskip
{\parbox {6.5 cm} {Зав. кафедри \dotfill Ю.В.~Крак}}\hspace {0.7cm} {\hbox to 6.5 cm { Екзаменатор \dotfill Т.О.~Карнаух}} \end{scriptsize}}
%end
\vspace{0.9cm}
%begin
{{\begin {scriptsize}\par
{ \center  Київський національний університет імені Тараса Шевченка \par}
{\parbox {5 cm}{Освітня програма \hfill}{\parbox {7 cm}{Інформатика\hfill}}\parbox {6 cm}{\hfill Семестр 1}}\par
{\parbox {5 cm} {Навчальний предмет \hfill}\parbox {7 cm}{Програмування \hfill}\parbox {6cm}{\hfill}\par}\vspace{-15pt} \end{scriptsize}}{\center Екзаменаційний білет \No \textbf { 35 }\par}}
{%beginitem
1. Цикли while та do-while, їх вигляд і семантика, умова продовження циклу.%enditem
\par
%beginitem
2. Написати програму, що запитує у користувача значення дійсних параметрів $x$ та $y$, обчислює для них значення виразу 
$$\frac{\log_{10} x + 11}{(\arccos x + 0.5) (y + 85)}$$ 
та повiдомляє введенi данi та результат обчислення.

%enditem
\par
%beginitem
3. Написати функцію, що для інтервалу $[a, b)$ знаходить кількість чисел, квадрати яких діляться на суму квадратів їх цифр. (Розглядаємо десятковий запис.)

По завданням 2 та 3 оцінюється правильність та якість коду, а також економність обчислень.

%enditem
}
{\smallskip\smallskip \begin {scriptsize} Затверджено на засіданні кафедри теоретичної кібернетики, протокол \No 2  від   07.10.2024 р.\par\smallskip\smallskip\smallskip\smallskip
{\parbox {6.5 cm} {Зав. кафедри \dotfill Ю.В.~Крак}}\hspace {0.7cm} {\hbox to 6.5 cm { Екзаменатор \dotfill Т.О.~Карнаух}} \end{scriptsize}}
%end
\vspace{0.9cm}
%begin
{{\begin {scriptsize}\par
{ \center  Київський національний університет імені Тараса Шевченка \par}
{\parbox {5 cm}{Освітня програма \hfill}{\parbox {7 cm}{Інформатика\hfill}}\parbox {6 cm}{\hfill Семестр 1}}\par
{\parbox {5 cm} {Навчальний предмет \hfill}\parbox {7 cm}{Програмування \hfill}\parbox {6cm}{\hfill}\par}\vspace{-15pt} \end{scriptsize}}{\center Екзаменаційний білет \No \textbf { 36 }\par}}
{%beginitem
1. Семантика функції та її виклику, параметри та аргументи.%enditem
\par
%beginitem
2. Написати програму, що запитує у користувача значення дійсних параметрів $x$ та $y$, обчислює для них значення виразу 
$$\frac{\pi x - 16}{(\arccos x - 0.4) (y - 51)}$$ 
та повiдомляє введенi данi та результат обчислення.

%enditem
\par
%beginitem
3. Написати функцію, що повертає найменше число з послідовності цілих (задається масивом), що має не більше трьох різних простих дільників. За відсутності має повертати None. Використовувати додаткові структури даних (у тому числі рядки) заборонено.

По завданням 2 та 3 оцінюється правильність та якість коду, а також економність обчислень.

%enditem
}
{\smallskip\smallskip \begin {scriptsize} Затверджено на засіданні кафедри теоретичної кібернетики, протокол \No 2  від   07.10.2024 р.\par\smallskip\smallskip\smallskip\smallskip
{\parbox {6.5 cm} {Зав. кафедри \dotfill Ю.В.~Крак}}\hspace {0.7cm} {\hbox to 6.5 cm { Екзаменатор \dotfill Т.О.~Карнаух}} \end{scriptsize}}
%end
\newpage
%begin
{{\begin {scriptsize}\par
{ \center  Київський національний університет імені Тараса Шевченка \par}
{\parbox {5 cm}{Освітня програма \hfill}{\parbox {7 cm}{Інформатика\hfill}}\parbox {6 cm}{\hfill Семестр 1}}\par
{\parbox {5 cm} {Навчальний предмет \hfill}\parbox {7 cm}{Програмування \hfill}\parbox {6cm}{\hfill}\par}\vspace{-15pt} \end{scriptsize}}{\center Екзаменаційний білет \No \textbf { 37 }\par}}
{%beginitem
1. Арифметичні операції. Вирази. Пріоритети операторів. Порядок обчислення виразів.%enditem
\par
%beginitem
2. Написати програму, що запитує у користувача значення дійсних параметрів $x$ та $y$, обчислює для них значення виразу 
$$\frac{\pi x - 32}{(\cos x - 0.4) (y - 27)}$$ 
та повiдомляє введенi данi та результат обчислення.

%enditem
\par
%beginitem
3. Написати функцію, що для інтервалу $[a, b)$ знаходить найближчу пару чисел з цього інтервалу, що кожне з них має парну кількість різних простих дільників. Якщо таких пар кілька, то цікавить пара з найменшим добутком.

По завданням 2 та 3 оцінюється правильність та якість коду, а також економність обчислень.

%enditem
}
{\smallskip\smallskip \begin {scriptsize} Затверджено на засіданні кафедри теоретичної кібернетики, протокол \No 2  від   07.10.2024 р.\par\smallskip\smallskip\smallskip\smallskip
{\parbox {6.5 cm} {Зав. кафедри \dotfill Ю.В.~Крак}}\hspace {0.7cm} {\hbox to 6.5 cm { Екзаменатор \dotfill Т.О.~Карнаух}} \end{scriptsize}}
%end
\vspace{0.9cm}
%begin
{{\begin {scriptsize}\par
{ \center  Київський національний університет імені Тараса Шевченка \par}
{\parbox {5 cm}{Освітня програма \hfill}{\parbox {7 cm}{Інформатика\hfill}}\parbox {6 cm}{\hfill Семестр 1}}\par
{\parbox {5 cm} {Навчальний предмет \hfill}\parbox {7 cm}{Програмування \hfill}\parbox {6cm}{\hfill}\par}\vspace{-15pt} \end{scriptsize}}{\center Екзаменаційний білет \No \textbf { 38 }\par}}
{%beginitem
1. Алгоритм двійкового пошуку та його складність.%enditem
\par
%beginitem
2. Написати програму, що запитує у користувача значення дійсних параметрів $x$ та $y$, обчислює для них значення виразу 
$$\frac{\sqrt x + 24}{(\arctg x + 0.9) (y + 29)}$$ 
та повiдомляє введенi данi та результат обчислення.

%enditem
\par
%beginitem
3. Написати функцію, що для інтервалу $[a, b)$ знаходить знаходить найбільший степінь числа 3, на який ділиться хоча б одне ціле число з цього інтервалу.

По завданням 2 та 3 оцінюється правильність та якість коду, а також економність обчислень.

%enditem
}
{\smallskip\smallskip \begin {scriptsize} Затверджено на засіданні кафедри теоретичної кібернетики, протокол \No 2  від   07.10.2024 р.\par\smallskip\smallskip\smallskip\smallskip
{\parbox {6.5 cm} {Зав. кафедри \dotfill Ю.В.~Крак}}\hspace {0.7cm} {\hbox to 6.5 cm { Екзаменатор \dotfill Т.О.~Карнаух}} \end{scriptsize}}
%end
\vspace{0.9cm}
%begin
{{\begin {scriptsize}\par
{ \center  Київський національний університет імені Тараса Шевченка \par}
{\parbox {5 cm}{Освітня програма \hfill}{\parbox {7 cm}{Інформатика\hfill}}\parbox {6 cm}{\hfill Семестр 1}}\par
{\parbox {5 cm} {Навчальний предмет \hfill}\parbox {7 cm}{Програмування \hfill}\parbox {6cm}{\hfill}\par}\vspace{-15pt} \end{scriptsize}}{\center Екзаменаційний білет \No \textbf { 39 }\par}}
{%beginitem
1. Цикл for, вигляд і семантика інструкції циклу for.%enditem
\par
%beginitem
2. Написати програму, що запитує у користувача значення дійсних параметрів $x$ та $y$, обчислює для них значення виразу 
$$\frac{\ln x - 62}{(\sin x - 0.3) (y + 36)}$$ 
та повiдомляє введенi данi та результат обчислення.

%enditem
\par
%beginitem
3. Написати функцію, що для послідовності цілих (задається масивом) знаходить найдовший відрізок, в якому всі сусідні числа різні. Якщо таких відрізків кілька, то повернути перший (як індекс початку та довжину). Використовувати додаткові структури даних (у тому числі рядки) заборонено.

По завданням 2 та 3 оцінюється правильність та якість коду, а також економність обчислень.

%enditem
}
{\smallskip\smallskip \begin {scriptsize} Затверджено на засіданні кафедри теоретичної кібернетики, протокол \No 2  від   07.10.2024 р.\par\smallskip\smallskip\smallskip\smallskip
{\parbox {6.5 cm} {Зав. кафедри \dotfill Ю.В.~Крак}}\hspace {0.7cm} {\hbox to 6.5 cm { Екзаменатор \dotfill Т.О.~Карнаух}} \end{scriptsize}}
%end
\newpage
%begin
{{\begin {scriptsize}\par
{ \center  Київський національний університет імені Тараса Шевченка \par}
{\parbox {5 cm}{Освітня програма \hfill}{\parbox {7 cm}{Інформатика\hfill}}\parbox {6 cm}{\hfill Семестр 1}}\par
{\parbox {5 cm} {Навчальний предмет \hfill}\parbox {7 cm}{Програмування \hfill}\parbox {6cm}{\hfill}\par}\vspace{-15pt} \end{scriptsize}}{\center Екзаменаційний білет \No \textbf { 40 }\par}}
{%beginitem
1. Емпіричний підхід до визначення складності обчислень.%enditem
\par
%beginitem
2. Написати програму, що запитує у користувача значення дійсних параметрів $x$ та $y$, обчислює для них значення виразу 
$$\frac{\ x + 17}{(\arcsin x + 0.6) (y - 80)}$$ 
та повiдомляє введенi данi та результат обчислення.

%enditem
\par
%beginitem
3. Написати функцію, що для інтервалу $[a, b)$ знаходить найближчу пару чисел з цього інтервалу, що мають від трьох до п'яти (включно) простих дільників.

По завданням 2 та 3 оцінюється правильність та якість коду, а також економність обчислень.

%enditem
}
{\smallskip\smallskip \begin {scriptsize} Затверджено на засіданні кафедри теоретичної кібернетики, протокол \No 2  від   07.10.2024 р.\par\smallskip\smallskip\smallskip\smallskip
{\parbox {6.5 cm} {Зав. кафедри \dotfill Ю.В.~Крак}}\hspace {0.7cm} {\hbox to 6.5 cm { Екзаменатор \dotfill Т.О.~Карнаух}} \end{scriptsize}}
%end
\vspace{0.9cm}
%begin
{{\begin {scriptsize}\par
{ \center  Київський національний університет імені Тараса Шевченка \par}
{\parbox {5 cm}{Освітня програма \hfill}{\parbox {7 cm}{Інформатика\hfill}}\parbox {6 cm}{\hfill Семестр 1}}\par
{\parbox {5 cm} {Навчальний предмет \hfill}\parbox {7 cm}{Програмування \hfill}\parbox {6cm}{\hfill}\par}\vspace{-15pt} \end{scriptsize}}{\center Екзаменаційний білет \No \textbf { 41 }\par}}
{%beginitem
1. Компільовані та інтерпретовані мови. Переваги та недоліки.%enditem
\par
%beginitem
2. Написати програму, що запитує у користувача значення дійсних параметрів $x$ та $y$, обчислює для них значення виразу 
$$\frac{\log_{2} x - 65}{(\tg x + 0.2) (y - 66)}$$ 
та повiдомляє введенi данi та результат обчислення.

%enditem
\par
%beginitem
3. Написати функцію, що для цілого числа знаходить кількість його простих дільників, що входять у його розклад не більш ніж у 5му степеню. Наприклад, для числа $3^8{\cdot}7^9{\cdot}11{\cdot}23^{2}$ таких дільників два – $11$ та $23$, тому функція має повернути $2$.

По завданням 2 та 3 оцінюється правильність та якість коду, а також економність обчислень.

%enditem
}
{\smallskip\smallskip \begin {scriptsize} Затверджено на засіданні кафедри теоретичної кібернетики, протокол \No 2  від   07.10.2024 р.\par\smallskip\smallskip\smallskip\smallskip
{\parbox {6.5 cm} {Зав. кафедри \dotfill Ю.В.~Крак}}\hspace {0.7cm} {\hbox to 6.5 cm { Екзаменатор \dotfill Т.О.~Карнаух}} \end{scriptsize}}
%end
\vspace{0.9cm}
%begin
{{\begin {scriptsize}\par
{ \center  Київський національний університет імені Тараса Шевченка \par}
{\parbox {5 cm}{Освітня програма \hfill}{\parbox {7 cm}{Інформатика\hfill}}\parbox {6 cm}{\hfill Семестр 1}}\par
{\parbox {5 cm} {Навчальний предмет \hfill}\parbox {7 cm}{Програмування \hfill}\parbox {6cm}{\hfill}\par}\vspace{-15pt} \end{scriptsize}}{\center Екзаменаційний білет \No \textbf { 42 }\par}}
{%beginitem
1. Цикли while та do-while, їх вигляд і семантика, умова продовження циклу.%enditem
\par
%beginitem
2. Написати програму, що запитує у користувача значення дійсних параметрів $x$ та $y$, обчислює для них значення виразу 
$$\frac{\log_{10} x + 44}{(\ctg x - 0.1) (y + 88)}$$ 
та повiдомляє введенi данi та результат обчислення.

%enditem
\par
%beginitem
3. Написати функцію, яка для цілого числа знаходить кількість його різних простих дільників, що мають у десятковому запису цифру 1.

По завданням 2 та 3 оцінюється правильність та якість коду, а також економність обчислень.

%enditem
}
{\smallskip\smallskip \begin {scriptsize} Затверджено на засіданні кафедри теоретичної кібернетики, протокол \No 2  від   07.10.2024 р.\par\smallskip\smallskip\smallskip\smallskip
{\parbox {6.5 cm} {Зав. кафедри \dotfill Ю.В.~Крак}}\hspace {0.7cm} {\hbox to 6.5 cm { Екзаменатор \dotfill Т.О.~Карнаух}} \end{scriptsize}}
%end
\newpage
%begin
{{\begin {scriptsize}\par
{ \center  Київський національний університет імені Тараса Шевченка \par}
{\parbox {5 cm}{Освітня програма \hfill}{\parbox {7 cm}{Інформатика\hfill}}\parbox {6 cm}{\hfill Семестр 1}}\par
{\parbox {5 cm} {Навчальний предмет \hfill}\parbox {7 cm}{Програмування \hfill}\parbox {6cm}{\hfill}\par}\vspace{-15pt} \end{scriptsize}}{\center Екзаменаційний білет \No \textbf { 43 }\par}}
{%beginitem
1. Засоби керування порядком обчислень. Інструкція if.%enditem
\par
%beginitem
2. Написати програму, що запитує у користувача значення дійсних параметрів $x$ та $y$, обчислює для них значення виразу 
$$\frac{\sqrt x - 78}{(\cos x - 0.7) (y - 35)}$$ 
та повiдомляє введенi данi та результат обчислення.

%enditem
\par
%beginitem
3. Написати функцію, що для цілого числа знаходить кількість простих дільників, які входять у його розклад у найбільшому степені. Наприклад, для числа $3^5{\cdot}7^9{\cdot}11{\cdot}23^{9}$ таких дільників два – $23$ та $7$, тому функція має повернути $2$.

По завданням 2 та 3 оцінюється правильність та якість коду, а також економність обчислень.

%enditem
}
{\smallskip\smallskip \begin {scriptsize} Затверджено на засіданні кафедри теоретичної кібернетики, протокол \No 2  від   07.10.2024 р.\par\smallskip\smallskip\smallskip\smallskip
{\parbox {6.5 cm} {Зав. кафедри \dotfill Ю.В.~Крак}}\hspace {0.7cm} {\hbox to 6.5 cm { Екзаменатор \dotfill Т.О.~Карнаух}} \end{scriptsize}}
%end
\vspace{0.9cm}
%begin
{{\begin {scriptsize}\par
{ \center  Київський національний університет імені Тараса Шевченка \par}
{\parbox {5 cm}{Освітня програма \hfill}{\parbox {7 cm}{Інформатика\hfill}}\parbox {6 cm}{\hfill Семестр 1}}\par
{\parbox {5 cm} {Навчальний предмет \hfill}\parbox {7 cm}{Програмування \hfill}\parbox {6cm}{\hfill}\par}\vspace{-15pt} \end{scriptsize}}{\center Екзаменаційний білет \No \textbf { 44 }\par}}
{%beginitem
1. Принципи зображення цілих та дійсних чисел у комп'ютері.%enditem
\par
%beginitem
2. Написати програму, що запитує у користувача значення дійсних параметрів $x$ та $y$, обчислює для них значення виразу 
$$\frac{\ x + 89}{(\arctg x + 0.5) (y + 99)}$$ 
та повiдомляє введенi данi та результат обчислення.

%enditem
\par
%beginitem
3. Написати функцію, що для інтервалу $[a, b)$ знаходить найближчу пару чисел з цього інтервалу, що кожне з них має непарну кількість різних простих дільників. Якщо таких пар кілька, то цікавить пара з найменшим добутком.

По завданням 2 та 3 оцінюється правильність та якість коду, а також економність обчислень.

%enditem
}
{\smallskip\smallskip \begin {scriptsize} Затверджено на засіданні кафедри теоретичної кібернетики, протокол \No 2  від   07.10.2024 р.\par\smallskip\smallskip\smallskip\smallskip
{\parbox {6.5 cm} {Зав. кафедри \dotfill Ю.В.~Крак}}\hspace {0.7cm} {\hbox to 6.5 cm { Екзаменатор \dotfill Т.О.~Карнаух}} \end{scriptsize}}
%end
\vspace{0.9cm}
%begin
{{\begin {scriptsize}\par
{ \center  Київський національний університет імені Тараса Шевченка \par}
{\parbox {5 cm}{Освітня програма \hfill}{\parbox {7 cm}{Інформатика\hfill}}\parbox {6 cm}{\hfill Семестр 1}}\par
{\parbox {5 cm} {Навчальний предмет \hfill}\parbox {7 cm}{Програмування \hfill}\parbox {6cm}{\hfill}\par}\vspace{-15pt} \end{scriptsize}}{\center Екзаменаційний білет \No \textbf { 45 }\par}}
{%beginitem
1. Рекурсивні означення, рекурсивні функції.%enditem
\par
%beginitem
2. Написати програму, що запитує у користувача значення дійсних параметрів $x$ та $y$, обчислює для них значення виразу 
$$\frac{\pi x + 29}{(\tg x + 0.8) (y - 34)}$$ 
та повiдомляє введенi данi та результат обчислення.

%enditem
\par
%beginitem
3. Написати функцію, що повертає найбільше число з послідовності цілих (задається масивом), що має рівно три різних простих дільники. За відсутності має повертати None. Використовувати додаткові структури даних (у тому числі рядки) заборонено.

По завданням 2 та 3 оцінюється правильність та якість коду, а також економність обчислень.

%enditem
}
{\smallskip\smallskip \begin {scriptsize} Затверджено на засіданні кафедри теоретичної кібернетики, протокол \No 2  від   07.10.2024 р.\par\smallskip\smallskip\smallskip\smallskip
{\parbox {6.5 cm} {Зав. кафедри \dotfill Ю.В.~Крак}}\hspace {0.7cm} {\hbox to 6.5 cm { Екзаменатор \dotfill Т.О.~Карнаух}} \end{scriptsize}}
%end
\newpage
%begin
{{\begin {scriptsize}\par
{ \center  Київський національний університет імені Тараса Шевченка \par}
{\parbox {5 cm}{Освітня програма \hfill}{\parbox {7 cm}{Інформатика\hfill}}\parbox {6 cm}{\hfill Семестр 1}}\par
{\parbox {5 cm} {Навчальний предмет \hfill}\parbox {7 cm}{Програмування \hfill}\parbox {6cm}{\hfill}\par}\vspace{-15pt} \end{scriptsize}}{\center Екзаменаційний білет \No \textbf { 46 }\par}}
{%beginitem
1. Математичний підхід до визначення складності обчислень.%enditem
\par
%beginitem
2. Написати програму, що запитує у користувача значення дійсних параметрів $x$ та $y$, обчислює для них значення виразу 
$$\frac{\log_{2} x - 65}{(\sin x - 0.3) (y + 77)}$$ 
та повiдомляє введенi данi та результат обчислення.

%enditem
\par
%beginitem
3. Написати функцію, що для цілого числа знаходить простий дільник, який входить у його розклад ц найбільшому степені. Якщо таких кілька, то повернути найбільший. Наприклад, для числа $2^{9}{\cdot}3^5{\cdot}7^9{\cdot}11$ функція має повернути $7$.

По завданням 2 та 3 оцінюється правильність та якість коду, а також економність обчислень.

%enditem
}
{\smallskip\smallskip \begin {scriptsize} Затверджено на засіданні кафедри теоретичної кібернетики, протокол \No 2  від   07.10.2024 р.\par\smallskip\smallskip\smallskip\smallskip
{\parbox {6.5 cm} {Зав. кафедри \dotfill Ю.В.~Крак}}\hspace {0.7cm} {\hbox to 6.5 cm { Екзаменатор \dotfill Т.О.~Карнаух}} \end{scriptsize}}
%end
\vspace{0.9cm}
%begin
{{\begin {scriptsize}\par
{ \center  Київський національний університет імені Тараса Шевченка \par}
{\parbox {5 cm}{Освітня програма \hfill}{\parbox {7 cm}{Інформатика\hfill}}\parbox {6 cm}{\hfill Семестр 1}}\par
{\parbox {5 cm} {Навчальний предмет \hfill}\parbox {7 cm}{Програмування \hfill}\parbox {6cm}{\hfill}\par}\vspace{-15pt} \end{scriptsize}}{\center Екзаменаційний білет \No \textbf { 47 }\par}}
{%beginitem
1. Поняття типу даних. Вбудовані арифметичні та цілі (intergral) типи мови С++.%enditem
\par
%beginitem
2. Написати програму, що запитує у користувача значення дійсних параметрів $x$ та $y$, обчислює для них значення виразу 
$$\frac{\log_{10} x + 85}{(\arcsin x + 0.8) (y - 72)}$$ 
та повiдомляє введенi данi та результат обчислення.

%enditem
\par
%beginitem
3. Написати функцію, що для послідовності цілих (задається масивом) знаходить найдовший відрізок, в якому всі сусідні числа різні. Якщо таких відрізків кілька, то повернути останній (як індекс початку та довжину). Використовувати додаткові структури даних (у тому числі рядки) заборонено.

По завданням 2 та 3 оцінюється правильність та якість коду, а також економність обчислень.

%enditem
}
{\smallskip\smallskip \begin {scriptsize} Затверджено на засіданні кафедри теоретичної кібернетики, протокол \No 2  від   07.10.2024 р.\par\smallskip\smallskip\smallskip\smallskip
{\parbox {6.5 cm} {Зав. кафедри \dotfill Ю.В.~Крак}}\hspace {0.7cm} {\hbox to 6.5 cm { Екзаменатор \dotfill Т.О.~Карнаух}} \end{scriptsize}}
%end
\vspace{0.9cm}
%begin
{{\begin {scriptsize}\par
{ \center  Київський національний університет імені Тараса Шевченка \par}
{\parbox {5 cm}{Освітня програма \hfill}{\parbox {7 cm}{Інформатика\hfill}}\parbox {6 cm}{\hfill Семестр 1}}\par
{\parbox {5 cm} {Навчальний предмет \hfill}\parbox {7 cm}{Програмування \hfill}\parbox {6cm}{\hfill}\par}\vspace{-15pt} \end{scriptsize}}{\center Екзаменаційний білет \No \textbf { 48 }\par}}
{%beginitem
1. Оператори порівняння (relational), рівності (equality) та булеві оператори.%enditem
\par
%beginitem
2. Написати програму, що запитує у користувача значення дійсних параметрів $x$ та $y$, обчислює для них значення виразу 
$$\frac{\ln x - 53}{(\arccos x - 0.6) (y + 16)}$$ 
та повiдомляє введенi данi та результат обчислення.

%enditem
\par
%beginitem
3. Написати функцію, що в інтервалі $[a, b)$ знаходить ціле число, що ділиться на найбільший степінь числа 3. Якщо таких чисел декілька, повернути найменше.

По завданням 2 та 3 оцінюється правильність та якість коду, а також економність обчислень.

%enditem
}
{\smallskip\smallskip \begin {scriptsize} Затверджено на засіданні кафедри теоретичної кібернетики, протокол \No 2  від   07.10.2024 р.\par\smallskip\smallskip\smallskip\smallskip
{\parbox {6.5 cm} {Зав. кафедри \dotfill Ю.В.~Крак}}\hspace {0.7cm} {\hbox to 6.5 cm { Екзаменатор \dotfill Т.О.~Карнаух}} \end{scriptsize}}
%end
\newpage
%begin
{{\begin {scriptsize}\par
{ \center  Київський національний університет імені Тараса Шевченка \par}
{\parbox {5 cm}{Освітня програма \hfill}{\parbox {7 cm}{Інформатика\hfill}}\parbox {6 cm}{\hfill Семестр 1}}\par
{\parbox {5 cm} {Навчальний предмет \hfill}\parbox {7 cm}{Програмування \hfill}\parbox {6cm}{\hfill}\par}\vspace{-15pt} \end{scriptsize}}{\center Екзаменаційний білет \No \textbf { 49 }\par}}
{%beginitem
1. Глибина рекурсії та загальна кількість рекурсивних викликів, їх вплив на розміри програмного стеку та на час виконання програми.%enditem
\par
%beginitem
2. Написати програму, що запитує у користувача значення дійсних параметрів $x$ та $y$, обчислює для них значення виразу 
$$\frac{\pi x - 69}{(\ctg x - 0.1) (y + 32)}$$ 
та повiдомляє введенi данi та результат обчислення.

%enditem
\par
%beginitem
3. Написати функцію, що знаходить найбільше число з послідовності цілих (задається масивом), у десятковому запису якого нема цифри 3. Повертає індекс останнього входження цього числа або $-1$ (за відсутності). Використовувати додаткові структури даних (у тому числі рядки) заборонено.

По завданням 2 та 3 оцінюється правильність та якість коду, а також економність обчислень.

%enditem
}
{\smallskip\smallskip \begin {scriptsize} Затверджено на засіданні кафедри теоретичної кібернетики, протокол \No 2  від   07.10.2024 р.\par\smallskip\smallskip\smallskip\smallskip
{\parbox {6.5 cm} {Зав. кафедри \dotfill Ю.В.~Крак}}\hspace {0.7cm} {\hbox to 6.5 cm { Екзаменатор \dotfill Т.О.~Карнаух}} \end{scriptsize}}
%end
\vspace{0.9cm}
%begin
{{\begin {scriptsize}\par
{ \center  Київський національний університет імені Тараса Шевченка \par}
{\parbox {5 cm}{Освітня програма \hfill}{\parbox {7 cm}{Інформатика\hfill}}\parbox {6 cm}{\hfill Семестр 1}}\par
{\parbox {5 cm} {Навчальний предмет \hfill}\parbox {7 cm}{Програмування \hfill}\parbox {6cm}{\hfill}\par}\vspace{-15pt} \end{scriptsize}}{\center Екзаменаційний білет \No \textbf { 50 }\par}}
{%beginitem
1. Змінні, їх властивості (у тому числі час існування, область дії оголошення імені) та операції над ними.%enditem
\par
%beginitem
2. Написати програму, що запитує у користувача значення дійсних параметрів $x$ та $y$, обчислює для них значення виразу 
$$\frac{\sqrt x + 31}{(\cos x + 0.5) (y - 20)}$$ 
та повiдомляє введенi данi та результат обчислення.

%enditem
\par
%beginitem
3. Написати функцію, що повертає число з послідовності цілих (задається масивом), що має найбільшу кількість різних простих дільників. Якщо таких чисел кілька, то повернути найменше. Використовувати додаткові структури даних (у тому числі рядки) заборонено.

По завданням 2 та 3 оцінюється правильність та якість коду, а також економність обчислень.

%enditem
}
{\smallskip\smallskip \begin {scriptsize} Затверджено на засіданні кафедри теоретичної кібернетики, протокол \No 2  від   07.10.2024 р.\par\smallskip\smallskip\smallskip\smallskip
{\parbox {6.5 cm} {Зав. кафедри \dotfill Ю.В.~Крак}}\hspace {0.7cm} {\hbox to 6.5 cm { Екзаменатор \dotfill Т.О.~Карнаух}} \end{scriptsize}}
%end
\vspace{0.9cm}
%begin
{{\begin {scriptsize}\par
{ \center  Київський національний університет імені Тараса Шевченка \par}
{\parbox {5 cm}{Освітня програма \hfill}{\parbox {7 cm}{Інформатика\hfill}}\parbox {6 cm}{\hfill Семестр 1}}\par
{\parbox {5 cm} {Навчальний предмет \hfill}\parbox {7 cm}{Програмування \hfill}\parbox {6cm}{\hfill}\par}\vspace{-15pt} \end{scriptsize}}{\center Екзаменаційний білет \No \textbf { 51 }\par}}
{%beginitem
1. Поняття масиву. Моделювання послідовностей масивами. Параметри, що зображують масиви у функціях.%enditem
\par
%beginitem
2. Написати програму, що запитує у користувача значення дійсних параметрів $x$ та $y$, обчислює для них значення виразу 
$$\frac{\ln x + 96}{(\arccos x + 0.7) (y + 84)}$$ 
та повiдомляє введенi данi та результат обчислення.

%enditem
\par
%beginitem
3. Написати функцію, що для послідовності цілих (задається масивом) знаходить довжину найдовшого відрізку, в якому всі сусідні числа відрізняються якнайменш на 2. Використовувати додаткові структури даних (у тому числі рядки) заборонено.

По завданням 2 та 3 оцінюється правильність та якість коду, а також економність обчислень.

%enditem
}
{\smallskip\smallskip \begin {scriptsize} Затверджено на засіданні кафедри теоретичної кібернетики, протокол \No 2  від   07.10.2024 р.\par\smallskip\smallskip\smallskip\smallskip
{\parbox {6.5 cm} {Зав. кафедри \dotfill Ю.В.~Крак}}\hspace {0.7cm} {\hbox to 6.5 cm { Екзаменатор \dotfill Т.О.~Карнаух}} \end{scriptsize}}
%end
\newpage
%begin
{{\begin {scriptsize}\par
{ \center  Київський національний університет імені Тараса Шевченка \par}
{\parbox {5 cm}{Освітня програма \hfill}{\parbox {7 cm}{Інформатика\hfill}}\parbox {6 cm}{\hfill Семестр 1}}\par
{\parbox {5 cm} {Навчальний предмет \hfill}\parbox {7 cm}{Програмування \hfill}\parbox {6cm}{\hfill}\par}\vspace{-15pt} \end{scriptsize}}{\center Екзаменаційний білет \No \textbf { 52 }\par}}
{%beginitem
1. Винятки та їх оброблення.%enditem
\par
%beginitem
2. Написати програму, що запитує у користувача значення дійсних параметрів $x$ та $y$, обчислює для них значення виразу 
$$\frac{\log_{10} x - 75}{(\tg x - 0.9) (y - 80)}$$ 
та повiдомляє введенi данi та результат обчислення.

%enditem
\par
%beginitem
3. Написати функцію, що знаходить найменше число з послідовності цілих (задається масивом), у десятковому запису якого нема цифри 7. Повертає індекс першого входження цього числа або $-1$ (за відсутності). Використовувати додаткові структури даних (у тому числі рядки) заборонено.

По завданням 2 та 3 оцінюється правильність та якість коду, а також економність обчислень.

%enditem
}
{\smallskip\smallskip \begin {scriptsize} Затверджено на засіданні кафедри теоретичної кібернетики, протокол \No 2  від   07.10.2024 р.\par\smallskip\smallskip\smallskip\smallskip
{\parbox {6.5 cm} {Зав. кафедри \dotfill Ю.В.~Крак}}\hspace {0.7cm} {\hbox to 6.5 cm { Екзаменатор \dotfill Т.О.~Карнаух}} \end{scriptsize}}
%end
\vspace{0.9cm}
%begin
{{\begin {scriptsize}\par
{ \center  Київський національний університет імені Тараса Шевченка \par}
{\parbox {5 cm}{Освітня програма \hfill}{\parbox {7 cm}{Інформатика\hfill}}\parbox {6 cm}{\hfill Семестр 1}}\par
{\parbox {5 cm} {Навчальний предмет \hfill}\parbox {7 cm}{Програмування \hfill}\parbox {6cm}{\hfill}\par}\vspace{-15pt} \end{scriptsize}}{\center Екзаменаційний білет \No \textbf { 53 }\par}}
{%beginitem
1. Поняття часової та просторової складності обчислень. Складність у середньому. Асимптотичні оцінки складності.%enditem
\par
%beginitem
2. Написати програму, що запитує у користувача значення дійсних параметрів $x$ та $y$, обчислює для них значення виразу 
$$\frac{\log_{2} x - 98}{(\arctg x + 0.2) (y - 34)}$$ 
та повiдомляє введенi данi та результат обчислення.

%enditem
\par
%beginitem
3. Написати функцію, що для інтервалу $[a, b)$ знаходить найближчу пару чисел з цього інтервалу, що кожне з них має непарну кількість різних простих дільників. Якщо таких пар кілька, то цікавить пара з найбільшою сумою.

По завданням 2 та 3 оцінюється правильність та якість коду, а також економність обчислень.

%enditem
}
{\smallskip\smallskip \begin {scriptsize} Затверджено на засіданні кафедри теоретичної кібернетики, протокол \No 2  від   07.10.2024 р.\par\smallskip\smallskip\smallskip\smallskip
{\parbox {6.5 cm} {Зав. кафедри \dotfill Ю.В.~Крак}}\hspace {0.7cm} {\hbox to 6.5 cm { Екзаменатор \dotfill Т.О.~Карнаух}} \end{scriptsize}}
%end
\vspace{0.9cm}
%begin
{{\begin {scriptsize}\par
{ \center  Київський національний університет імені Тараса Шевченка \par}
{\parbox {5 cm}{Освітня програма \hfill}{\parbox {7 cm}{Інформатика\hfill}}\parbox {6 cm}{\hfill Семестр 1}}\par
{\parbox {5 cm} {Навчальний предмет \hfill}\parbox {7 cm}{Програмування \hfill}\parbox {6cm}{\hfill}\par}\vspace{-15pt} \end{scriptsize}}{\center Екзаменаційний білет \No \textbf { 54 }\par}}
{%beginitem
1. Семантика функції та її виклику, параметри та аргументи.%enditem
\par
%beginitem
2. Написати програму, що запитує у користувача значення дійсних параметрів $x$ та $y$, обчислює для них значення виразу 
$$\frac{\ x + 86}{(\arcsin x - 0.4) (y + 36)}$$ 
та повiдомляє введенi данi та результат обчислення.

%enditem
\par
%beginitem
3. Написати функцію, що для цілого числа знаходить суму його простих дільників, що входять у його розклад не більш ніж у 4му степеню. Наприклад, для числа $3^5{\cdot}7^9{\cdot}11{\cdot}23^{2}$ таких дільників два – $11$ та $23$, тому функція має повернути $34$.

По завданням 2 та 3 оцінюється правильність та якість коду, а також економність обчислень.

%enditem
}
{\smallskip\smallskip \begin {scriptsize} Затверджено на засіданні кафедри теоретичної кібернетики, протокол \No 2  від   07.10.2024 р.\par\smallskip\smallskip\smallskip\smallskip
{\parbox {6.5 cm} {Зав. кафедри \dotfill Ю.В.~Крак}}\hspace {0.7cm} {\hbox to 6.5 cm { Екзаменатор \dotfill Т.О.~Карнаух}} \end{scriptsize}}
%end
\newpage
%begin
{{\begin {scriptsize}\par
{ \center  Київський національний університет імені Тараса Шевченка \par}
{\parbox {5 cm}{Освітня програма \hfill}{\parbox {7 cm}{Інформатика\hfill}}\parbox {6 cm}{\hfill Семестр 1}}\par
{\parbox {5 cm} {Навчальний предмет \hfill}\parbox {7 cm}{Програмування \hfill}\parbox {6cm}{\hfill}\par}\vspace{-15pt} \end{scriptsize}}{\center Екзаменаційний білет \No \textbf { 55 }\par}}
{%beginitem
1. Математичний підхід до визначення складності обчислень.%enditem
\par
%beginitem
2. Написати програму, що запитує у користувача значення дійсних параметрів $x$ та $y$, обчислює для них значення виразу 
$$\frac{\log_{2} x - 51}{(\ctg x - 0.8) (y - 48)}$$ 
та повiдомляє введенi данi та результат обчислення.

%enditem
\par
%beginitem
3. Написати функцію, що повертає найбільше число з послідовності цілих (задається масивом), що має не більше трьох різних простих дільників. За відсутності має повертати None. Використовувати додаткові структури даних (у тому числі рядки) заборонено.

По завданням 2 та 3 оцінюється правильність та якість коду, а також економність обчислень.

%enditem
}
{\smallskip\smallskip \begin {scriptsize} Затверджено на засіданні кафедри теоретичної кібернетики, протокол \No 2  від   07.10.2024 р.\par\smallskip\smallskip\smallskip\smallskip
{\parbox {6.5 cm} {Зав. кафедри \dotfill Ю.В.~Крак}}\hspace {0.7cm} {\hbox to 6.5 cm { Екзаменатор \dotfill Т.О.~Карнаух}} \end{scriptsize}}
%end
\vspace{0.9cm}
%begin
{{\begin {scriptsize}\par
{ \center  Київський національний університет імені Тараса Шевченка \par}
{\parbox {5 cm}{Освітня програма \hfill}{\parbox {7 cm}{Інформатика\hfill}}\parbox {6 cm}{\hfill Семестр 1}}\par
{\parbox {5 cm} {Навчальний предмет \hfill}\parbox {7 cm}{Програмування \hfill}\parbox {6cm}{\hfill}\par}\vspace{-15pt} \end{scriptsize}}{\center Екзаменаційний білет \No \textbf { 56 }\par}}
{%beginitem
1. Рекурсивні означення, рекурсивні функції.%enditem
\par
%beginitem
2. Написати програму, що запитує у користувача значення дійсних параметрів $x$ та $y$, обчислює для них значення виразу 
$$\frac{\ln x + 14}{(\sin x + 0.4) (y + 71)}$$ 
та повiдомляє введенi данi та результат обчислення.

%enditem
\par
%beginitem
3. Написати функцію, що для послідовності цілих (задається масивом) знаходить найдовший відрізок, в якому всі сусідні числа відрізняються якнайменш на 3. Якщо таких відрізків кілька, то повернути перший (як індекс початку та довжину). Використовувати додаткові структури даних (у тому числі рядки) заборонено.

По завданням 2 та 3 оцінюється правильність та якість коду, а також економність обчислень.

%enditem
}
{\smallskip\smallskip \begin {scriptsize} Затверджено на засіданні кафедри теоретичної кібернетики, протокол \No 2  від   07.10.2024 р.\par\smallskip\smallskip\smallskip\smallskip
{\parbox {6.5 cm} {Зав. кафедри \dotfill Ю.В.~Крак}}\hspace {0.7cm} {\hbox to 6.5 cm { Екзаменатор \dotfill Т.О.~Карнаух}} \end{scriptsize}}
%end
\vspace{0.9cm}
%begin
{{\begin {scriptsize}\par
{ \center  Київський національний університет імені Тараса Шевченка \par}
{\parbox {5 cm}{Освітня програма \hfill}{\parbox {7 cm}{Інформатика\hfill}}\parbox {6 cm}{\hfill Семестр 1}}\par
{\parbox {5 cm} {Навчальний предмет \hfill}\parbox {7 cm}{Програмування \hfill}\parbox {6cm}{\hfill}\par}\vspace{-15pt} \end{scriptsize}}{\center Екзаменаційний білет \No \textbf { 57 }\par}}
{%beginitem
1. Засоби керування порядком обчислень. Інструкція if.%enditem
\par
%beginitem
2. Написати програму, що запитує у користувача значення дійсних параметрів $x$ та $y$, обчислює для них значення виразу 
$$\frac{\ x - 38}{(\arcsin x + 0.2) (y - 47)}$$ 
та повiдомляє введенi данi та результат обчислення.

%enditem
\par
%beginitem
3. Написати функцію, що для послідовності цілих (задається масивом) знаходить довжину найдовшого відрізку, в якому всі сусідні числа різні. Використовувати додаткові структури даних (у тому числі рядки) заборонено.

По завданням 2 та 3 оцінюється правильність та якість коду, а також економність обчислень.

%enditem
}
{\smallskip\smallskip \begin {scriptsize} Затверджено на засіданні кафедри теоретичної кібернетики, протокол \No 2  від   07.10.2024 р.\par\smallskip\smallskip\smallskip\smallskip
{\parbox {6.5 cm} {Зав. кафедри \dotfill Ю.В.~Крак}}\hspace {0.7cm} {\hbox to 6.5 cm { Екзаменатор \dotfill Т.О.~Карнаух}} \end{scriptsize}}
%end
\newpage
%begin
{{\begin {scriptsize}\par
{ \center  Київський національний університет імені Тараса Шевченка \par}
{\parbox {5 cm}{Освітня програма \hfill}{\parbox {7 cm}{Інформатика\hfill}}\parbox {6 cm}{\hfill Семестр 1}}\par
{\parbox {5 cm} {Навчальний предмет \hfill}\parbox {7 cm}{Програмування \hfill}\parbox {6cm}{\hfill}\par}\vspace{-15pt} \end{scriptsize}}{\center Екзаменаційний білет \No \textbf { 58 }\par}}
{%beginitem
1. Арифметичні операції. Вирази. Пріоритети операторів. Порядок обчислення виразів.%enditem
\par
%beginitem
2. Написати програму, що запитує у користувача значення дійсних параметрів $x$ та $y$, обчислює для них значення виразу 
$$\frac{\pi x + 37}{(\arccos x - 0.7) (y + 30)}$$ 
та повiдомляє введенi данi та результат обчислення.

%enditem
\par
%beginitem
3. Написати функцію, що повертає найбільше число з послідовності цілих (задається масивом), у сімковому запису якого нема цифри 2. За відсутності має повертати None. Використовувати додаткові структури даних (у тому числі рядки) заборонено.

По завданням 2 та 3 оцінюється правильність та якість коду, а також економність обчислень.

%enditem
}
{\smallskip\smallskip \begin {scriptsize} Затверджено на засіданні кафедри теоретичної кібернетики, протокол \No 2  від   07.10.2024 р.\par\smallskip\smallskip\smallskip\smallskip
{\parbox {6.5 cm} {Зав. кафедри \dotfill Ю.В.~Крак}}\hspace {0.7cm} {\hbox to 6.5 cm { Екзаменатор \dotfill Т.О.~Карнаух}} \end{scriptsize}}
%end
\vspace{0.9cm}
%begin
{{\begin {scriptsize}\par
{ \center  Київський національний університет імені Тараса Шевченка \par}
{\parbox {5 cm}{Освітня програма \hfill}{\parbox {7 cm}{Інформатика\hfill}}\parbox {6 cm}{\hfill Семестр 1}}\par
{\parbox {5 cm} {Навчальний предмет \hfill}\parbox {7 cm}{Програмування \hfill}\parbox {6cm}{\hfill}\par}\vspace{-15pt} \end{scriptsize}}{\center Екзаменаційний білет \No \textbf { 59 }\par}}
{%beginitem
1. Поняття часової та просторової складності обчислень. Складність у середньому. Асимптотичні оцінки складності.%enditem
\par
%beginitem
2. Написати програму, що запитує у користувача значення дійсних параметрів $x$ та $y$, обчислює для них значення виразу 
$$\frac{\log_{10} x - 45}{(\cos x + 0.5) (y - 35)}$$ 
та повiдомляє введенi данi та результат обчислення.

%enditem
\par
%beginitem
3. Написати функцію, що для послідовності цілих (задається масивом) знаходить довжину найдовшого відрізку, що складається з чисел, у сімковому запису якого нема цифри 2. За відсутності має повертати $0$. Використовувати додаткові структури даних (у тому числі рядки) заборонено.

По завданням 2 та 3 оцінюється правильність та якість коду, а також економність обчислень.

%enditem
}
{\smallskip\smallskip \begin {scriptsize} Затверджено на засіданні кафедри теоретичної кібернетики, протокол \No 2  від   07.10.2024 р.\par\smallskip\smallskip\smallskip\smallskip
{\parbox {6.5 cm} {Зав. кафедри \dotfill Ю.В.~Крак}}\hspace {0.7cm} {\hbox to 6.5 cm { Екзаменатор \dotfill Т.О.~Карнаух}} \end{scriptsize}}
%end
\vspace{0.9cm}
%begin
{{\begin {scriptsize}\par
{ \center  Київський національний університет імені Тараса Шевченка \par}
{\parbox {5 cm}{Освітня програма \hfill}{\parbox {7 cm}{Інформатика\hfill}}\parbox {6 cm}{\hfill Семестр 1}}\par
{\parbox {5 cm} {Навчальний предмет \hfill}\parbox {7 cm}{Програмування \hfill}\parbox {6cm}{\hfill}\par}\vspace{-15pt} \end{scriptsize}}{\center Екзаменаційний білет \No \textbf { 60 }\par}}
{%beginitem
1. Оператори порівняння (relational), рівності (equality) та булеві оператори.%enditem
\par
%beginitem
2. Написати програму, що запитує у користувача значення дійсних параметрів $x$ та $y$, обчислює для них значення виразу 
$$\frac{\sqrt x + 88}{(\sin x - 0.1) (y + 64)}$$ 
та повiдомляє введенi данi та результат обчислення.

%enditem
\par
%beginitem
3. Написати функцію, що знаходить найменше число з послідовності цілих (задається масивом), у десятковому запису якого нема цифри 5. Повертає індекс останнього входження цього числа або $-1$ (за відсутності). Використовувати додаткові структури даних (у тому числі рядки) заборонено.

По завданням 2 та 3 оцінюється правильність та якість коду, а також економність обчислень.

%enditem
}
{\smallskip\smallskip \begin {scriptsize} Затверджено на засіданні кафедри теоретичної кібернетики, протокол \No 2  від   07.10.2024 р.\par\smallskip\smallskip\smallskip\smallskip
{\parbox {6.5 cm} {Зав. кафедри \dotfill Ю.В.~Крак}}\hspace {0.7cm} {\hbox to 6.5 cm { Екзаменатор \dotfill Т.О.~Карнаух}} \end{scriptsize}}
%end
\newpage
%begin
{{\begin {scriptsize}\par
{ \center  Київський національний університет імені Тараса Шевченка \par}
{\parbox {5 cm}{Освітня програма \hfill}{\parbox {7 cm}{Інформатика\hfill}}\parbox {6 cm}{\hfill Семестр 1}}\par
{\parbox {5 cm} {Навчальний предмет \hfill}\parbox {7 cm}{Програмування \hfill}\parbox {6cm}{\hfill}\par}\vspace{-15pt} \end{scriptsize}}{\center Екзаменаційний білет \No \textbf { 61 }\par}}
{%beginitem
1. Компільовані та інтерпретовані мови. Переваги та недоліки.%enditem
\par
%beginitem
2. Написати програму, що запитує у користувача значення дійсних параметрів $x$ та $y$, обчислює для них значення виразу 
$$\frac{\sqrt x + 41}{(\ctg x - 0.9) (y + 74)}$$ 
та повiдомляє введенi данi та результат обчислення.

%enditem
\par
%beginitem
3. Написати функцію, що повертає число з послідовності цілих (задається масивом), що має найменшу кількість різних простих дільників. Якщо таких чисел кілька, то повернути найбільше. Використовувати додаткові структури даних (у тому числі рядки) заборонено.

По завданням 2 та 3 оцінюється правильність та якість коду, а також економність обчислень.

%enditem
}
{\smallskip\smallskip \begin {scriptsize} Затверджено на засіданні кафедри теоретичної кібернетики, протокол \No 2  від   07.10.2024 р.\par\smallskip\smallskip\smallskip\smallskip
{\parbox {6.5 cm} {Зав. кафедри \dotfill Ю.В.~Крак}}\hspace {0.7cm} {\hbox to 6.5 cm { Екзаменатор \dotfill Т.О.~Карнаух}} \end{scriptsize}}
%end
\vspace{0.9cm}
%begin
{{\begin {scriptsize}\par
{ \center  Київський національний університет імені Тараса Шевченка \par}
{\parbox {5 cm}{Освітня програма \hfill}{\parbox {7 cm}{Інформатика\hfill}}\parbox {6 cm}{\hfill Семестр 1}}\par
{\parbox {5 cm} {Навчальний предмет \hfill}\parbox {7 cm}{Програмування \hfill}\parbox {6cm}{\hfill}\par}\vspace{-15pt} \end{scriptsize}}{\center Екзаменаційний білет \No \textbf { 62 }\par}}
{%beginitem
1. Цикли while та do-while, їх вигляд і семантика, умова продовження циклу.%enditem
\par
%beginitem
2. Написати програму, що запитує у користувача значення дійсних параметрів $x$ та $y$, обчислює для них значення виразу 
$$\frac{\log_{10} x - 70}{(\tg x + 0.3) (y - 66)}$$ 
та повiдомляє введенi данi та результат обчислення.

%enditem
\par
%beginitem
3. Написати функцію, що для інтервалу $[a, b)$ знаходить найближчу пару чисел з цього інтервалу, що кожне з них має парну кількість різних простих дільників. Якщо таких пар кілька, то цікавить пара з найбільшою сумою.

По завданням 2 та 3 оцінюється правильність та якість коду, а також економність обчислень.

%enditem
}
{\smallskip\smallskip \begin {scriptsize} Затверджено на засіданні кафедри теоретичної кібернетики, протокол \No 2  від   07.10.2024 р.\par\smallskip\smallskip\smallskip\smallskip
{\parbox {6.5 cm} {Зав. кафедри \dotfill Ю.В.~Крак}}\hspace {0.7cm} {\hbox to 6.5 cm { Екзаменатор \dotfill Т.О.~Карнаух}} \end{scriptsize}}
%end
\vspace{0.9cm}
%begin
{{\begin {scriptsize}\par
{ \center  Київський національний університет імені Тараса Шевченка \par}
{\parbox {5 cm}{Освітня програма \hfill}{\parbox {7 cm}{Інформатика\hfill}}\parbox {6 cm}{\hfill Семестр 1}}\par
{\parbox {5 cm} {Навчальний предмет \hfill}\parbox {7 cm}{Програмування \hfill}\parbox {6cm}{\hfill}\par}\vspace{-15pt} \end{scriptsize}}{\center Екзаменаційний білет \No \textbf { 63 }\par}}
{%beginitem
1. Глибина рекурсії та загальна кількість рекурсивних викликів, їх вплив на розміри програмного стеку та на час виконання програми.%enditem
\par
%beginitem
2. Написати програму, що запитує у користувача значення дійсних параметрів $x$ та $y$, обчислює для них значення виразу 
$$\frac{\log_{2} x - 61}{(\arctg x + 0.6) (y - 76)}$$ 
та повiдомляє введенi данi та результат обчислення.

%enditem
\par
%beginitem
3. Написати функцію, що в інтервалі $[a, b)$ знаходить ціле число, що ділиться на найбільший степінь числа 3. Якщо таких чисел декілька, повернути найбільше.

По завданням 2 та 3 оцінюється правильність та якість коду, а також економність обчислень.

%enditem
}
{\smallskip\smallskip \begin {scriptsize} Затверджено на засіданні кафедри теоретичної кібернетики, протокол \No 2  від   07.10.2024 р.\par\smallskip\smallskip\smallskip\smallskip
{\parbox {6.5 cm} {Зав. кафедри \dotfill Ю.В.~Крак}}\hspace {0.7cm} {\hbox to 6.5 cm { Екзаменатор \dotfill Т.О.~Карнаух}} \end{scriptsize}}
%end
\newpage
%begin
{{\begin {scriptsize}\par
{ \center  Київський національний університет імені Тараса Шевченка \par}
{\parbox {5 cm}{Освітня програма \hfill}{\parbox {7 cm}{Інформатика\hfill}}\parbox {6 cm}{\hfill Семестр 1}}\par
{\parbox {5 cm} {Навчальний предмет \hfill}\parbox {7 cm}{Програмування \hfill}\parbox {6cm}{\hfill}\par}\vspace{-15pt} \end{scriptsize}}{\center Екзаменаційний білет \No \textbf { 64 }\par}}
{%beginitem
1. Принципи зображення цілих та дійсних чисел у комп'ютері.%enditem
\par
%beginitem
2. Написати програму, що запитує у користувача значення дійсних параметрів $x$ та $y$, обчислює для них значення виразу 
$$\frac{\pi x + 15}{(\arctg x - 0.7) (y + 73)}$$ 
та повiдомляє введенi данi та результат обчислення.

%enditem
\par
%beginitem
3. Написати функцію, що для цілого числа знаходить простий дільник, який входить у його розклад у найбільшому степені. Якщо таких кілька, то повернути найменший. Наприклад, для числа $2^{9}{\cdot}3^5{\cdot}7^9{\cdot}11$ функція має повернути $2$.

По завданням 2 та 3 оцінюється правильність та якість коду, а також економність обчислень.

%enditem
}
{\smallskip\smallskip \begin {scriptsize} Затверджено на засіданні кафедри теоретичної кібернетики, протокол \No 2  від   07.10.2024 р.\par\smallskip\smallskip\smallskip\smallskip
{\parbox {6.5 cm} {Зав. кафедри \dotfill Ю.В.~Крак}}\hspace {0.7cm} {\hbox to 6.5 cm { Екзаменатор \dotfill Т.О.~Карнаух}} \end{scriptsize}}
%end
\vspace{0.9cm}
%begin
{{\begin {scriptsize}\par
{ \center  Київський національний університет імені Тараса Шевченка \par}
{\parbox {5 cm}{Освітня програма \hfill}{\parbox {7 cm}{Інформатика\hfill}}\parbox {6 cm}{\hfill Семестр 1}}\par
{\parbox {5 cm} {Навчальний предмет \hfill}\parbox {7 cm}{Програмування \hfill}\parbox {6cm}{\hfill}\par}\vspace{-15pt} \end{scriptsize}}{\center Екзаменаційний білет \No \textbf { 65 }\par}}
{%beginitem
1. Поняття масиву. Моделювання послідовностей масивами. Параметри, що зображують масиви у функціях.%enditem
\par
%beginitem
2. Написати програму, що запитує у користувача значення дійсних параметрів $x$ та $y$, обчислює для них значення виразу 
$$\frac{\ x - 24}{(\cos x - 0.1) (y - 28)}$$ 
та повiдомляє введенi данi та результат обчислення.

%enditem
\par
%beginitem
3. Написати функцію, що для інтервалу $[a, b)$ знаходить найближчу пару чисел з цього інтервалу, що мають не менше трьох простих дільників та мають цифру 3 у своєму десятковому запису.

По завданням 2 та 3 оцінюється правильність та якість коду, а також економність обчислень.

%enditem
}
{\smallskip\smallskip \begin {scriptsize} Затверджено на засіданні кафедри теоретичної кібернетики, протокол \No 2  від   07.10.2024 р.\par\smallskip\smallskip\smallskip\smallskip
{\parbox {6.5 cm} {Зав. кафедри \dotfill Ю.В.~Крак}}\hspace {0.7cm} {\hbox to 6.5 cm { Екзаменатор \dotfill Т.О.~Карнаух}} \end{scriptsize}}
%end
\vspace{0.9cm}
%begin
{{\begin {scriptsize}\par
{ \center  Київський національний університет імені Тараса Шевченка \par}
{\parbox {5 cm}{Освітня програма \hfill}{\parbox {7 cm}{Інформатика\hfill}}\parbox {6 cm}{\hfill Семестр 1}}\par
{\parbox {5 cm} {Навчальний предмет \hfill}\parbox {7 cm}{Програмування \hfill}\parbox {6cm}{\hfill}\par}\vspace{-15pt} \end{scriptsize}}{\center Екзаменаційний білет \No \textbf { 66 }\par}}
{%beginitem
1. Поняття типу даних. Вбудовані арифметичні та цілі (intergral) типи мови С++.%enditem
\par
%beginitem
2. Написати програму, що запитує у користувача значення дійсних параметрів $x$ та $y$, обчислює для них значення виразу 
$$\frac{\ln x + 56}{(\ctg x + 0.3) (y + 60)}$$ 
та повiдомляє введенi данi та результат обчислення.

%enditem
\par
%beginitem
3. Написати функцію, що повертає число з послідовності цілих (задається масивом), що має парну кількість різних простих дільників. Якщо таких чисел кілька, то повернути найбільше. Використовувати додаткові структури даних (у тому числі рядки) заборонено.

По завданням 2 та 3 оцінюється правильність та якість коду, а також економність обчислень.

%enditem
}
{\smallskip\smallskip \begin {scriptsize} Затверджено на засіданні кафедри теоретичної кібернетики, протокол \No 2  від   07.10.2024 р.\par\smallskip\smallskip\smallskip\smallskip
{\parbox {6.5 cm} {Зав. кафедри \dotfill Ю.В.~Крак}}\hspace {0.7cm} {\hbox to 6.5 cm { Екзаменатор \dotfill Т.О.~Карнаух}} \end{scriptsize}}
%end
\newpage
%begin
{{\begin {scriptsize}\par
{ \center  Київський національний університет імені Тараса Шевченка \par}
{\parbox {5 cm}{Освітня програма \hfill}{\parbox {7 cm}{Інформатика\hfill}}\parbox {6 cm}{\hfill Семестр 1}}\par
{\parbox {5 cm} {Навчальний предмет \hfill}\parbox {7 cm}{Програмування \hfill}\parbox {6cm}{\hfill}\par}\vspace{-15pt} \end{scriptsize}}{\center Екзаменаційний білет \No \textbf { 67 }\par}}
{%beginitem
1. Винятки та їх оброблення.%enditem
\par
%beginitem
2. Написати програму, що запитує у користувача значення дійсних параметрів $x$ та $y$, обчислює для них значення виразу 
$$\frac{\log_{2} x + 33}{(\arcsin x + 0.9) (y - 83)}$$ 
та повiдомляє введенi данi та результат обчислення.

%enditem
\par
%beginitem
3. Написати функцію, що для цілого числа знаходить кількість його простих дільників, що входять у його розклад рівно у другому степеню. Наприклад, для числа $3^2{\cdot}7^9{\cdot}11{\cdot}23^{2}$ таких дільників два – $3$ та $23$, тому функція має повернути $2$.

По завданням 2 та 3 оцінюється правильність та якість коду, а також економність обчислень.

%enditem
}
{\smallskip\smallskip \begin {scriptsize} Затверджено на засіданні кафедри теоретичної кібернетики, протокол \No 2  від   07.10.2024 р.\par\smallskip\smallskip\smallskip\smallskip
{\parbox {6.5 cm} {Зав. кафедри \dotfill Ю.В.~Крак}}\hspace {0.7cm} {\hbox to 6.5 cm { Екзаменатор \dotfill Т.О.~Карнаух}} \end{scriptsize}}
%end
\vspace{0.9cm}
%begin
{{\begin {scriptsize}\par
{ \center  Київський національний університет імені Тараса Шевченка \par}
{\parbox {5 cm}{Освітня програма \hfill}{\parbox {7 cm}{Інформатика\hfill}}\parbox {6 cm}{\hfill Семестр 1}}\par
{\parbox {5 cm} {Навчальний предмет \hfill}\parbox {7 cm}{Програмування \hfill}\parbox {6cm}{\hfill}\par}\vspace{-15pt} \end{scriptsize}}{\center Екзаменаційний білет \No \textbf { 68 }\par}}
{%beginitem
1. Семантика функції та її виклику, параметри та аргументи.%enditem
\par
%beginitem
2. Написати програму, що запитує у користувача значення дійсних параметрів $x$ та $y$, обчислює для них значення виразу 
$$\frac{\ x - 81}{(\sin x - 0.8) (y + 91)}$$ 
та повiдомляє введенi данi та результат обчислення.

%enditem
\par
%beginitem
3. Написати функцію, що знаходить найбільше число з послідовності цілих (задається масивом), у десятковому запису якого нема цифри 4. Повертає індекс першого входження цього числа або $-1$ (за відсутності). Використовувати додаткові структури даних (у тому числі рядки) заборонено.

По завданням 2 та 3 оцінюється правильність та якість коду, а також економність обчислень.

%enditem
}
{\smallskip\smallskip \begin {scriptsize} Затверджено на засіданні кафедри теоретичної кібернетики, протокол \No 2  від   07.10.2024 р.\par\smallskip\smallskip\smallskip\smallskip
{\parbox {6.5 cm} {Зав. кафедри \dotfill Ю.В.~Крак}}\hspace {0.7cm} {\hbox to 6.5 cm { Екзаменатор \dotfill Т.О.~Карнаух}} \end{scriptsize}}
%end
\vspace{0.9cm}
%begin
{{\begin {scriptsize}\par
{ \center  Київський національний університет імені Тараса Шевченка \par}
{\parbox {5 cm}{Освітня програма \hfill}{\parbox {7 cm}{Інформатика\hfill}}\parbox {6 cm}{\hfill Семестр 1}}\par
{\parbox {5 cm} {Навчальний предмет \hfill}\parbox {7 cm}{Програмування \hfill}\parbox {6cm}{\hfill}\par}\vspace{-15pt} \end{scriptsize}}{\center Екзаменаційний білет \No \textbf { 69 }\par}}
{%beginitem
1. Алгоритм двійкового пошуку та його складність.%enditem
\par
%beginitem
2. Написати програму, що запитує у користувача значення дійсних параметрів $x$ та $y$, обчислює для них значення виразу 
$$\frac{\log_{10} x - 26}{(\tg x + 0.5) (y - 78)}$$ 
та повiдомляє введенi данi та результат обчислення.

%enditem
\par
%beginitem
3. Написати функцію, що повертає число з послідовності цілих (задається масивом), що має парну кількість різних простих дільників. Якщо таких чисел кілька, то повернути найбільше. Використовувати додаткові структури даних (у тому числі рядки) заборонено.

По завданням 2 та 3 оцінюється правильність та якість коду, а також економність обчислень.

%enditem
}
{\smallskip\smallskip \begin {scriptsize} Затверджено на засіданні кафедри теоретичної кібернетики, протокол \No 2  від   07.10.2024 р.\par\smallskip\smallskip\smallskip\smallskip
{\parbox {6.5 cm} {Зав. кафедри \dotfill Ю.В.~Крак}}\hspace {0.7cm} {\hbox to 6.5 cm { Екзаменатор \dotfill Т.О.~Карнаух}} \end{scriptsize}}
%end
\newpage
%begin
{{\begin {scriptsize}\par
{ \center  Київський національний університет імені Тараса Шевченка \par}
{\parbox {5 cm}{Освітня програма \hfill}{\parbox {7 cm}{Інформатика\hfill}}\parbox {6 cm}{\hfill Семестр 1}}\par
{\parbox {5 cm} {Навчальний предмет \hfill}\parbox {7 cm}{Програмування \hfill}\parbox {6cm}{\hfill}\par}\vspace{-15pt} \end{scriptsize}}{\center Екзаменаційний білет \No \textbf { 70 }\par}}
{%beginitem
1. Змінні, їх властивості (у тому числі час існування, область дії оголошення імені) та операції над ними.%enditem
\par
%beginitem
2. Написати програму, що запитує у користувача значення дійсних параметрів $x$ та $y$, обчислює для них значення виразу 
$$\frac{\sqrt x + 23}{(\arccos x - 0.6) (y + 99)}$$ 
та повiдомляє введенi данi та результат обчислення.

%enditem
\par
%beginitem
3. Написати функцію, що повертає найменше число з послідовності цілих (задається масивом), що має не більше трьох різних простих дільників. За відсутності має повертати None. Використовувати додаткові структури даних (у тому числі рядки) заборонено.

По завданням 2 та 3 оцінюється правильність та якість коду, а також економність обчислень.

%enditem
}
{\smallskip\smallskip \begin {scriptsize} Затверджено на засіданні кафедри теоретичної кібернетики, протокол \No 2  від   07.10.2024 р.\par\smallskip\smallskip\smallskip\smallskip
{\parbox {6.5 cm} {Зав. кафедри \dotfill Ю.В.~Крак}}\hspace {0.7cm} {\hbox to 6.5 cm { Екзаменатор \dotfill Т.О.~Карнаух}} \end{scriptsize}}
%end
\vspace{0.9cm}
%begin
{{\begin {scriptsize}\par
{ \center  Київський національний університет імені Тараса Шевченка \par}
{\parbox {5 cm}{Освітня програма \hfill}{\parbox {7 cm}{Інформатика\hfill}}\parbox {6 cm}{\hfill Семестр 1}}\par
{\parbox {5 cm} {Навчальний предмет \hfill}\parbox {7 cm}{Програмування \hfill}\parbox {6cm}{\hfill}\par}\vspace{-15pt} \end{scriptsize}}{\center Екзаменаційний білет \No \textbf { 71 }\par}}
{%beginitem
1. Цикл for, вигляд і семантика інструкції циклу for.%enditem
\par
%beginitem
2. Написати програму, що запитує у користувача значення дійсних параметрів $x$ та $y$, обчислює для них значення виразу 
$$\frac{\ln x + 13}{(\arctg x - 0.4) (y + 49)}$$ 
та повiдомляє введенi данi та результат обчислення.

%enditem
\par
%beginitem
3. Написати функцію, що знаходить найменше число з послідовності цілих (задається масивом), у десятковому запису якого нема цифри 5. Повертає індекс останнього входження цього числа або $-1$ (за відсутності). Використовувати додаткові структури даних (у тому числі рядки) заборонено.

По завданням 2 та 3 оцінюється правильність та якість коду, а також економність обчислень.

%enditem
}
{\smallskip\smallskip \begin {scriptsize} Затверджено на засіданні кафедри теоретичної кібернетики, протокол \No 2  від   07.10.2024 р.\par\smallskip\smallskip\smallskip\smallskip
{\parbox {6.5 cm} {Зав. кафедри \dotfill Ю.В.~Крак}}\hspace {0.7cm} {\hbox to 6.5 cm { Екзаменатор \dotfill Т.О.~Карнаух}} \end{scriptsize}}
%end
\vspace{0.9cm}
%begin
{{\begin {scriptsize}\par
{ \center  Київський національний університет імені Тараса Шевченка \par}
{\parbox {5 cm}{Освітня програма \hfill}{\parbox {7 cm}{Інформатика\hfill}}\parbox {6 cm}{\hfill Семестр 1}}\par
{\parbox {5 cm} {Навчальний предмет \hfill}\parbox {7 cm}{Програмування \hfill}\parbox {6cm}{\hfill}\par}\vspace{-15pt} \end{scriptsize}}{\center Екзаменаційний білет \No \textbf { 72 }\par}}
{%beginitem
1. Емпіричний підхід до визначення складності обчислень.%enditem
\par
%beginitem
2. Написати програму, що запитує у користувача значення дійсних параметрів $x$ та $y$, обчислює для них значення виразу 
$$\frac{\pi x - 40}{(\arccos x + 0.2) (y - 57)}$$ 
та повiдомляє введенi данi та результат обчислення.

%enditem
\par
%beginitem
3. Написати функцію, що для інтервалу $[a, b)$ знаходить найближчу пару чисел з цього інтервалу, що кожне з них має непарну кількість різних простих дільників. Якщо таких пар кілька, то цікавить пара з найменшим добутком.

По завданням 2 та 3 оцінюється правильність та якість коду, а також економність обчислень.

%enditem
}
{\smallskip\smallskip \begin {scriptsize} Затверджено на засіданні кафедри теоретичної кібернетики, протокол \No 2  від   07.10.2024 р.\par\smallskip\smallskip\smallskip\smallskip
{\parbox {6.5 cm} {Зав. кафедри \dotfill Ю.В.~Крак}}\hspace {0.7cm} {\hbox to 6.5 cm { Екзаменатор \dotfill Т.О.~Карнаух}} \end{scriptsize}}
%end
\newpage
%begin
{{\begin {scriptsize}\par
{ \center  Київський національний університет імені Тараса Шевченка \par}
{\parbox {5 cm}{Освітня програма \hfill}{\parbox {7 cm}{Інформатика\hfill}}\parbox {6 cm}{\hfill Семестр 1}}\par
{\parbox {5 cm} {Навчальний предмет \hfill}\parbox {7 cm}{Програмування \hfill}\parbox {6cm}{\hfill}\par}\vspace{-15pt} \end{scriptsize}}{\center Екзаменаційний білет \No \textbf { 73 }\par}}
{%beginitem
1. Оператори порівняння (relational), рівності (equality) та булеві оператори.%enditem
\par
%beginitem
2. Написати програму, що запитує у користувача значення дійсних параметрів $x$ та $y$, обчислює для них значення виразу 
$$\frac{\sqrt x + 19}{(\cos x + 0.9) (y - 43)}$$ 
та повiдомляє введенi данi та результат обчислення.

%enditem
\par
%beginitem
3. Написати функцію, що повертає найбільше число з послідовності цілих (задається масивом), що має не більше трьох різних простих дільників. За відсутності має повертати None. Використовувати додаткові структури даних (у тому числі рядки) заборонено.

По завданням 2 та 3 оцінюється правильність та якість коду, а також економність обчислень.

%enditem
}
{\smallskip\smallskip \begin {scriptsize} Затверджено на засіданні кафедри теоретичної кібернетики, протокол \No 2  від   07.10.2024 р.\par\smallskip\smallskip\smallskip\smallskip
{\parbox {6.5 cm} {Зав. кафедри \dotfill Ю.В.~Крак}}\hspace {0.7cm} {\hbox to 6.5 cm { Екзаменатор \dotfill Т.О.~Карнаух}} \end{scriptsize}}
%end
\vspace{0.9cm}
%begin
{{\begin {scriptsize}\par
{ \center  Київський національний університет імені Тараса Шевченка \par}
{\parbox {5 cm}{Освітня програма \hfill}{\parbox {7 cm}{Інформатика\hfill}}\parbox {6 cm}{\hfill Семестр 1}}\par
{\parbox {5 cm} {Навчальний предмет \hfill}\parbox {7 cm}{Програмування \hfill}\parbox {6cm}{\hfill}\par}\vspace{-15pt} \end{scriptsize}}{\center Екзаменаційний білет \No \textbf { 74 }\par}}
{%beginitem
1. Семантика функції та її виклику, параметри та аргументи.%enditem
\par
%beginitem
2. Написати програму, що запитує у користувача значення дійсних параметрів $x$ та $y$, обчислює для них значення виразу 
$$\frac{\log_{10} x - 68}{(\ctg x - 0.4) (y + 18)}$$ 
та повiдомляє введенi данi та результат обчислення.

%enditem
\par
%beginitem
3. Написати функцію, що в інтервалі $[a, b)$ знаходить ціле число, що ділиться на найбільший степінь числа 3. Якщо таких чисел декілька, повернути найбільше.

По завданням 2 та 3 оцінюється правильність та якість коду, а також економність обчислень.

%enditem
}
{\smallskip\smallskip \begin {scriptsize} Затверджено на засіданні кафедри теоретичної кібернетики, протокол \No 2  від   07.10.2024 р.\par\smallskip\smallskip\smallskip\smallskip
{\parbox {6.5 cm} {Зав. кафедри \dotfill Ю.В.~Крак}}\hspace {0.7cm} {\hbox to 6.5 cm { Екзаменатор \dotfill Т.О.~Карнаух}} \end{scriptsize}}
%end
\vspace{0.9cm}
%begin
{{\begin {scriptsize}\par
{ \center  Київський національний університет імені Тараса Шевченка \par}
{\parbox {5 cm}{Освітня програма \hfill}{\parbox {7 cm}{Інформатика\hfill}}\parbox {6 cm}{\hfill Семестр 1}}\par
{\parbox {5 cm} {Навчальний предмет \hfill}\parbox {7 cm}{Програмування \hfill}\parbox {6cm}{\hfill}\par}\vspace{-15pt} \end{scriptsize}}{\center Екзаменаційний білет \No \textbf { 75 }\par}}
{%beginitem
1. Поняття типу даних. Вбудовані арифметичні та цілі (intergral) типи мови С++.%enditem
\par
%beginitem
2. Написати програму, що запитує у користувача значення дійсних параметрів $x$ та $y$, обчислює для них значення виразу 
$$\frac{\pi x + 44}{(\arcsin x + 0.8) (y - 92)}$$ 
та повiдомляє введенi данi та результат обчислення.

%enditem
\par
%beginitem
3. Написати функцію, що для цілого числа знаходить кількість його простих дільників, що входять у його розклад рівно у другому степеню. Наприклад, для числа $3^2{\cdot}7^9{\cdot}11{\cdot}23^{2}$ таких дільників два – $3$ та $23$, тому функція має повернути $2$.

По завданням 2 та 3 оцінюється правильність та якість коду, а також економність обчислень.

%enditem
}
{\smallskip\smallskip \begin {scriptsize} Затверджено на засіданні кафедри теоретичної кібернетики, протокол \No 2  від   07.10.2024 р.\par\smallskip\smallskip\smallskip\smallskip
{\parbox {6.5 cm} {Зав. кафедри \dotfill Ю.В.~Крак}}\hspace {0.7cm} {\hbox to 6.5 cm { Екзаменатор \dotfill Т.О.~Карнаух}} \end{scriptsize}}
%end
\newpage
%begin
{{\begin {scriptsize}\par
{ \center  Київський національний університет імені Тараса Шевченка \par}
{\parbox {5 cm}{Освітня програма \hfill}{\parbox {7 cm}{Інформатика\hfill}}\parbox {6 cm}{\hfill Семестр 1}}\par
{\parbox {5 cm} {Навчальний предмет \hfill}\parbox {7 cm}{Програмування \hfill}\parbox {6cm}{\hfill}\par}\vspace{-15pt} \end{scriptsize}}{\center Екзаменаційний білет \No \textbf { 76 }\par}}
{%beginitem
1. Арифметичні операції. Вирази. Пріоритети операторів. Порядок обчислення виразів.%enditem
\par
%beginitem
2. Написати програму, що запитує у користувача значення дійсних параметрів $x$ та $y$, обчислює для них значення виразу 
$$\frac{\ x - 12}{(\sin x - 0.6) (y + 50)}$$ 
та повiдомляє введенi данi та результат обчислення.

%enditem
\par
%beginitem
3. Написати функцію, що для цілого числа знаходить суму його простих дільників, що входять у його розклад не більш ніж у 4му степеню. Наприклад, для числа $3^5{\cdot}7^9{\cdot}11{\cdot}23^{2}$ таких дільників два – $11$ та $23$, тому функція має повернути $34$.

По завданням 2 та 3 оцінюється правильність та якість коду, а також економність обчислень.

%enditem
}
{\smallskip\smallskip \begin {scriptsize} Затверджено на засіданні кафедри теоретичної кібернетики, протокол \No 2  від   07.10.2024 р.\par\smallskip\smallskip\smallskip\smallskip
{\parbox {6.5 cm} {Зав. кафедри \dotfill Ю.В.~Крак}}\hspace {0.7cm} {\hbox to 6.5 cm { Екзаменатор \dotfill Т.О.~Карнаух}} \end{scriptsize}}
%end
\vspace{0.9cm}
%begin
{{\begin {scriptsize}\par
{ \center  Київський національний університет імені Тараса Шевченка \par}
{\parbox {5 cm}{Освітня програма \hfill}{\parbox {7 cm}{Інформатика\hfill}}\parbox {6 cm}{\hfill Семестр 1}}\par
{\parbox {5 cm} {Навчальний предмет \hfill}\parbox {7 cm}{Програмування \hfill}\parbox {6cm}{\hfill}\par}\vspace{-15pt} \end{scriptsize}}{\center Екзаменаційний білет \No \textbf { 77 }\par}}
{%beginitem
1. Поняття часової та просторової складності обчислень. Складність у середньому. Асимптотичні оцінки складності.%enditem
\par
%beginitem
2. Написати програму, що запитує у користувача значення дійсних параметрів $x$ та $y$, обчислює для них значення виразу 
$$\frac{\log_{2} x - 46}{(\tg x + 0.2) (y - 93)}$$ 
та повiдомляє введенi данi та результат обчислення.

%enditem
\par
%beginitem
3. Написати функцію, що для інтервалу $[a, b)$ знаходить найближчу пару чисел з цього інтервалу, що мають від трьох до п'яти (включно) простих дільників.

По завданням 2 та 3 оцінюється правильність та якість коду, а також економність обчислень.

%enditem
}
{\smallskip\smallskip \begin {scriptsize} Затверджено на засіданні кафедри теоретичної кібернетики, протокол \No 2  від   07.10.2024 р.\par\smallskip\smallskip\smallskip\smallskip
{\parbox {6.5 cm} {Зав. кафедри \dotfill Ю.В.~Крак}}\hspace {0.7cm} {\hbox to 6.5 cm { Екзаменатор \dotfill Т.О.~Карнаух}} \end{scriptsize}}
%end
\vspace{0.9cm}
%begin
{{\begin {scriptsize}\par
{ \center  Київський національний університет імені Тараса Шевченка \par}
{\parbox {5 cm}{Освітня програма \hfill}{\parbox {7 cm}{Інформатика\hfill}}\parbox {6 cm}{\hfill Семестр 1}}\par
{\parbox {5 cm} {Навчальний предмет \hfill}\parbox {7 cm}{Програмування \hfill}\parbox {6cm}{\hfill}\par}\vspace{-15pt} \end{scriptsize}}{\center Екзаменаційний білет \No \textbf { 78 }\par}}
{%beginitem
1. Поняття масиву. Моделювання послідовностей масивами. Параметри, що зображують масиви у функціях.%enditem
\par
%beginitem
2. Написати програму, що запитує у користувача значення дійсних параметрів $x$ та $y$, обчислює для них значення виразу 
$$\frac{\ln x + 67}{(\arccos x - 0.3) (y + 25)}$$ 
та повiдомляє введенi данi та результат обчислення.

%enditem
\par
%beginitem
3. Написати функцію, що для послідовності цілих (задається масивом) знаходить довжину найдовшого відрізку, що складається з чисел, у сімковому запису якого нема цифри 2. За відсутності має повертати $0$. Використовувати додаткові структури даних (у тому числі рядки) заборонено.

По завданням 2 та 3 оцінюється правильність та якість коду, а також економність обчислень.

%enditem
}
{\smallskip\smallskip \begin {scriptsize} Затверджено на засіданні кафедри теоретичної кібернетики, протокол \No 2  від   07.10.2024 р.\par\smallskip\smallskip\smallskip\smallskip
{\parbox {6.5 cm} {Зав. кафедри \dotfill Ю.В.~Крак}}\hspace {0.7cm} {\hbox to 6.5 cm { Екзаменатор \dotfill Т.О.~Карнаух}} \end{scriptsize}}
%end
\newpage
%begin
{{\begin {scriptsize}\par
{ \center  Київський національний університет імені Тараса Шевченка \par}
{\parbox {5 cm}{Освітня програма \hfill}{\parbox {7 cm}{Інформатика\hfill}}\parbox {6 cm}{\hfill Семестр 1}}\par
{\parbox {5 cm} {Навчальний предмет \hfill}\parbox {7 cm}{Програмування \hfill}\parbox {6cm}{\hfill}\par}\vspace{-15pt} \end{scriptsize}}{\center Екзаменаційний білет \No \textbf { 79 }\par}}
{%beginitem
1. Емпіричний підхід до визначення складності обчислень.%enditem
\par
%beginitem
2. Написати програму, що запитує у користувача значення дійсних параметрів $x$ та $y$, обчислює для них значення виразу 
$$\frac{\ x - 62}{(\tg x - 0.1) (y - 52)}$$ 
та повiдомляє введенi данi та результат обчислення.

%enditem
\par
%beginitem
3. Написати функцію, що в інтервалі $[a, b)$ знаходить ціле число, що ділиться на найбільший степінь числа 3. Якщо таких чисел декілька, повернути найменше.

По завданням 2 та 3 оцінюється правильність та якість коду, а також економність обчислень.

%enditem
}
{\smallskip\smallskip \begin {scriptsize} Затверджено на засіданні кафедри теоретичної кібернетики, протокол \No 2  від   07.10.2024 р.\par\smallskip\smallskip\smallskip\smallskip
{\parbox {6.5 cm} {Зав. кафедри \dotfill Ю.В.~Крак}}\hspace {0.7cm} {\hbox to 6.5 cm { Екзаменатор \dotfill Т.О.~Карнаух}} \end{scriptsize}}
%end
\vspace{0.9cm}
%begin
{{\begin {scriptsize}\par
{ \center  Київський національний університет імені Тараса Шевченка \par}
{\parbox {5 cm}{Освітня програма \hfill}{\parbox {7 cm}{Інформатика\hfill}}\parbox {6 cm}{\hfill Семестр 1}}\par
{\parbox {5 cm} {Навчальний предмет \hfill}\parbox {7 cm}{Програмування \hfill}\parbox {6cm}{\hfill}\par}\vspace{-15pt} \end{scriptsize}}{\center Екзаменаційний білет \No \textbf { 80 }\par}}
{%beginitem
1. Змінні, їх властивості (у тому числі час існування, область дії оголошення імені) та операції над ними.%enditem
\par
%beginitem
2. Написати програму, що запитує у користувача значення дійсних параметрів $x$ та $y$, обчислює для них значення виразу 
$$\frac{\ln x + 63}{(\ctg x + 0.7) (y + 94)}$$ 
та повiдомляє введенi данi та результат обчислення.

%enditem
\par
%beginitem
3. Написати функцію, що для послідовності цілих (задається масивом) знаходить довжину найдовшого відрізку, в якому всі сусідні числа різні. Використовувати додаткові структури даних (у тому числі рядки) заборонено.

По завданням 2 та 3 оцінюється правильність та якість коду, а також економність обчислень.

%enditem
}
{\smallskip\smallskip \begin {scriptsize} Затверджено на засіданні кафедри теоретичної кібернетики, протокол \No 2  від   07.10.2024 р.\par\smallskip\smallskip\smallskip\smallskip
{\parbox {6.5 cm} {Зав. кафедри \dotfill Ю.В.~Крак}}\hspace {0.7cm} {\hbox to 6.5 cm { Екзаменатор \dotfill Т.О.~Карнаух}} \end{scriptsize}}
%end
\vspace{0.9cm}
%begin
{{\begin {scriptsize}\par
{ \center  Київський національний університет імені Тараса Шевченка \par}
{\parbox {5 cm}{Освітня програма \hfill}{\parbox {7 cm}{Інформатика\hfill}}\parbox {6 cm}{\hfill Семестр 1}}\par
{\parbox {5 cm} {Навчальний предмет \hfill}\parbox {7 cm}{Програмування \hfill}\parbox {6cm}{\hfill}\par}\vspace{-15pt} \end{scriptsize}}{\center Екзаменаційний білет \No \textbf { 81 }\par}}
{%beginitem
1. Рекурсивні означення, рекурсивні функції.%enditem
\par
%beginitem
2. Написати програму, що запитує у користувача значення дійсних параметрів $x$ та $y$, обчислює для них значення виразу 
$$\frac{\log_{2} x - 58}{(\arcsin x + 0.5) (y + 82)}$$ 
та повiдомляє введенi данi та результат обчислення.

%enditem
\par
%beginitem
3. Написати функцію, що для інтервалу $[a, b)$ знаходить найближчу пару чисел з цього інтервалу, що мають не менше трьох простих дільників та мають цифру 3 у своєму десятковому запису.

По завданням 2 та 3 оцінюється правильність та якість коду, а також економність обчислень.

%enditem
}
{\smallskip\smallskip \begin {scriptsize} Затверджено на засіданні кафедри теоретичної кібернетики, протокол \No 2  від   07.10.2024 р.\par\smallskip\smallskip\smallskip\smallskip
{\parbox {6.5 cm} {Зав. кафедри \dotfill Ю.В.~Крак}}\hspace {0.7cm} {\hbox to 6.5 cm { Екзаменатор \dotfill Т.О.~Карнаух}} \end{scriptsize}}
%end
\newpage
%begin
{{\begin {scriptsize}\par
{ \center  Київський національний університет імені Тараса Шевченка \par}
{\parbox {5 cm}{Освітня програма \hfill}{\parbox {7 cm}{Інформатика\hfill}}\parbox {6 cm}{\hfill Семестр 1}}\par
{\parbox {5 cm} {Навчальний предмет \hfill}\parbox {7 cm}{Програмування \hfill}\parbox {6cm}{\hfill}\par}\vspace{-15pt} \end{scriptsize}}{\center Екзаменаційний білет \No \textbf { 82 }\par}}
{%beginitem
1. Принципи зображення цілих та дійсних чисел у комп'ютері.%enditem
\par
%beginitem
2. Написати програму, що запитує у користувача значення дійсних параметрів $x$ та $y$, обчислює для них значення виразу 
$$\frac{\log_{10} x + 11}{(\sin x - 0.2) (y - 17)}$$ 
та повiдомляє введенi данi та результат обчислення.

%enditem
\par
%beginitem
3. Написати функцію, яка для цілого числа знаходить кількість його різних простих дільників, що мають у десятковому запису цифру 1.

По завданням 2 та 3 оцінюється правильність та якість коду, а також економність обчислень.

%enditem
}
{\smallskip\smallskip \begin {scriptsize} Затверджено на засіданні кафедри теоретичної кібернетики, протокол \No 2  від   07.10.2024 р.\par\smallskip\smallskip\smallskip\smallskip
{\parbox {6.5 cm} {Зав. кафедри \dotfill Ю.В.~Крак}}\hspace {0.7cm} {\hbox to 6.5 cm { Екзаменатор \dotfill Т.О.~Карнаух}} \end{scriptsize}}
%end
\vspace{0.9cm}
%begin
{{\begin {scriptsize}\par
{ \center  Київський національний університет імені Тараса Шевченка \par}
{\parbox {5 cm}{Освітня програма \hfill}{\parbox {7 cm}{Інформатика\hfill}}\parbox {6 cm}{\hfill Семестр 1}}\par
{\parbox {5 cm} {Навчальний предмет \hfill}\parbox {7 cm}{Програмування \hfill}\parbox {6cm}{\hfill}\par}\vspace{-15pt} \end{scriptsize}}{\center Екзаменаційний білет \No \textbf { 83 }\par}}
{%beginitem
1. Цикли while та do-while, їх вигляд і семантика, умова продовження циклу.%enditem
\par
%beginitem
2. Написати програму, що запитує у користувача значення дійсних параметрів $x$ та $y$, обчислює для них значення виразу 
$$\frac{\sqrt x - 42}{(\cos x - 0.7) (y + 59)}$$ 
та повiдомляє введенi данi та результат обчислення.

%enditem
\par
%beginitem
3. Написати функцію, що для інтервалу $[a, b)$ знаходить найближчу пару чисел з цього інтервалу, що кожне з них має непарну кількість різних простих дільників. Якщо таких пар кілька, то цікавить пара з найбільшою сумою.

По завданням 2 та 3 оцінюється правильність та якість коду, а також економність обчислень.

%enditem
}
{\smallskip\smallskip \begin {scriptsize} Затверджено на засіданні кафедри теоретичної кібернетики, протокол \No 2  від   07.10.2024 р.\par\smallskip\smallskip\smallskip\smallskip
{\parbox {6.5 cm} {Зав. кафедри \dotfill Ю.В.~Крак}}\hspace {0.7cm} {\hbox to 6.5 cm { Екзаменатор \dotfill Т.О.~Карнаух}} \end{scriptsize}}
%end
\vspace{0.9cm}
%begin
{{\begin {scriptsize}\par
{ \center  Київський національний університет імені Тараса Шевченка \par}
{\parbox {5 cm}{Освітня програма \hfill}{\parbox {7 cm}{Інформатика\hfill}}\parbox {6 cm}{\hfill Семестр 1}}\par
{\parbox {5 cm} {Навчальний предмет \hfill}\parbox {7 cm}{Програмування \hfill}\parbox {6cm}{\hfill}\par}\vspace{-15pt} \end{scriptsize}}{\center Екзаменаційний білет \No \textbf { 84 }\par}}
{%beginitem
1. Алгоритм двійкового пошуку та його складність.%enditem
\par
%beginitem
2. Написати програму, що запитує у користувача значення дійсних параметрів $x$ та $y$, обчислює для них значення виразу 
$$\frac{\pi x + 22}{(\arctg x + 0.4) (y - 55)}$$ 
та повiдомляє введенi данi та результат обчислення.

%enditem
\par
%beginitem
3. Написати функцію, що знаходить найбільше число з послідовності цілих (задається масивом), у десятковому запису якого нема цифри 4. Повертає індекс першого входження цього числа або $-1$ (за відсутності). Використовувати додаткові структури даних (у тому числі рядки) заборонено.

По завданням 2 та 3 оцінюється правильність та якість коду, а також економність обчислень.

%enditem
}
{\smallskip\smallskip \begin {scriptsize} Затверджено на засіданні кафедри теоретичної кібернетики, протокол \No 2  від   07.10.2024 р.\par\smallskip\smallskip\smallskip\smallskip
{\parbox {6.5 cm} {Зав. кафедри \dotfill Ю.В.~Крак}}\hspace {0.7cm} {\hbox to 6.5 cm { Екзаменатор \dotfill Т.О.~Карнаух}} \end{scriptsize}}
%end
\newpage
%begin
{{\begin {scriptsize}\par
{ \center  Київський національний університет імені Тараса Шевченка \par}
{\parbox {5 cm}{Освітня програма \hfill}{\parbox {7 cm}{Інформатика\hfill}}\parbox {6 cm}{\hfill Семестр 1}}\par
{\parbox {5 cm} {Навчальний предмет \hfill}\parbox {7 cm}{Програмування \hfill}\parbox {6cm}{\hfill}\par}\vspace{-15pt} \end{scriptsize}}{\center Екзаменаційний білет \No \textbf { 85 }\par}}
{%beginitem
1. Математичний підхід до визначення складності обчислень.%enditem
\par
%beginitem
2. Написати програму, що запитує у користувача значення дійсних параметрів $x$ та $y$, обчислює для них значення виразу 
$$\frac{\ x - 21}{(\arctg x + 0.9) (y + 89)}$$ 
та повiдомляє введенi данi та результат обчислення.

%enditem
\par
%beginitem
3. Написати функцію, що для послідовності цілих (задається масивом) знаходить найдовший відрізок, в якому всі сусідні числа відрізняються якнайменш на 3. Якщо таких відрізків кілька, то повернути перший (як індекс початку та довжину). Використовувати додаткові структури даних (у тому числі рядки) заборонено.

По завданням 2 та 3 оцінюється правильність та якість коду, а також економність обчислень.

%enditem
}
{\smallskip\smallskip \begin {scriptsize} Затверджено на засіданні кафедри теоретичної кібернетики, протокол \No 2  від   07.10.2024 р.\par\smallskip\smallskip\smallskip\smallskip
{\parbox {6.5 cm} {Зав. кафедри \dotfill Ю.В.~Крак}}\hspace {0.7cm} {\hbox to 6.5 cm { Екзаменатор \dotfill Т.О.~Карнаух}} \end{scriptsize}}
%end
\vspace{0.9cm}
%begin
{{\begin {scriptsize}\par
{ \center  Київський національний університет імені Тараса Шевченка \par}
{\parbox {5 cm}{Освітня програма \hfill}{\parbox {7 cm}{Інформатика\hfill}}\parbox {6 cm}{\hfill Семестр 1}}\par
{\parbox {5 cm} {Навчальний предмет \hfill}\parbox {7 cm}{Програмування \hfill}\parbox {6cm}{\hfill}\par}\vspace{-15pt} \end{scriptsize}}{\center Екзаменаційний білет \No \textbf { 86 }\par}}
{%beginitem
1. Глибина рекурсії та загальна кількість рекурсивних викликів, їх вплив на розміри програмного стеку та на час виконання програми.%enditem
\par
%beginitem
2. Написати програму, що запитує у користувача значення дійсних параметрів $x$ та $y$, обчислює для них значення виразу 
$$\frac{\log_{2} x + 27}{(\ctg x - 0.5) (y - 90)}$$ 
та повiдомляє введенi данi та результат обчислення.

%enditem
\par
%beginitem
3. Написати функцію, що повертає найбільше число з послідовності цілих (задається масивом), що має рівно три різних простих дільники. За відсутності має повертати None. Використовувати додаткові структури даних (у тому числі рядки) заборонено.

По завданням 2 та 3 оцінюється правильність та якість коду, а також економність обчислень.

%enditem
}
{\smallskip\smallskip \begin {scriptsize} Затверджено на засіданні кафедри теоретичної кібернетики, протокол \No 2  від   07.10.2024 р.\par\smallskip\smallskip\smallskip\smallskip
{\parbox {6.5 cm} {Зав. кафедри \dotfill Ю.В.~Крак}}\hspace {0.7cm} {\hbox to 6.5 cm { Екзаменатор \dotfill Т.О.~Карнаух}} \end{scriptsize}}
%end
\vspace{0.9cm}
%begin
{{\begin {scriptsize}\par
{ \center  Київський національний університет імені Тараса Шевченка \par}
{\parbox {5 cm}{Освітня програма \hfill}{\parbox {7 cm}{Інформатика\hfill}}\parbox {6 cm}{\hfill Семестр 1}}\par
{\parbox {5 cm} {Навчальний предмет \hfill}\parbox {7 cm}{Програмування \hfill}\parbox {6cm}{\hfill}\par}\vspace{-15pt} \end{scriptsize}}{\center Екзаменаційний білет \No \textbf { 87 }\par}}
{%beginitem
1. Винятки та їх оброблення.%enditem
\par
%beginitem
2. Написати програму, що запитує у користувача значення дійсних параметрів $x$ та $y$, обчислює для них значення виразу 
$$\frac{\ln x + 79}{(\arcsin x + 0.8) (y - 97)}$$ 
та повiдомляє введенi данi та результат обчислення.

%enditem
\par
%beginitem
3. Написати функцію, що для інтервалу $[a, b)$ знаходить найближчу пару чисел з цього інтервалу, що кожне з них має парну кількість різних простих дільників. Якщо таких пар кілька, то цікавить пара з найбільшою сумою.

По завданням 2 та 3 оцінюється правильність та якість коду, а також економність обчислень.

%enditem
}
{\smallskip\smallskip \begin {scriptsize} Затверджено на засіданні кафедри теоретичної кібернетики, протокол \No 2  від   07.10.2024 р.\par\smallskip\smallskip\smallskip\smallskip
{\parbox {6.5 cm} {Зав. кафедри \dotfill Ю.В.~Крак}}\hspace {0.7cm} {\hbox to 6.5 cm { Екзаменатор \dotfill Т.О.~Карнаух}} \end{scriptsize}}
%end
\newpage
%begin
{{\begin {scriptsize}\par
{ \center  Київський національний університет імені Тараса Шевченка \par}
{\parbox {5 cm}{Освітня програма \hfill}{\parbox {7 cm}{Інформатика\hfill}}\parbox {6 cm}{\hfill Семестр 1}}\par
{\parbox {5 cm} {Навчальний предмет \hfill}\parbox {7 cm}{Програмування \hfill}\parbox {6cm}{\hfill}\par}\vspace{-15pt} \end{scriptsize}}{\center Екзаменаційний білет \No \textbf { 88 }\par}}
{%beginitem
1. Цикл for, вигляд і семантика інструкції циклу for.%enditem
\par
%beginitem
2. Написати програму, що запитує у користувача значення дійсних параметрів $x$ та $y$, обчислює для них значення виразу 
$$\frac{\log_{10} x - 87}{(\cos x - 0.1) (y + 39)}$$ 
та повiдомляє введенi данi та результат обчислення.

%enditem
\par
%beginitem
3. Написати функцію, що для інтервалу $[a, b)$ знаходить найближчу пару чисел з цього інтервалу, що кожне з них має парну кількість різних простих дільників. Якщо таких пар кілька, то цікавить пара з найменшим добутком.

По завданням 2 та 3 оцінюється правильність та якість коду, а також економність обчислень.

%enditem
}
{\smallskip\smallskip \begin {scriptsize} Затверджено на засіданні кафедри теоретичної кібернетики, протокол \No 2  від   07.10.2024 р.\par\smallskip\smallskip\smallskip\smallskip
{\parbox {6.5 cm} {Зав. кафедри \dotfill Ю.В.~Крак}}\hspace {0.7cm} {\hbox to 6.5 cm { Екзаменатор \dotfill Т.О.~Карнаух}} \end{scriptsize}}
%end
\vspace{0.9cm}
%begin
{{\begin {scriptsize}\par
{ \center  Київський національний університет імені Тараса Шевченка \par}
{\parbox {5 cm}{Освітня програма \hfill}{\parbox {7 cm}{Інформатика\hfill}}\parbox {6 cm}{\hfill Семестр 1}}\par
{\parbox {5 cm} {Навчальний предмет \hfill}\parbox {7 cm}{Програмування \hfill}\parbox {6cm}{\hfill}\par}\vspace{-15pt} \end{scriptsize}}{\center Екзаменаційний білет \No \textbf { 89 }\par}}
{%beginitem
1. Компільовані та інтерпретовані мови. Переваги та недоліки.%enditem
\par
%beginitem
2. Написати програму, що запитує у користувача значення дійсних параметрів $x$ та $y$, обчислює для них значення виразу 
$$\frac{\pi x - 54}{(\tg x - 0.6) (y + 95)}$$ 
та повiдомляє введенi данi та результат обчислення.

%enditem
\par
%beginitem
3. Написати функцію, що для цілого числа знаходить кількість його простих дільників, що входять у його розклад не більш ніж у 5му степеню. Наприклад, для числа $3^8{\cdot}7^9{\cdot}11{\cdot}23^{2}$ таких дільників два – $11$ та $23$, тому функція має повернути $2$.

По завданням 2 та 3 оцінюється правильність та якість коду, а також економність обчислень.

%enditem
}
{\smallskip\smallskip \begin {scriptsize} Затверджено на засіданні кафедри теоретичної кібернетики, протокол \No 2  від   07.10.2024 р.\par\smallskip\smallskip\smallskip\smallskip
{\parbox {6.5 cm} {Зав. кафедри \dotfill Ю.В.~Крак}}\hspace {0.7cm} {\hbox to 6.5 cm { Екзаменатор \dotfill Т.О.~Карнаух}} \end{scriptsize}}
%end
\vspace{0.9cm}
%begin
{{\begin {scriptsize}\par
{ \center  Київський національний університет імені Тараса Шевченка \par}
{\parbox {5 cm}{Освітня програма \hfill}{\parbox {7 cm}{Інформатика\hfill}}\parbox {6 cm}{\hfill Семестр 1}}\par
{\parbox {5 cm} {Навчальний предмет \hfill}\parbox {7 cm}{Програмування \hfill}\parbox {6cm}{\hfill}\par}\vspace{-15pt} \end{scriptsize}}{\center Екзаменаційний білет \No \textbf { 90 }\par}}
{%beginitem
1. Засоби керування порядком обчислень. Інструкція if.%enditem
\par
%beginitem
2. Написати програму, що запитує у користувача значення дійсних параметрів $x$ та $y$, обчислює для них значення виразу 
$$\frac{\sqrt x + 48}{(\sin x + 0.3) (y - 26)}$$ 
та повiдомляє введенi данi та результат обчислення.

%enditem
\par
%beginitem
3. Написати функцію, що повертає число з послідовності цілих (задається масивом), що має найбільшу кількість різних простих дільників. Якщо таких чисел кілька, то повернути найменше. Використовувати додаткові структури даних (у тому числі рядки) заборонено.

По завданням 2 та 3 оцінюється правильність та якість коду, а також економність обчислень.

%enditem
}
{\smallskip\smallskip \begin {scriptsize} Затверджено на засіданні кафедри теоретичної кібернетики, протокол \No 2  від   07.10.2024 р.\par\smallskip\smallskip\smallskip\smallskip
{\parbox {6.5 cm} {Зав. кафедри \dotfill Ю.В.~Крак}}\hspace {0.7cm} {\hbox to 6.5 cm { Екзаменатор \dotfill Т.О.~Карнаух}} \end{scriptsize}}
%end
\newpage
%begin
{{\begin {scriptsize}\par
{ \center  Київський національний університет імені Тараса Шевченка \par}
{\parbox {5 cm}{Освітня програма \hfill}{\parbox {7 cm}{Інформатика\hfill}}\parbox {6 cm}{\hfill Семестр 1}}\par
{\parbox {5 cm} {Навчальний предмет \hfill}\parbox {7 cm}{Програмування \hfill}\parbox {6cm}{\hfill}\par}\vspace{-15pt} \end{scriptsize}}{\center Екзаменаційний білет \No \textbf { 91 }\par}}
{%beginitem
1. Семантика функції та її виклику, параметри та аргументи.%enditem
\par
%beginitem
2. Написати програму, що запитує у користувача значення дійсних параметрів $x$ та $y$, обчислює для них значення виразу 
$$\frac{\ln x - 92}{(\arccos x - 0.6) (y - 18)}$$ 
та повiдомляє введенi данi та результат обчислення.

%enditem
\par
%beginitem
3. Написати функцію, що для послідовності цілих (задається масивом) знаходить довжину найдовшого відрізку, в якому всі сусідні числа відрізняються якнайменш на 2. Використовувати додаткові структури даних (у тому числі рядки) заборонено.

По завданням 2 та 3 оцінюється правильність та якість коду, а також економність обчислень.

%enditem
}
{\smallskip\smallskip \begin {scriptsize} Затверджено на засіданні кафедри теоретичної кібернетики, протокол \No 2  від   07.10.2024 р.\par\smallskip\smallskip\smallskip\smallskip
{\parbox {6.5 cm} {Зав. кафедри \dotfill Ю.В.~Крак}}\hspace {0.7cm} {\hbox to 6.5 cm { Екзаменатор \dotfill Т.О.~Карнаух}} \end{scriptsize}}
%end
\vspace{0.9cm}
%begin
{{\begin {scriptsize}\par
{ \center  Київський національний університет імені Тараса Шевченка \par}
{\parbox {5 cm}{Освітня програма \hfill}{\parbox {7 cm}{Інформатика\hfill}}\parbox {6 cm}{\hfill Семестр 1}}\par
{\parbox {5 cm} {Навчальний предмет \hfill}\parbox {7 cm}{Програмування \hfill}\parbox {6cm}{\hfill}\par}\vspace{-15pt} \end{scriptsize}}{\center Екзаменаційний білет \No \textbf { 92 }\par}}
{%beginitem
1. Алгоритм двійкового пошуку та його складність.%enditem
\par
%beginitem
2. Написати програму, що запитує у користувача значення дійсних параметрів $x$ та $y$, обчислює для них значення виразу 
$$\frac{\pi x + 93}{(\arctg x + 0.1) (y + 73)}$$ 
та повiдомляє введенi данi та результат обчислення.

%enditem
\par
%beginitem
3. Написати функцію, що для цілого числа знаходить простий дільник, який входить у його розклад у найбільшому степені. Якщо таких кілька, то повернути найменший. Наприклад, для числа $2^{9}{\cdot}3^5{\cdot}7^9{\cdot}11$ функція має повернути $2$.

По завданням 2 та 3 оцінюється правильність та якість коду, а також економність обчислень.

%enditem
}
{\smallskip\smallskip \begin {scriptsize} Затверджено на засіданні кафедри теоретичної кібернетики, протокол \No 2  від   07.10.2024 р.\par\smallskip\smallskip\smallskip\smallskip
{\parbox {6.5 cm} {Зав. кафедри \dotfill Ю.В.~Крак}}\hspace {0.7cm} {\hbox to 6.5 cm { Екзаменатор \dotfill Т.О.~Карнаух}} \end{scriptsize}}
%end
\vspace{0.9cm}
%begin
{{\begin {scriptsize}\par
{ \center  Київський національний університет імені Тараса Шевченка \par}
{\parbox {5 cm}{Освітня програма \hfill}{\parbox {7 cm}{Інформатика\hfill}}\parbox {6 cm}{\hfill Семестр 1}}\par
{\parbox {5 cm} {Навчальний предмет \hfill}\parbox {7 cm}{Програмування \hfill}\parbox {6cm}{\hfill}\par}\vspace{-15pt} \end{scriptsize}}{\center Екзаменаційний білет \No \textbf { 93 }\par}}
{%beginitem
1. Математичний підхід до визначення складності обчислень.%enditem
\par
%beginitem
2. Написати програму, що запитує у користувача значення дійсних параметрів $x$ та $y$, обчислює для них значення виразу 
$$\frac{\sqrt x - 22}{(\ctg x - 0.4) (y - 82)}$$ 
та повiдомляє введенi данi та результат обчислення.

%enditem
\par
%beginitem
3. Написати функцію, що повертає найбільше число з послідовності цілих (задається масивом), у сімковому запису якого нема цифри 2. За відсутності має повертати None. Використовувати додаткові структури даних (у тому числі рядки) заборонено.

По завданням 2 та 3 оцінюється правильність та якість коду, а також економність обчислень.

%enditem
}
{\smallskip\smallskip \begin {scriptsize} Затверджено на засіданні кафедри теоретичної кібернетики, протокол \No 2  від   07.10.2024 р.\par\smallskip\smallskip\smallskip\smallskip
{\parbox {6.5 cm} {Зав. кафедри \dotfill Ю.В.~Крак}}\hspace {0.7cm} {\hbox to 6.5 cm { Екзаменатор \dotfill Т.О.~Карнаух}} \end{scriptsize}}
%end
\newpage
%begin
{{\begin {scriptsize}\par
{ \center  Київський національний університет імені Тараса Шевченка \par}
{\parbox {5 cm}{Освітня програма \hfill}{\parbox {7 cm}{Інформатика\hfill}}\parbox {6 cm}{\hfill Семестр 1}}\par
{\parbox {5 cm} {Навчальний предмет \hfill}\parbox {7 cm}{Програмування \hfill}\parbox {6cm}{\hfill}\par}\vspace{-15pt} \end{scriptsize}}{\center Екзаменаційний білет \No \textbf { 94 }\par}}
{%beginitem
1. Винятки та їх оброблення.%enditem
\par
%beginitem
2. Написати програму, що запитує у користувача значення дійсних параметрів $x$ та $y$, обчислює для них значення виразу 
$$\frac{\log_{10} x + 83}{(\sin x + 0.5) (y + 76)}$$ 
та повiдомляє введенi данi та результат обчислення.

%enditem
\par
%beginitem
3. Написати функцію, що знаходить найменше число з послідовності цілих (задається масивом), у десятковому запису якого нема цифри 7. Повертає індекс першого входження цього числа або $-1$ (за відсутності). Використовувати додаткові структури даних (у тому числі рядки) заборонено.

По завданням 2 та 3 оцінюється правильність та якість коду, а також економність обчислень.

%enditem
}
{\smallskip\smallskip \begin {scriptsize} Затверджено на засіданні кафедри теоретичної кібернетики, протокол \No 2  від   07.10.2024 р.\par\smallskip\smallskip\smallskip\smallskip
{\parbox {6.5 cm} {Зав. кафедри \dotfill Ю.В.~Крак}}\hspace {0.7cm} {\hbox to 6.5 cm { Екзаменатор \dotfill Т.О.~Карнаух}} \end{scriptsize}}
%end
\vspace{0.9cm}
%begin
{{\begin {scriptsize}\par
{ \center  Київський національний університет імені Тараса Шевченка \par}
{\parbox {5 cm}{Освітня програма \hfill}{\parbox {7 cm}{Інформатика\hfill}}\parbox {6 cm}{\hfill Семестр 1}}\par
{\parbox {5 cm} {Навчальний предмет \hfill}\parbox {7 cm}{Програмування \hfill}\parbox {6cm}{\hfill}\par}\vspace{-15pt} \end{scriptsize}}{\center Екзаменаційний білет \No \textbf { 95 }\par}}
{%beginitem
1. Емпіричний підхід до визначення складності обчислень.%enditem
\par
%beginitem
2. Написати програму, що запитує у користувача значення дійсних параметрів $x$ та $y$, обчислює для них значення виразу 
$$\frac{\log_{2} x - 58}{(\arccos x - 0.3) (y - 14)}$$ 
та повiдомляє введенi данi та результат обчислення.

%enditem
\par
%beginitem
3. Написати функцію, що повертає число з послідовності цілих (задається масивом), що має найменшу кількість різних простих дільників. Якщо таких чисел кілька, то повернути найбільше. Використовувати додаткові структури даних (у тому числі рядки) заборонено.

По завданням 2 та 3 оцінюється правильність та якість коду, а також економність обчислень.

%enditem
}
{\smallskip\smallskip \begin {scriptsize} Затверджено на засіданні кафедри теоретичної кібернетики, протокол \No 2  від   07.10.2024 р.\par\smallskip\smallskip\smallskip\smallskip
{\parbox {6.5 cm} {Зав. кафедри \dotfill Ю.В.~Крак}}\hspace {0.7cm} {\hbox to 6.5 cm { Екзаменатор \dotfill Т.О.~Карнаух}} \end{scriptsize}}
%end
\vspace{0.9cm}
%begin
{{\begin {scriptsize}\par
{ \center  Київський національний університет імені Тараса Шевченка \par}
{\parbox {5 cm}{Освітня програма \hfill}{\parbox {7 cm}{Інформатика\hfill}}\parbox {6 cm}{\hfill Семестр 1}}\par
{\parbox {5 cm} {Навчальний предмет \hfill}\parbox {7 cm}{Програмування \hfill}\parbox {6cm}{\hfill}\par}\vspace{-15pt} \end{scriptsize}}{\center Екзаменаційний білет \No \textbf { 96 }\par}}
{%beginitem
1. Компільовані та інтерпретовані мови. Переваги та недоліки.%enditem
\par
%beginitem
2. Написати програму, що запитує у користувача значення дійсних параметрів $x$ та $y$, обчислює для них значення виразу 
$$\frac{\ x + 74}{(\tg x + 0.2) (y + 95)}$$ 
та повiдомляє введенi данi та результат обчислення.

%enditem
\par
%beginitem
3. Написати функцію, що для цілого числа знаходить кількість простих дільників, які входять у його розклад у найбільшому степені. Наприклад, для числа $3^5{\cdot}7^9{\cdot}11{\cdot}23^{9}$ таких дільників два – $23$ та $7$, тому функція має повернути $2$.

По завданням 2 та 3 оцінюється правильність та якість коду, а також економність обчислень.

%enditem
}
{\smallskip\smallskip \begin {scriptsize} Затверджено на засіданні кафедри теоретичної кібернетики, протокол \No 2  від   07.10.2024 р.\par\smallskip\smallskip\smallskip\smallskip
{\parbox {6.5 cm} {Зав. кафедри \dotfill Ю.В.~Крак}}\hspace {0.7cm} {\hbox to 6.5 cm { Екзаменатор \dotfill Т.О.~Карнаух}} \end{scriptsize}}
%end
\newpage
%begin
{{\begin {scriptsize}\par
{ \center  Київський національний університет імені Тараса Шевченка \par}
{\parbox {5 cm}{Освітня програма \hfill}{\parbox {7 cm}{Інформатика\hfill}}\parbox {6 cm}{\hfill Семестр 1}}\par
{\parbox {5 cm} {Навчальний предмет \hfill}\parbox {7 cm}{Програмування \hfill}\parbox {6cm}{\hfill}\par}\vspace{-15pt} \end{scriptsize}}{\center Екзаменаційний білет \No \textbf { 97 }\par}}
{%beginitem
1. Поняття масиву. Моделювання послідовностей масивами. Параметри, що зображують масиви у функціях.%enditem
\par
%beginitem
2. Написати програму, що запитує у користувача значення дійсних параметрів $x$ та $y$, обчислює для них значення виразу 
$$\frac{\ln x - 40}{(\cos x - 0.9) (y - 25)}$$ 
та повiдомляє введенi данi та результат обчислення.

%enditem
\par
%beginitem
3. Написати функцію, що для послідовності цілих (задається масивом) знаходить найдовший відрізок, в якому всі сусідні числа різні. Якщо таких відрізків кілька, то повернути перший (як індекс початку та довжину). Використовувати додаткові структури даних (у тому числі рядки) заборонено.

По завданням 2 та 3 оцінюється правильність та якість коду, а також економність обчислень.

%enditem
}
{\smallskip\smallskip \begin {scriptsize} Затверджено на засіданні кафедри теоретичної кібернетики, протокол \No 2  від   07.10.2024 р.\par\smallskip\smallskip\smallskip\smallskip
{\parbox {6.5 cm} {Зав. кафедри \dotfill Ю.В.~Крак}}\hspace {0.7cm} {\hbox to 6.5 cm { Екзаменатор \dotfill Т.О.~Карнаух}} \end{scriptsize}}
%end
\vspace{0.9cm}
%begin
{{\begin {scriptsize}\par
{ \center  Київський національний університет імені Тараса Шевченка \par}
{\parbox {5 cm}{Освітня програма \hfill}{\parbox {7 cm}{Інформатика\hfill}}\parbox {6 cm}{\hfill Семестр 1}}\par
{\parbox {5 cm} {Навчальний предмет \hfill}\parbox {7 cm}{Програмування \hfill}\parbox {6cm}{\hfill}\par}\vspace{-15pt} \end{scriptsize}}{\center Екзаменаційний білет \No \textbf { 98 }\par}}
{%beginitem
1. Поняття типу даних. Вбудовані арифметичні та цілі (intergral) типи мови С++.%enditem
\par
%beginitem
2. Написати програму, що запитує у користувача значення дійсних параметрів $x$ та $y$, обчислює для них значення виразу 
$$\frac{\log_{10} x + 29}{(\arcsin x + 0.7) (y + 16)}$$ 
та повiдомляє введенi данi та результат обчислення.

%enditem
\par
%beginitem
3. Написати функцію, що для інтервалу $[a, b)$ знаходить знаходить найбільший степінь числа 3, на який ділиться хоча б одне ціле число з цього інтервалу.

По завданням 2 та 3 оцінюється правильність та якість коду, а також економність обчислень.

%enditem
}
{\smallskip\smallskip \begin {scriptsize} Затверджено на засіданні кафедри теоретичної кібернетики, протокол \No 2  від   07.10.2024 р.\par\smallskip\smallskip\smallskip\smallskip
{\parbox {6.5 cm} {Зав. кафедри \dotfill Ю.В.~Крак}}\hspace {0.7cm} {\hbox to 6.5 cm { Екзаменатор \dotfill Т.О.~Карнаух}} \end{scriptsize}}
%end
\vspace{0.9cm}
%begin
{{\begin {scriptsize}\par
{ \center  Київський національний університет імені Тараса Шевченка \par}
{\parbox {5 cm}{Освітня програма \hfill}{\parbox {7 cm}{Інформатика\hfill}}\parbox {6 cm}{\hfill Семестр 1}}\par
{\parbox {5 cm} {Навчальний предмет \hfill}\parbox {7 cm}{Програмування \hfill}\parbox {6cm}{\hfill}\par}\vspace{-15pt} \end{scriptsize}}{\center Екзаменаційний білет \No \textbf { 99 }\par}}
{%beginitem
1. Оператори порівняння (relational), рівності (equality) та булеві оператори.%enditem
\par
%beginitem
2. Написати програму, що запитує у користувача значення дійсних параметрів $x$ та $y$, обчислює для них значення виразу 
$$\frac{\log_{2} x - 60}{(\ctg x - 0.8) (y - 90)}$$ 
та повiдомляє введенi данi та результат обчислення.

%enditem
\par
%beginitem
3. Написати функцію, що для послідовності цілих (задається масивом) знаходить найдовший відрізок, в якому всі сусідні числа різні. Якщо таких відрізків кілька, то повернути останній (як індекс початку та довжину). Використовувати додаткові структури даних (у тому числі рядки) заборонено.

По завданням 2 та 3 оцінюється правильність та якість коду, а також економність обчислень.

%enditem
}
{\smallskip\smallskip \begin {scriptsize} Затверджено на засіданні кафедри теоретичної кібернетики, протокол \No 2  від   07.10.2024 р.\par\smallskip\smallskip\smallskip\smallskip
{\parbox {6.5 cm} {Зав. кафедри \dotfill Ю.В.~Крак}}\hspace {0.7cm} {\hbox to 6.5 cm { Екзаменатор \dotfill Т.О.~Карнаух}} \end{scriptsize}}
%end
\newpage
%begin
{{\begin {scriptsize}\par
{ \center  Київський національний університет імені Тараса Шевченка \par}
{\parbox {5 cm}{Освітня програма \hfill}{\parbox {7 cm}{Інформатика\hfill}}\parbox {6 cm}{\hfill Семестр 1}}\par
{\parbox {5 cm} {Навчальний предмет \hfill}\parbox {7 cm}{Програмування \hfill}\parbox {6cm}{\hfill}\par}\vspace{-15pt} \end{scriptsize}}{\center Екзаменаційний білет \No \textbf { 100 }\par}}
{%beginitem
1. Засоби керування порядком обчислень. Інструкція if.%enditem
\par
%beginitem
2. Написати програму, що запитує у користувача значення дійсних параметрів $x$ та $y$, обчислює для них значення виразу 
$$\frac{\sqrt x + 61}{(\arccos x + 0.1) (y + 37)}$$ 
та повiдомляє введенi данi та результат обчислення.

%enditem
\par
%beginitem
3. Написати функцію, що для цілого числа знаходить простий дільник, який входить у його розклад ц найбільшому степені. Якщо таких кілька, то повернути найбільший. Наприклад, для числа $2^{9}{\cdot}3^5{\cdot}7^9{\cdot}11$ функція має повернути $7$.

По завданням 2 та 3 оцінюється правильність та якість коду, а також економність обчислень.

%enditem
}
{\smallskip\smallskip \begin {scriptsize} Затверджено на засіданні кафедри теоретичної кібернетики, протокол \No 2  від   07.10.2024 р.\par\smallskip\smallskip\smallskip\smallskip
{\parbox {6.5 cm} {Зав. кафедри \dotfill Ю.В.~Крак}}\hspace {0.7cm} {\hbox to 6.5 cm { Екзаменатор \dotfill Т.О.~Карнаух}} \end{scriptsize}}
%end
\vspace{0.9cm}
%begin
{{\begin {scriptsize}\par
{ \center  Київський національний університет імені Тараса Шевченка \par}
{\parbox {5 cm}{Освітня програма \hfill}{\parbox {7 cm}{Інформатика\hfill}}\parbox {6 cm}{\hfill Семестр 1}}\par
{\parbox {5 cm} {Навчальний предмет \hfill}\parbox {7 cm}{Програмування \hfill}\parbox {6cm}{\hfill}\par}\vspace{-15pt} \end{scriptsize}}{\center Екзаменаційний білет \No \textbf { 101 }\par}}
{%beginitem
1. Змінні, їх властивості (у тому числі час існування, область дії оголошення імені) та операції над ними.%enditem
\par
%beginitem
2. Написати програму, що запитує у користувача значення дійсних параметрів $x$ та $y$, обчислює для них значення виразу 
$$\frac{\pi x - 49}{(\tg x - 0.4) (y + 94)}$$ 
та повiдомляє введенi данi та результат обчислення.

%enditem
\par
%beginitem
3. Написати функцію, що знаходить найбільше число з послідовності цілих (задається масивом), у десятковому запису якого нема цифри 3. Повертає індекс останнього входження цього числа або $-1$ (за відсутності). Використовувати додаткові структури даних (у тому числі рядки) заборонено.

По завданням 2 та 3 оцінюється правильність та якість коду, а також економність обчислень.

%enditem
}
{\smallskip\smallskip \begin {scriptsize} Затверджено на засіданні кафедри теоретичної кібернетики, протокол \No 2  від   07.10.2024 р.\par\smallskip\smallskip\smallskip\smallskip
{\parbox {6.5 cm} {Зав. кафедри \dotfill Ю.В.~Крак}}\hspace {0.7cm} {\hbox to 6.5 cm { Екзаменатор \dotfill Т.О.~Карнаух}} \end{scriptsize}}
%end
\vspace{0.9cm}
%begin
{{\begin {scriptsize}\par
{ \center  Київський національний університет імені Тараса Шевченка \par}
{\parbox {5 cm}{Освітня програма \hfill}{\parbox {7 cm}{Інформатика\hfill}}\parbox {6 cm}{\hfill Семестр 1}}\par
{\parbox {5 cm} {Навчальний предмет \hfill}\parbox {7 cm}{Програмування \hfill}\parbox {6cm}{\hfill}\par}\vspace{-15pt} \end{scriptsize}}{\center Екзаменаційний білет \No \textbf { 102 }\par}}
{%beginitem
1. Глибина рекурсії та загальна кількість рекурсивних викликів, їх вплив на розміри програмного стеку та на час виконання програми.%enditem
\par
%beginitem
2. Написати програму, що запитує у користувача значення дійсних параметрів $x$ та $y$, обчислює для них значення виразу 
$$\frac{\ x + 87}{(\sin x + 0.5) (y - 85)}$$ 
та повiдомляє введенi данi та результат обчислення.

%enditem
\par
%beginitem
3. Написати функцію, що для інтервалу $[a, b)$ знаходить кількість чисел, квадрати яких діляться на суму квадратів їх цифр. (Розглядаємо десятковий запис.)

По завданням 2 та 3 оцінюється правильність та якість коду, а також економність обчислень.

%enditem
}
{\smallskip\smallskip \begin {scriptsize} Затверджено на засіданні кафедри теоретичної кібернетики, протокол \No 2  від   07.10.2024 р.\par\smallskip\smallskip\smallskip\smallskip
{\parbox {6.5 cm} {Зав. кафедри \dotfill Ю.В.~Крак}}\hspace {0.7cm} {\hbox to 6.5 cm { Екзаменатор \dotfill Т.О.~Карнаух}} \end{scriptsize}}
%end
\newpage
%begin
{{\begin {scriptsize}\par
{ \center  Київський національний університет імені Тараса Шевченка \par}
{\parbox {5 cm}{Освітня програма \hfill}{\parbox {7 cm}{Інформатика\hfill}}\parbox {6 cm}{\hfill Семестр 1}}\par
{\parbox {5 cm} {Навчальний предмет \hfill}\parbox {7 cm}{Програмування \hfill}\parbox {6cm}{\hfill}\par}\vspace{-15pt} \end{scriptsize}}{\center Екзаменаційний білет \No \textbf { 103 }\par}}
{%beginitem
1. Принципи зображення цілих та дійсних чисел у комп'ютері.%enditem
\par
%beginitem
2. Написати програму, що запитує у користувача значення дійсних параметрів $x$ та $y$, обчислює для них значення виразу 
$$\frac{\pi x + 91}{(\arcsin x + 0.7) (y + 56)}$$ 
та повiдомляє введенi данi та результат обчислення.

%enditem
\par
%beginitem
3. Написати функцію, що повертає найбільше число з послідовності цілих (задається масивом), що має рівно три різних простих дільники. За відсутності має повертати None. Використовувати додаткові структури даних (у тому числі рядки) заборонено.

По завданням 2 та 3 оцінюється правильність та якість коду, а також економність обчислень.

%enditem
}
{\smallskip\smallskip \begin {scriptsize} Затверджено на засіданні кафедри теоретичної кібернетики, протокол \No 2  від   07.10.2024 р.\par\smallskip\smallskip\smallskip\smallskip
{\parbox {6.5 cm} {Зав. кафедри \dotfill Ю.В.~Крак}}\hspace {0.7cm} {\hbox to 6.5 cm { Екзаменатор \dotfill Т.О.~Карнаух}} \end{scriptsize}}
%end
\vspace{0.9cm}
%begin
{{\begin {scriptsize}\par
{ \center  Київський національний університет імені Тараса Шевченка \par}
{\parbox {5 cm}{Освітня програма \hfill}{\parbox {7 cm}{Інформатика\hfill}}\parbox {6 cm}{\hfill Семестр 1}}\par
{\parbox {5 cm} {Навчальний предмет \hfill}\parbox {7 cm}{Програмування \hfill}\parbox {6cm}{\hfill}\par}\vspace{-15pt} \end{scriptsize}}{\center Екзаменаційний білет \No \textbf { 104 }\par}}
{%beginitem
1. Рекурсивні означення, рекурсивні функції.%enditem
\par
%beginitem
2. Написати програму, що запитує у користувача значення дійсних параметрів $x$ та $y$, обчислює для них значення виразу 
$$\frac{\sqrt x - 44}{(\arctg x - 0.3) (y - 17)}$$ 
та повiдомляє введенi данi та результат обчислення.

%enditem
\par
%beginitem
3. Написати функцію, що повертає число з послідовності цілих (задається масивом), що має найменшу кількість різних простих дільників. Якщо таких чисел кілька, то повернути найбільше. Використовувати додаткові структури даних (у тому числі рядки) заборонено.

По завданням 2 та 3 оцінюється правильність та якість коду, а також економність обчислень.

%enditem
}
{\smallskip\smallskip \begin {scriptsize} Затверджено на засіданні кафедри теоретичної кібернетики, протокол \No 2  від   07.10.2024 р.\par\smallskip\smallskip\smallskip\smallskip
{\parbox {6.5 cm} {Зав. кафедри \dotfill Ю.В.~Крак}}\hspace {0.7cm} {\hbox to 6.5 cm { Екзаменатор \dotfill Т.О.~Карнаух}} \end{scriptsize}}
%end
\vspace{0.9cm}
%begin
{{\begin {scriptsize}\par
{ \center  Київський національний університет імені Тараса Шевченка \par}
{\parbox {5 cm}{Освітня програма \hfill}{\parbox {7 cm}{Інформатика\hfill}}\parbox {6 cm}{\hfill Семестр 1}}\par
{\parbox {5 cm} {Навчальний предмет \hfill}\parbox {7 cm}{Програмування \hfill}\parbox {6cm}{\hfill}\par}\vspace{-15pt} \end{scriptsize}}{\center Екзаменаційний білет \No \textbf { 105 }\par}}
{%beginitem
1. Арифметичні операції. Вирази. Пріоритети операторів. Порядок обчислення виразів.%enditem
\par
%beginitem
2. Написати програму, що запитує у користувача значення дійсних параметрів $x$ та $y$, обчислює для них значення виразу 
$$\frac{\log_{2} x + 27}{(\cos x + 0.8) (y + 39)}$$ 
та повiдомляє введенi данi та результат обчислення.

%enditem
\par
%beginitem
3. Написати функцію, що для інтервалу $[a, b)$ знаходить найближчу пару чисел з цього інтервалу, що кожне з них має непарну кількість різних простих дільників. Якщо таких пар кілька, то цікавить пара з найменшим добутком.

По завданням 2 та 3 оцінюється правильність та якість коду, а також економність обчислень.

%enditem
}
{\smallskip\smallskip \begin {scriptsize} Затверджено на засіданні кафедри теоретичної кібернетики, протокол \No 2  від   07.10.2024 р.\par\smallskip\smallskip\smallskip\smallskip
{\parbox {6.5 cm} {Зав. кафедри \dotfill Ю.В.~Крак}}\hspace {0.7cm} {\hbox to 6.5 cm { Екзаменатор \dotfill Т.О.~Карнаух}} \end{scriptsize}}
%end
\newpage
%begin
{{\begin {scriptsize}\par
{ \center  Київський національний університет імені Тараса Шевченка \par}
{\parbox {5 cm}{Освітня програма \hfill}{\parbox {7 cm}{Інформатика\hfill}}\parbox {6 cm}{\hfill Семестр 1}}\par
{\parbox {5 cm} {Навчальний предмет \hfill}\parbox {7 cm}{Програмування \hfill}\parbox {6cm}{\hfill}\par}\vspace{-15pt} \end{scriptsize}}{\center Екзаменаційний білет \No \textbf { 106 }\par}}
{%beginitem
1. Цикли while та do-while, їх вигляд і семантика, умова продовження циклу.%enditem
\par
%beginitem
2. Написати програму, що запитує у користувача значення дійсних параметрів $x$ та $y$, обчислює для них значення виразу 
$$\frac{\log_{10} x - 67}{(\arcsin x - 0.9) (y - 52)}$$ 
та повiдомляє введенi данi та результат обчислення.

%enditem
\par
%beginitem
3. Написати функцію, що повертає число з послідовності цілих (задається масивом), що має парну кількість різних простих дільників. Якщо таких чисел кілька, то повернути найбільше. Використовувати додаткові структури даних (у тому числі рядки) заборонено.

По завданням 2 та 3 оцінюється правильність та якість коду, а також економність обчислень.

%enditem
}
{\smallskip\smallskip \begin {scriptsize} Затверджено на засіданні кафедри теоретичної кібернетики, протокол \No 2  від   07.10.2024 р.\par\smallskip\smallskip\smallskip\smallskip
{\parbox {6.5 cm} {Зав. кафедри \dotfill Ю.В.~Крак}}\hspace {0.7cm} {\hbox to 6.5 cm { Екзаменатор \dotfill Т.О.~Карнаух}} \end{scriptsize}}
%end
\vspace{0.9cm}
%begin
{{\begin {scriptsize}\par
{ \center  Київський національний університет імені Тараса Шевченка \par}
{\parbox {5 cm}{Освітня програма \hfill}{\parbox {7 cm}{Інформатика\hfill}}\parbox {6 cm}{\hfill Семестр 1}}\par
{\parbox {5 cm} {Навчальний предмет \hfill}\parbox {7 cm}{Програмування \hfill}\parbox {6cm}{\hfill}\par}\vspace{-15pt} \end{scriptsize}}{\center Екзаменаційний білет \No \textbf { 107 }\par}}
{%beginitem
1. Поняття часової та просторової складності обчислень. Складність у середньому. Асимптотичні оцінки складності.%enditem
\par
%beginitem
2. Написати програму, що запитує у користувача значення дійсних параметрів $x$ та $y$, обчислює для них значення виразу 
$$\frac{\ln x - 45}{(\arccos x - 0.6) (y + 32)}$$ 
та повiдомляє введенi данi та результат обчислення.

%enditem
\par
%beginitem
3. Написати функцію, що для цілого числа знаходить суму його простих дільників, що входять у його розклад не більш ніж у 4му степеню. Наприклад, для числа $3^5{\cdot}7^9{\cdot}11{\cdot}23^{2}$ таких дільників два – $11$ та $23$, тому функція має повернути $34$.

По завданням 2 та 3 оцінюється правильність та якість коду, а також економність обчислень.

%enditem
}
{\smallskip\smallskip \begin {scriptsize} Затверджено на засіданні кафедри теоретичної кібернетики, протокол \No 2  від   07.10.2024 р.\par\smallskip\smallskip\smallskip\smallskip
{\parbox {6.5 cm} {Зав. кафедри \dotfill Ю.В.~Крак}}\hspace {0.7cm} {\hbox to 6.5 cm { Екзаменатор \dotfill Т.О.~Карнаух}} \end{scriptsize}}
%end
\vspace{0.9cm}
%begin
{{\begin {scriptsize}\par
{ \center  Київський національний університет імені Тараса Шевченка \par}
{\parbox {5 cm}{Освітня програма \hfill}{\parbox {7 cm}{Інформатика\hfill}}\parbox {6 cm}{\hfill Семестр 1}}\par
{\parbox {5 cm} {Навчальний предмет \hfill}\parbox {7 cm}{Програмування \hfill}\parbox {6cm}{\hfill}\par}\vspace{-15pt} \end{scriptsize}}{\center Екзаменаційний білет \No \textbf { 108 }\par}}
{%beginitem
1. Цикл for, вигляд і семантика інструкції циклу for.%enditem
\par
%beginitem
2. Написати програму, що запитує у користувача значення дійсних параметрів $x$ та $y$, обчислює для них значення виразу 
$$\frac{\ x + 20}{(\sin x + 0.2) (y - 13)}$$ 
та повiдомляє введенi данi та результат обчислення.

%enditem
\par
%beginitem
3. Написати функцію, що для інтервалу $[a, b)$ знаходить найближчу пару чисел з цього інтервалу, що мають від трьох до п'яти (включно) простих дільників.

По завданням 2 та 3 оцінюється правильність та якість коду, а також економність обчислень.

%enditem
}
{\smallskip\smallskip \begin {scriptsize} Затверджено на засіданні кафедри теоретичної кібернетики, протокол \No 2  від   07.10.2024 р.\par\smallskip\smallskip\smallskip\smallskip
{\parbox {6.5 cm} {Зав. кафедри \dotfill Ю.В.~Крак}}\hspace {0.7cm} {\hbox to 6.5 cm { Екзаменатор \dotfill Т.О.~Карнаух}} \end{scriptsize}}
%end
\newpage
%begin
{{\begin {scriptsize}\par
{ \center  Київський національний університет імені Тараса Шевченка \par}
{\parbox {5 cm}{Освітня програма \hfill}{\parbox {7 cm}{Інформатика\hfill}}\parbox {6 cm}{\hfill Семестр 1}}\par
{\parbox {5 cm} {Навчальний предмет \hfill}\parbox {7 cm}{Програмування \hfill}\parbox {6cm}{\hfill}\par}\vspace{-15pt} \end{scriptsize}}{\center Екзаменаційний білет \No \textbf { 109 }\par}}
{%beginitem
1. Глибина рекурсії та загальна кількість рекурсивних викликів, їх вплив на розміри програмного стеку та на час виконання програми.%enditem
\par
%beginitem
2. Написати програму, що запитує у користувача значення дійсних параметрів $x$ та $y$, обчислює для них значення виразу 
$$\frac{\log_{10} x - 79}{(\cos x + 0.6) (y + 68)}$$ 
та повiдомляє введенi данi та результат обчислення.

%enditem
\par
%beginitem
3. Написати функцію, що знаходить найбільше число з послідовності цілих (задається масивом), у десятковому запису якого нема цифри 3. Повертає індекс останнього входження цього числа або $-1$ (за відсутності). Використовувати додаткові структури даних (у тому числі рядки) заборонено.

По завданням 2 та 3 оцінюється правильність та якість коду, а також економність обчислень.

%enditem
}
{\smallskip\smallskip \begin {scriptsize} Затверджено на засіданні кафедри теоретичної кібернетики, протокол \No 2  від   07.10.2024 р.\par\smallskip\smallskip\smallskip\smallskip
{\parbox {6.5 cm} {Зав. кафедри \dotfill Ю.В.~Крак}}\hspace {0.7cm} {\hbox to 6.5 cm { Екзаменатор \dotfill Т.О.~Карнаух}} \end{scriptsize}}
%end
\vspace{0.9cm}
%begin
{{\begin {scriptsize}\par
{ \center  Київський національний університет імені Тараса Шевченка \par}
{\parbox {5 cm}{Освітня програма \hfill}{\parbox {7 cm}{Інформатика\hfill}}\parbox {6 cm}{\hfill Семестр 1}}\par
{\parbox {5 cm} {Навчальний предмет \hfill}\parbox {7 cm}{Програмування \hfill}\parbox {6cm}{\hfill}\par}\vspace{-15pt} \end{scriptsize}}{\center Екзаменаційний білет \No \textbf { 110 }\par}}
{%beginitem
1. Математичний підхід до визначення складності обчислень.%enditem
\par
%beginitem
2. Написати програму, що запитує у користувача значення дійсних параметрів $x$ та $y$, обчислює для них значення виразу 
$$\frac{\log_{2} x + 57}{(\ctg x - 0.2) (y - 66)}$$ 
та повiдомляє введенi данi та результат обчислення.

%enditem
\par
%beginitem
3. Написати функцію, що для цілого числа знаходить кількість простих дільників, які входять у його розклад у найбільшому степені. Наприклад, для числа $3^5{\cdot}7^9{\cdot}11{\cdot}23^{9}$ таких дільників два – $23$ та $7$, тому функція має повернути $2$.

По завданням 2 та 3 оцінюється правильність та якість коду, а також економність обчислень.

%enditem
}
{\smallskip\smallskip \begin {scriptsize} Затверджено на засіданні кафедри теоретичної кібернетики, протокол \No 2  від   07.10.2024 р.\par\smallskip\smallskip\smallskip\smallskip
{\parbox {6.5 cm} {Зав. кафедри \dotfill Ю.В.~Крак}}\hspace {0.7cm} {\hbox to 6.5 cm { Екзаменатор \dotfill Т.О.~Карнаух}} \end{scriptsize}}
%end
\vspace{0.9cm}
%begin
{{\begin {scriptsize}\par
{ \center  Київський національний університет імені Тараса Шевченка \par}
{\parbox {5 cm}{Освітня програма \hfill}{\parbox {7 cm}{Інформатика\hfill}}\parbox {6 cm}{\hfill Семестр 1}}\par
{\parbox {5 cm} {Навчальний предмет \hfill}\parbox {7 cm}{Програмування \hfill}\parbox {6cm}{\hfill}\par}\vspace{-15pt} \end{scriptsize}}{\center Екзаменаційний білет \No \textbf { 111 }\par}}
{%beginitem
1. Рекурсивні означення, рекурсивні функції.%enditem
\par
%beginitem
2. Написати програму, що запитує у користувача значення дійсних параметрів $x$ та $y$, обчислює для них значення виразу 
$$\frac{\ln x + 75}{(\arctg x - 0.1) (y + 89)}$$ 
та повiдомляє введенi данi та результат обчислення.

%enditem
\par
%beginitem
3. Написати функцію, яка для цілого числа знаходить кількість його різних простих дільників, що мають у десятковому запису цифру 1.

По завданням 2 та 3 оцінюється правильність та якість коду, а також економність обчислень.

%enditem
}
{\smallskip\smallskip \begin {scriptsize} Затверджено на засіданні кафедри теоретичної кібернетики, протокол \No 2  від   07.10.2024 р.\par\smallskip\smallskip\smallskip\smallskip
{\parbox {6.5 cm} {Зав. кафедри \dotfill Ю.В.~Крак}}\hspace {0.7cm} {\hbox to 6.5 cm { Екзаменатор \dotfill Т.О.~Карнаух}} \end{scriptsize}}
%end
\newpage
%begin
{{\begin {scriptsize}\par
{ \center  Київський національний університет імені Тараса Шевченка \par}
{\parbox {5 cm}{Освітня програма \hfill}{\parbox {7 cm}{Інформатика\hfill}}\parbox {6 cm}{\hfill Семестр 1}}\par
{\parbox {5 cm} {Навчальний предмет \hfill}\parbox {7 cm}{Програмування \hfill}\parbox {6cm}{\hfill}\par}\vspace{-15pt} \end{scriptsize}}{\center Екзаменаційний білет \No \textbf { 112 }\par}}
{%beginitem
1. Цикл for, вигляд і семантика інструкції циклу for.%enditem
\par
%beginitem
2. Написати програму, що запитує у користувача значення дійсних параметрів $x$ та $y$, обчислює для них значення виразу 
$$\frac{\sqrt x - 15}{(\tg x + 0.7) (y - 53)}$$ 
та повiдомляє введенi данi та результат обчислення.

%enditem
\par
%beginitem
3. Написати функцію, що для інтервалу $[a, b)$ знаходить кількість чисел, квадрати яких діляться на суму квадратів їх цифр. (Розглядаємо десятковий запис.)

По завданням 2 та 3 оцінюється правильність та якість коду, а також економність обчислень.

%enditem
}
{\smallskip\smallskip \begin {scriptsize} Затверджено на засіданні кафедри теоретичної кібернетики, протокол \No 2  від   07.10.2024 р.\par\smallskip\smallskip\smallskip\smallskip
{\parbox {6.5 cm} {Зав. кафедри \dotfill Ю.В.~Крак}}\hspace {0.7cm} {\hbox to 6.5 cm { Екзаменатор \dotfill Т.О.~Карнаух}} \end{scriptsize}}
%end
\vspace{0.9cm}
%begin
{{\begin {scriptsize}\par
{ \center  Київський національний університет імені Тараса Шевченка \par}
{\parbox {5 cm}{Освітня програма \hfill}{\parbox {7 cm}{Інформатика\hfill}}\parbox {6 cm}{\hfill Семестр 1}}\par
{\parbox {5 cm} {Навчальний предмет \hfill}\parbox {7 cm}{Програмування \hfill}\parbox {6cm}{\hfill}\par}\vspace{-15pt} \end{scriptsize}}{\center Екзаменаційний білет \No \textbf { 113 }\par}}
{%beginitem
1. Алгоритм двійкового пошуку та його складність.%enditem
\par
%beginitem
2. Написати програму, що запитує у користувача значення дійсних параметрів $x$ та $y$, обчислює для них значення виразу 
$$\frac{\pi x - 98}{(\sin x + 0.5) (y + 36)}$$ 
та повiдомляє введенi данi та результат обчислення.

%enditem
\par
%beginitem
3. Написати функцію, що для інтервалу $[a, b)$ знаходить найближчу пару чисел з цього інтервалу, що кожне з них має непарну кількість різних простих дільників. Якщо таких пар кілька, то цікавить пара з найбільшою сумою.

По завданням 2 та 3 оцінюється правильність та якість коду, а також економність обчислень.

%enditem
}
{\smallskip\smallskip \begin {scriptsize} Затверджено на засіданні кафедри теоретичної кібернетики, протокол \No 2  від   07.10.2024 р.\par\smallskip\smallskip\smallskip\smallskip
{\parbox {6.5 cm} {Зав. кафедри \dotfill Ю.В.~Крак}}\hspace {0.7cm} {\hbox to 6.5 cm { Екзаменатор \dotfill Т.О.~Карнаух}} \end{scriptsize}}
%end
\vspace{0.9cm}
%begin
{{\begin {scriptsize}\par
{ \center  Київський національний університет імені Тараса Шевченка \par}
{\parbox {5 cm}{Освітня програма \hfill}{\parbox {7 cm}{Інформатика\hfill}}\parbox {6 cm}{\hfill Семестр 1}}\par
{\parbox {5 cm} {Навчальний предмет \hfill}\parbox {7 cm}{Програмування \hfill}\parbox {6cm}{\hfill}\par}\vspace{-15pt} \end{scriptsize}}{\center Екзаменаційний білет \No \textbf { 114 }\par}}
{%beginitem
1. Принципи зображення цілих та дійсних чисел у комп'ютері.%enditem
\par
%beginitem
2. Написати програму, що запитує у користувача значення дійсних параметрів $x$ та $y$, обчислює для них значення виразу 
$$\frac{\ x + 86}{(\arctg x - 0.4) (y - 59)}$$ 
та повiдомляє введенi данi та результат обчислення.

%enditem
\par
%beginitem
3. Написати функцію, що для послідовності цілих (задається масивом) знаходить найдовший відрізок, в якому всі сусідні числа різні. Якщо таких відрізків кілька, то повернути останній (як індекс початку та довжину). Використовувати додаткові структури даних (у тому числі рядки) заборонено.

По завданням 2 та 3 оцінюється правильність та якість коду, а також економність обчислень.

%enditem
}
{\smallskip\smallskip \begin {scriptsize} Затверджено на засіданні кафедри теоретичної кібернетики, протокол \No 2  від   07.10.2024 р.\par\smallskip\smallskip\smallskip\smallskip
{\parbox {6.5 cm} {Зав. кафедри \dotfill Ю.В.~Крак}}\hspace {0.7cm} {\hbox to 6.5 cm { Екзаменатор \dotfill Т.О.~Карнаух}} \end{scriptsize}}
%end
\newpage
%begin
{{\begin {scriptsize}\par
{ \center  Київський національний університет імені Тараса Шевченка \par}
{\parbox {5 cm}{Освітня програма \hfill}{\parbox {7 cm}{Інформатика\hfill}}\parbox {6 cm}{\hfill Семестр 1}}\par
{\parbox {5 cm} {Навчальний предмет \hfill}\parbox {7 cm}{Програмування \hfill}\parbox {6cm}{\hfill}\par}\vspace{-15pt} \end{scriptsize}}{\center Екзаменаційний білет \No \textbf { 115 }\par}}
{%beginitem
1. Винятки та їх оброблення.%enditem
\par
%beginitem
2. Написати програму, що запитує у користувача значення дійсних параметрів $x$ та $y$, обчислює для них значення виразу 
$$\frac{\ln x + 78}{(\cos x - 0.9) (y - 51)}$$ 
та повiдомляє введенi данi та результат обчислення.

%enditem
\par
%beginitem
3. Написати функцію, що для інтервалу $[a, b)$ знаходить знаходить найбільший степінь числа 3, на який ділиться хоча б одне ціле число з цього інтервалу.

По завданням 2 та 3 оцінюється правильність та якість коду, а також економність обчислень.

%enditem
}
{\smallskip\smallskip \begin {scriptsize} Затверджено на засіданні кафедри теоретичної кібернетики, протокол \No 2  від   07.10.2024 р.\par\smallskip\smallskip\smallskip\smallskip
{\parbox {6.5 cm} {Зав. кафедри \dotfill Ю.В.~Крак}}\hspace {0.7cm} {\hbox to 6.5 cm { Екзаменатор \dotfill Т.О.~Карнаух}} \end{scriptsize}}
%end
\vspace{0.9cm}
%begin
{{\begin {scriptsize}\par
{ \center  Київський національний університет імені Тараса Шевченка \par}
{\parbox {5 cm}{Освітня програма \hfill}{\parbox {7 cm}{Інформатика\hfill}}\parbox {6 cm}{\hfill Семестр 1}}\par
{\parbox {5 cm} {Навчальний предмет \hfill}\parbox {7 cm}{Програмування \hfill}\parbox {6cm}{\hfill}\par}\vspace{-15pt} \end{scriptsize}}{\center Екзаменаційний білет \No \textbf { 116 }\par}}
{%beginitem
1. Емпіричний підхід до визначення складності обчислень.%enditem
\par
%beginitem
2. Написати програму, що запитує у користувача значення дійсних параметрів $x$ та $y$, обчислює для них значення виразу 
$$\frac{\log_{10} x - 46}{(\arcsin x + 0.3) (y + 34)}$$ 
та повiдомляє введенi данi та результат обчислення.

%enditem
\par
%beginitem
3. Написати функцію, що для цілого числа знаходить простий дільник, який входить у його розклад у найбільшому степені. Якщо таких кілька, то повернути найменший. Наприклад, для числа $2^{9}{\cdot}3^5{\cdot}7^9{\cdot}11$ функція має повернути $2$.

По завданням 2 та 3 оцінюється правильність та якість коду, а також економність обчислень.

%enditem
}
{\smallskip\smallskip \begin {scriptsize} Затверджено на засіданні кафедри теоретичної кібернетики, протокол \No 2  від   07.10.2024 р.\par\smallskip\smallskip\smallskip\smallskip
{\parbox {6.5 cm} {Зав. кафедри \dotfill Ю.В.~Крак}}\hspace {0.7cm} {\hbox to 6.5 cm { Екзаменатор \dotfill Т.О.~Карнаух}} \end{scriptsize}}
%end
\vspace{0.9cm}
%begin
{{\begin {scriptsize}\par
{ \center  Київський національний університет імені Тараса Шевченка \par}
{\parbox {5 cm}{Освітня програма \hfill}{\parbox {7 cm}{Інформатика\hfill}}\parbox {6 cm}{\hfill Семестр 1}}\par
{\parbox {5 cm} {Навчальний предмет \hfill}\parbox {7 cm}{Програмування \hfill}\parbox {6cm}{\hfill}\par}\vspace{-15pt} \end{scriptsize}}{\center Екзаменаційний білет \No \textbf { 117 }\par}}
{%beginitem
1. Оператори порівняння (relational), рівності (equality) та булеві оператори.%enditem
\par
%beginitem
2. Написати програму, що запитує у користувача значення дійсних параметрів $x$ та $y$, обчислює для них значення виразу 
$$\frac{\sqrt x - 31}{(\arccos x - 0.8) (y - 19)}$$ 
та повiдомляє введенi данi та результат обчислення.

%enditem
\par
%beginitem
3. Написати функцію, що для інтервалу $[a, b)$ знаходить найближчу пару чисел з цього інтервалу, що кожне з них має парну кількість різних простих дільників. Якщо таких пар кілька, то цікавить пара з найменшим добутком.

По завданням 2 та 3 оцінюється правильність та якість коду, а також економність обчислень.

%enditem
}
{\smallskip\smallskip \begin {scriptsize} Затверджено на засіданні кафедри теоретичної кібернетики, протокол \No 2  від   07.10.2024 р.\par\smallskip\smallskip\smallskip\smallskip
{\parbox {6.5 cm} {Зав. кафедри \dotfill Ю.В.~Крак}}\hspace {0.7cm} {\hbox to 6.5 cm { Екзаменатор \dotfill Т.О.~Карнаух}} \end{scriptsize}}
%end
\newpage
%begin
{{\begin {scriptsize}\par
{ \center  Київський національний університет імені Тараса Шевченка \par}
{\parbox {5 cm}{Освітня програма \hfill}{\parbox {7 cm}{Інформатика\hfill}}\parbox {6 cm}{\hfill Семестр 1}}\par
{\parbox {5 cm} {Навчальний предмет \hfill}\parbox {7 cm}{Програмування \hfill}\parbox {6cm}{\hfill}\par}\vspace{-15pt} \end{scriptsize}}{\center Екзаменаційний білет \No \textbf { 118 }\par}}
{%beginitem
1. Цикли while та do-while, їх вигляд і семантика, умова продовження циклу.%enditem
\par
%beginitem
2. Написати програму, що запитує у користувача значення дійсних параметрів $x$ та $y$, обчислює для них значення виразу 
$$\frac{\log_{2} x + 11}{(\tg x + 0.4) (y + 50)}$$ 
та повiдомляє введенi данi та результат обчислення.

%enditem
\par
%beginitem
3. Написати функцію, що повертає число з послідовності цілих (задається масивом), що має найбільшу кількість різних простих дільників. Якщо таких чисел кілька, то повернути найменше. Використовувати додаткові структури даних (у тому числі рядки) заборонено.

По завданням 2 та 3 оцінюється правильність та якість коду, а також економність обчислень.

%enditem
}
{\smallskip\smallskip \begin {scriptsize} Затверджено на засіданні кафедри теоретичної кібернетики, протокол \No 2  від   07.10.2024 р.\par\smallskip\smallskip\smallskip\smallskip
{\parbox {6.5 cm} {Зав. кафедри \dotfill Ю.В.~Крак}}\hspace {0.7cm} {\hbox to 6.5 cm { Екзаменатор \dotfill Т.О.~Карнаух}} \end{scriptsize}}
%end
\vspace{0.9cm}
%begin
{{\begin {scriptsize}\par
{ \center  Київський національний університет імені Тараса Шевченка \par}
{\parbox {5 cm}{Освітня програма \hfill}{\parbox {7 cm}{Інформатика\hfill}}\parbox {6 cm}{\hfill Семестр 1}}\par
{\parbox {5 cm} {Навчальний предмет \hfill}\parbox {7 cm}{Програмування \hfill}\parbox {6cm}{\hfill}\par}\vspace{-15pt} \end{scriptsize}}{\center Екзаменаційний білет \No \textbf { 119 }\par}}
{%beginitem
1. Арифметичні операції. Вирази. Пріоритети операторів. Порядок обчислення виразів.%enditem
\par
%beginitem
2. Написати програму, що запитує у користувача значення дійсних параметрів $x$ та $y$, обчислює для них значення виразу 
$$\frac{\ x - 12}{(\ctg x + 0.6) (y + 30)}$$ 
та повiдомляє введенi данi та результат обчислення.

%enditem
\par
%beginitem
3. Написати функцію, що знаходить найменше число з послідовності цілих (задається масивом), у десятковому запису якого нема цифри 5. Повертає індекс останнього входження цього числа або $-1$ (за відсутності). Використовувати додаткові структури даних (у тому числі рядки) заборонено.

По завданням 2 та 3 оцінюється правильність та якість коду, а також економність обчислень.

%enditem
}
{\smallskip\smallskip \begin {scriptsize} Затверджено на засіданні кафедри теоретичної кібернетики, протокол \No 2  від   07.10.2024 р.\par\smallskip\smallskip\smallskip\smallskip
{\parbox {6.5 cm} {Зав. кафедри \dotfill Ю.В.~Крак}}\hspace {0.7cm} {\hbox to 6.5 cm { Екзаменатор \dotfill Т.О.~Карнаух}} \end{scriptsize}}
%end
\vspace{0.9cm}
%begin
{{\begin {scriptsize}\par
{ \center  Київський національний університет імені Тараса Шевченка \par}
{\parbox {5 cm}{Освітня програма \hfill}{\parbox {7 cm}{Інформатика\hfill}}\parbox {6 cm}{\hfill Семестр 1}}\par
{\parbox {5 cm} {Навчальний предмет \hfill}\parbox {7 cm}{Програмування \hfill}\parbox {6cm}{\hfill}\par}\vspace{-15pt} \end{scriptsize}}{\center Екзаменаційний білет \No \textbf { 120 }\par}}
{%beginitem
1. Компільовані та інтерпретовані мови. Переваги та недоліки.%enditem
\par
%beginitem
2. Написати програму, що запитує у користувача значення дійсних параметрів $x$ та $y$, обчислює для них значення виразу 
$$\frac{\pi x + 97}{(\ctg x - 0.3) (y - 63)}$$ 
та повiдомляє введенi данi та результат обчислення.

%enditem
\par
%beginitem
3. Написати функцію, що повертає найбільше число з послідовності цілих (задається масивом), що має не більше трьох різних простих дільників. За відсутності має повертати None. Використовувати додаткові структури даних (у тому числі рядки) заборонено.

По завданням 2 та 3 оцінюється правильність та якість коду, а також економність обчислень.

%enditem
}
{\smallskip\smallskip \begin {scriptsize} Затверджено на засіданні кафедри теоретичної кібернетики, протокол \No 2  від   07.10.2024 р.\par\smallskip\smallskip\smallskip\smallskip
{\parbox {6.5 cm} {Зав. кафедри \dotfill Ю.В.~Крак}}\hspace {0.7cm} {\hbox to 6.5 cm { Екзаменатор \dotfill Т.О.~Карнаух}} \end{scriptsize}}
%end
\newpage
%begin
{{\begin {scriptsize}\par
{ \center  Київський національний університет імені Тараса Шевченка \par}
{\parbox {5 cm}{Освітня програма \hfill}{\parbox {7 cm}{Інформатика\hfill}}\parbox {6 cm}{\hfill Семестр 1}}\par
{\parbox {5 cm} {Навчальний предмет \hfill}\parbox {7 cm}{Програмування \hfill}\parbox {6cm}{\hfill}\par}\vspace{-15pt} \end{scriptsize}}{\center Екзаменаційний білет \No \textbf { 121 }\par}}
{%beginitem
1. Засоби керування порядком обчислень. Інструкція if.%enditem
\par
%beginitem
2. Написати програму, що запитує у користувача значення дійсних параметрів $x$ та $y$, обчислює для них значення виразу 
$$\frac{\ x - 69}{(\arctg x - 0.1) (y + 55)}$$ 
та повiдомляє введенi данi та результат обчислення.

%enditem
\par
%beginitem
3. Написати функцію, що для цілого числа знаходить кількість його простих дільників, що входять у його розклад не більш ніж у 5му степеню. Наприклад, для числа $3^8{\cdot}7^9{\cdot}11{\cdot}23^{2}$ таких дільників два – $11$ та $23$, тому функція має повернути $2$.

По завданням 2 та 3 оцінюється правильність та якість коду, а також економність обчислень.

%enditem
}
{\smallskip\smallskip \begin {scriptsize} Затверджено на засіданні кафедри теоретичної кібернетики, протокол \No 2  від   07.10.2024 р.\par\smallskip\smallskip\smallskip\smallskip
{\parbox {6.5 cm} {Зав. кафедри \dotfill Ю.В.~Крак}}\hspace {0.7cm} {\hbox to 6.5 cm { Екзаменатор \dotfill Т.О.~Карнаух}} \end{scriptsize}}
%end
\vspace{0.9cm}
%begin
{{\begin {scriptsize}\par
{ \center  Київський національний університет імені Тараса Шевченка \par}
{\parbox {5 cm}{Освітня програма \hfill}{\parbox {7 cm}{Інформатика\hfill}}\parbox {6 cm}{\hfill Семестр 1}}\par
{\parbox {5 cm} {Навчальний предмет \hfill}\parbox {7 cm}{Програмування \hfill}\parbox {6cm}{\hfill}\par}\vspace{-15pt} \end{scriptsize}}{\center Екзаменаційний білет \No \textbf { 122 }\par}}
{%beginitem
1. Поняття типу даних. Вбудовані арифметичні та цілі (intergral) типи мови С++.%enditem
\par
%beginitem
2. Написати програму, що запитує у користувача значення дійсних параметрів $x$ та $y$, обчислює для них значення виразу 
$$\frac{\log_{2} x + 21}{(\tg x + 0.8) (y - 35)}$$ 
та повiдомляє введенi данi та результат обчислення.

%enditem
\par
%beginitem
3. Написати функцію, що для цілого числа знаходить кількість його простих дільників, що входять у його розклад рівно у другому степеню. Наприклад, для числа $3^2{\cdot}7^9{\cdot}11{\cdot}23^{2}$ таких дільників два – $3$ та $23$, тому функція має повернути $2$.

По завданням 2 та 3 оцінюється правильність та якість коду, а також економність обчислень.

%enditem
}
{\smallskip\smallskip \begin {scriptsize} Затверджено на засіданні кафедри теоретичної кібернетики, протокол \No 2  від   07.10.2024 р.\par\smallskip\smallskip\smallskip\smallskip
{\parbox {6.5 cm} {Зав. кафедри \dotfill Ю.В.~Крак}}\hspace {0.7cm} {\hbox to 6.5 cm { Екзаменатор \dotfill Т.О.~Карнаух}} \end{scriptsize}}
%end
\vspace{0.9cm}
%begin
{{\begin {scriptsize}\par
{ \center  Київський національний університет імені Тараса Шевченка \par}
{\parbox {5 cm}{Освітня програма \hfill}{\parbox {7 cm}{Інформатика\hfill}}\parbox {6 cm}{\hfill Семестр 1}}\par
{\parbox {5 cm} {Навчальний предмет \hfill}\parbox {7 cm}{Програмування \hfill}\parbox {6cm}{\hfill}\par}\vspace{-15pt} \end{scriptsize}}{\center Екзаменаційний білет \No \textbf { 123 }\par}}
{%beginitem
1. Поняття масиву. Моделювання послідовностей масивами. Параметри, що зображують масиви у функціях.%enditem
\par
%beginitem
2. Написати програму, що запитує у користувача значення дійсних параметрів $x$ та $y$, обчислює для них значення виразу 
$$\frac{\ln x - 72}{(\arccos x + 0.7) (y + 99)}$$ 
та повiдомляє введенi данi та результат обчислення.

%enditem
\par
%beginitem
3. Написати функцію, що для послідовності цілих (задається масивом) знаходить довжину найдовшого відрізку, в якому всі сусідні числа відрізняються якнайменш на 2. Використовувати додаткові структури даних (у тому числі рядки) заборонено.

По завданням 2 та 3 оцінюється правильність та якість коду, а також економність обчислень.

%enditem
}
{\smallskip\smallskip \begin {scriptsize} Затверджено на засіданні кафедри теоретичної кібернетики, протокол \No 2  від   07.10.2024 р.\par\smallskip\smallskip\smallskip\smallskip
{\parbox {6.5 cm} {Зав. кафедри \dotfill Ю.В.~Крак}}\hspace {0.7cm} {\hbox to 6.5 cm { Екзаменатор \dotfill Т.О.~Карнаух}} \end{scriptsize}}
%end
\newpage
%begin
{{\begin {scriptsize}\par
{ \center  Київський національний університет імені Тараса Шевченка \par}
{\parbox {5 cm}{Освітня програма \hfill}{\parbox {7 cm}{Інформатика\hfill}}\parbox {6 cm}{\hfill Семестр 1}}\par
{\parbox {5 cm} {Навчальний предмет \hfill}\parbox {7 cm}{Програмування \hfill}\parbox {6cm}{\hfill}\par}\vspace{-15pt} \end{scriptsize}}{\center Екзаменаційний білет \No \textbf { 124 }\par}}
{%beginitem
1. Поняття часової та просторової складності обчислень. Складність у середньому. Асимптотичні оцінки складності.%enditem
\par
%beginitem
2. Написати програму, що запитує у користувача значення дійсних параметрів $x$ та $y$, обчислює для них значення виразу 
$$\frac{\pi x + 81}{(\sin x - 0.9) (y - 41)}$$ 
та повiдомляє введенi данi та результат обчислення.

%enditem
\par
%beginitem
3. Написати функцію, що для інтервалу $[a, b)$ знаходить найближчу пару чисел з цього інтервалу, що кожне з них має парну кількість різних простих дільників. Якщо таких пар кілька, то цікавить пара з найбільшою сумою.

По завданням 2 та 3 оцінюється правильність та якість коду, а також економність обчислень.

%enditem
}
{\smallskip\smallskip \begin {scriptsize} Затверджено на засіданні кафедри теоретичної кібернетики, протокол \No 2  від   07.10.2024 р.\par\smallskip\smallskip\smallskip\smallskip
{\parbox {6.5 cm} {Зав. кафедри \dotfill Ю.В.~Крак}}\hspace {0.7cm} {\hbox to 6.5 cm { Екзаменатор \dotfill Т.О.~Карнаух}} \end{scriptsize}}
%end
\vspace{0.9cm}
%begin
{{\begin {scriptsize}\par
{ \center  Київський національний університет імені Тараса Шевченка \par}
{\parbox {5 cm}{Освітня програма \hfill}{\parbox {7 cm}{Інформатика\hfill}}\parbox {6 cm}{\hfill Семестр 1}}\par
{\parbox {5 cm} {Навчальний предмет \hfill}\parbox {7 cm}{Програмування \hfill}\parbox {6cm}{\hfill}\par}\vspace{-15pt} \end{scriptsize}}{\center Екзаменаційний білет \No \textbf { 125 }\par}}
{%beginitem
1. Змінні, їх властивості (у тому числі час існування, область дії оголошення імені) та операції над ними.%enditem
\par
%beginitem
2. Написати програму, що запитує у користувача значення дійсних параметрів $x$ та $y$, обчислює для них значення виразу 
$$\frac{\log_{10} x + 28}{(\cos x - 0.5) (y - 42)}$$ 
та повiдомляє введенi данi та результат обчислення.

%enditem
\par
%beginitem
3. Написати функцію, що в інтервалі $[a, b)$ знаходить ціле число, що ділиться на найбільший степінь числа 3. Якщо таких чисел декілька, повернути найменше.

По завданням 2 та 3 оцінюється правильність та якість коду, а також економність обчислень.

%enditem
}
{\smallskip\smallskip \begin {scriptsize} Затверджено на засіданні кафедри теоретичної кібернетики, протокол \No 2  від   07.10.2024 р.\par\smallskip\smallskip\smallskip\smallskip
{\parbox {6.5 cm} {Зав. кафедри \dotfill Ю.В.~Крак}}\hspace {0.7cm} {\hbox to 6.5 cm { Екзаменатор \dotfill Т.О.~Карнаух}} \end{scriptsize}}
%end
\vspace{0.9cm}
%begin
{{\begin {scriptsize}\par
{ \center  Київський національний університет імені Тараса Шевченка \par}
{\parbox {5 cm}{Освітня програма \hfill}{\parbox {7 cm}{Інформатика\hfill}}\parbox {6 cm}{\hfill Семестр 1}}\par
{\parbox {5 cm} {Навчальний предмет \hfill}\parbox {7 cm}{Програмування \hfill}\parbox {6cm}{\hfill}\par}\vspace{-15pt} \end{scriptsize}}{\center Екзаменаційний білет \No \textbf { 126 }\par}}
{%beginitem
1. Семантика функції та її виклику, параметри та аргументи.%enditem
\par
%beginitem
2. Написати програму, що запитує у користувача значення дійсних параметрів $x$ та $y$, обчислює для них значення виразу 
$$\frac{\sqrt x - 47}{(\arcsin x + 0.2) (y + 77)}$$ 
та повiдомляє введенi данi та результат обчислення.

%enditem
\par
%beginitem
3. Написати функцію, що знаходить найменше число з послідовності цілих (задається масивом), у десятковому запису якого нема цифри 7. Повертає індекс першого входження цього числа або $-1$ (за відсутності). Використовувати додаткові структури даних (у тому числі рядки) заборонено.

По завданням 2 та 3 оцінюється правильність та якість коду, а також економність обчислень.

%enditem
}
{\smallskip\smallskip \begin {scriptsize} Затверджено на засіданні кафедри теоретичної кібернетики, протокол \No 2  від   07.10.2024 р.\par\smallskip\smallskip\smallskip\smallskip
{\parbox {6.5 cm} {Зав. кафедри \dotfill Ю.В.~Крак}}\hspace {0.7cm} {\hbox to 6.5 cm { Екзаменатор \dotfill Т.О.~Карнаух}} \end{scriptsize}}
%end
\newpage
%begin
{{\begin {scriptsize}\par
{ \center  Київський національний університет імені Тараса Шевченка \par}
{\parbox {5 cm}{Освітня програма \hfill}{\parbox {7 cm}{Інформатика\hfill}}\parbox {6 cm}{\hfill Семестр 1}}\par
{\parbox {5 cm} {Навчальний предмет \hfill}\parbox {7 cm}{Програмування \hfill}\parbox {6cm}{\hfill}\par}\vspace{-15pt} \end{scriptsize}}{\center Екзаменаційний білет \No \textbf { 127 }\par}}
{%beginitem
1. Компільовані та інтерпретовані мови. Переваги та недоліки.%enditem
\par
%beginitem
2. Написати програму, що запитує у користувача значення дійсних параметрів $x$ та $y$, обчислює для них значення виразу 
$$\frac{\ln x - 43}{(\arccos x + 0.3) (y + 70)}$$ 
та повiдомляє введенi данi та результат обчислення.

%enditem
\par
%beginitem
3. Написати функцію, що знаходить найбільше число з послідовності цілих (задається масивом), у десятковому запису якого нема цифри 4. Повертає індекс першого входження цього числа або $-1$ (за відсутності). Використовувати додаткові структури даних (у тому числі рядки) заборонено.

По завданням 2 та 3 оцінюється правильність та якість коду, а також економність обчислень.

%enditem
}
{\smallskip\smallskip \begin {scriptsize} Затверджено на засіданні кафедри теоретичної кібернетики, протокол \No 2  від   07.10.2024 р.\par\smallskip\smallskip\smallskip\smallskip
{\parbox {6.5 cm} {Зав. кафедри \dotfill Ю.В.~Крак}}\hspace {0.7cm} {\hbox to 6.5 cm { Екзаменатор \dotfill Т.О.~Карнаух}} \end{scriptsize}}
%end
\vspace{0.9cm}
%begin
{{\begin {scriptsize}\par
{ \center  Київський національний університет імені Тараса Шевченка \par}
{\parbox {5 cm}{Освітня програма \hfill}{\parbox {7 cm}{Інформатика\hfill}}\parbox {6 cm}{\hfill Семестр 1}}\par
{\parbox {5 cm} {Навчальний предмет \hfill}\parbox {7 cm}{Програмування \hfill}\parbox {6cm}{\hfill}\par}\vspace{-15pt} \end{scriptsize}}{\center Екзаменаційний білет \No \textbf { 128 }\par}}
{%beginitem
1. Поняття типу даних. Вбудовані арифметичні та цілі (intergral) типи мови С++.%enditem
\par
%beginitem
2. Написати програму, що запитує у користувача значення дійсних параметрів $x$ та $y$, обчислює для них значення виразу 
$$\frac{\ x + 80}{(\cos x - 0.7) (y - 23)}$$ 
та повiдомляє введенi данi та результат обчислення.

%enditem
\par
%beginitem
3. Написати функцію, що для послідовності цілих (задається масивом) знаходить довжину найдовшого відрізку, що складається з чисел, у сімковому запису якого нема цифри 2. За відсутності має повертати $0$. Використовувати додаткові структури даних (у тому числі рядки) заборонено.

По завданням 2 та 3 оцінюється правильність та якість коду, а також економність обчислень.

%enditem
}
{\smallskip\smallskip \begin {scriptsize} Затверджено на засіданні кафедри теоретичної кібернетики, протокол \No 2  від   07.10.2024 р.\par\smallskip\smallskip\smallskip\smallskip
{\parbox {6.5 cm} {Зав. кафедри \dotfill Ю.В.~Крак}}\hspace {0.7cm} {\hbox to 6.5 cm { Екзаменатор \dotfill Т.О.~Карнаух}} \end{scriptsize}}
%end
\vspace{0.9cm}
%begin
{{\begin {scriptsize}\par
{ \center  Київський національний університет імені Тараса Шевченка \par}
{\parbox {5 cm}{Освітня програма \hfill}{\parbox {7 cm}{Інформатика\hfill}}\parbox {6 cm}{\hfill Семестр 1}}\par
{\parbox {5 cm} {Навчальний предмет \hfill}\parbox {7 cm}{Програмування \hfill}\parbox {6cm}{\hfill}\par}\vspace{-15pt} \end{scriptsize}}{\center Екзаменаційний білет \No \textbf { 129 }\par}}
{%beginitem
1. Семантика функції та її виклику, параметри та аргументи.%enditem
\par
%beginitem
2. Написати програму, що запитує у користувача значення дійсних параметрів $x$ та $y$, обчислює для них значення виразу 
$$\frac{\pi x + 38}{(\arcsin x - 0.1) (y + 65)}$$ 
та повiдомляє введенi данi та результат обчислення.

%enditem
\par
%beginitem
3. Написати функцію, що для послідовності цілих (задається масивом) знаходить довжину найдовшого відрізку, в якому всі сусідні числа різні. Використовувати додаткові структури даних (у тому числі рядки) заборонено.

По завданням 2 та 3 оцінюється правильність та якість коду, а також економність обчислень.

%enditem
}
{\smallskip\smallskip \begin {scriptsize} Затверджено на засіданні кафедри теоретичної кібернетики, протокол \No 2  від   07.10.2024 р.\par\smallskip\smallskip\smallskip\smallskip
{\parbox {6.5 cm} {Зав. кафедри \dotfill Ю.В.~Крак}}\hspace {0.7cm} {\hbox to 6.5 cm { Екзаменатор \dotfill Т.О.~Карнаух}} \end{scriptsize}}
%end
\newpage
%begin
{{\begin {scriptsize}\par
{ \center  Київський національний університет імені Тараса Шевченка \par}
{\parbox {5 cm}{Освітня програма \hfill}{\parbox {7 cm}{Інформатика\hfill}}\parbox {6 cm}{\hfill Семестр 1}}\par
{\parbox {5 cm} {Навчальний предмет \hfill}\parbox {7 cm}{Програмування \hfill}\parbox {6cm}{\hfill}\par}\vspace{-15pt} \end{scriptsize}}{\center Екзаменаційний білет \No \textbf { 130 }\par}}
{%beginitem
1. Глибина рекурсії та загальна кількість рекурсивних викликів, їх вплив на розміри програмного стеку та на час виконання програми.%enditem
\par
%beginitem
2. Написати програму, що запитує у користувача значення дійсних параметрів $x$ та $y$, обчислює для них значення виразу 
$$\frac{\log_{10} x - 54}{(\sin x + 0.9) (y - 24)}$$ 
та повiдомляє введенi данi та результат обчислення.

%enditem
\par
%beginitem
3. Написати функцію, що для цілого числа знаходить простий дільник, який входить у його розклад ц найбільшому степені. Якщо таких кілька, то повернути найбільший. Наприклад, для числа $2^{9}{\cdot}3^5{\cdot}7^9{\cdot}11$ функція має повернути $7$.

По завданням 2 та 3 оцінюється правильність та якість коду, а також економність обчислень.

%enditem
}
{\smallskip\smallskip \begin {scriptsize} Затверджено на засіданні кафедри теоретичної кібернетики, протокол \No 2  від   07.10.2024 р.\par\smallskip\smallskip\smallskip\smallskip
{\parbox {6.5 cm} {Зав. кафедри \dotfill Ю.В.~Крак}}\hspace {0.7cm} {\hbox to 6.5 cm { Екзаменатор \dotfill Т.О.~Карнаух}} \end{scriptsize}}
%end
\vspace{0.9cm}
%begin
{{\begin {scriptsize}\par
{ \center  Київський національний університет імені Тараса Шевченка \par}
{\parbox {5 cm}{Освітня програма \hfill}{\parbox {7 cm}{Інформатика\hfill}}\parbox {6 cm}{\hfill Семестр 1}}\par
{\parbox {5 cm} {Навчальний предмет \hfill}\parbox {7 cm}{Програмування \hfill}\parbox {6cm}{\hfill}\par}\vspace{-15pt} \end{scriptsize}}{\center Екзаменаційний білет \No \textbf { 131 }\par}}
{%beginitem
1. Принципи зображення цілих та дійсних чисел у комп'ютері.%enditem
\par
%beginitem
2. Написати програму, що запитує у користувача значення дійсних параметрів $x$ та $y$, обчислює для них значення виразу 
$$\frac{\sqrt x + 33}{(\ctg x - 0.6) (y - 71)}$$ 
та повiдомляє введенi данi та результат обчислення.

%enditem
\par
%beginitem
3. Написати функцію, що повертає найбільше число з послідовності цілих (задається масивом), у сімковому запису якого нема цифри 2. За відсутності має повертати None. Використовувати додаткові структури даних (у тому числі рядки) заборонено.

По завданням 2 та 3 оцінюється правильність та якість коду, а також економність обчислень.

%enditem
}
{\smallskip\smallskip \begin {scriptsize} Затверджено на засіданні кафедри теоретичної кібернетики, протокол \No 2  від   07.10.2024 р.\par\smallskip\smallskip\smallskip\smallskip
{\parbox {6.5 cm} {Зав. кафедри \dotfill Ю.В.~Крак}}\hspace {0.7cm} {\hbox to 6.5 cm { Екзаменатор \dotfill Т.О.~Карнаух}} \end{scriptsize}}
%end
\vspace{0.9cm}
%begin
{{\begin {scriptsize}\par
{ \center  Київський національний університет імені Тараса Шевченка \par}
{\parbox {5 cm}{Освітня програма \hfill}{\parbox {7 cm}{Інформатика\hfill}}\parbox {6 cm}{\hfill Семестр 1}}\par
{\parbox {5 cm} {Навчальний предмет \hfill}\parbox {7 cm}{Програмування \hfill}\parbox {6cm}{\hfill}\par}\vspace{-15pt} \end{scriptsize}}{\center Екзаменаційний білет \No \textbf { 132 }\par}}
{%beginitem
1. Засоби керування порядком обчислень. Інструкція if.%enditem
\par
%beginitem
2. Написати програму, що запитує у користувача значення дійсних параметрів $x$ та $y$, обчислює для них значення виразу 
$$\frac{\log_{2} x - 84}{(\arctg x + 0.8) (y + 64)}$$ 
та повiдомляє введенi данi та результат обчислення.

%enditem
\par
%beginitem
3. Написати функцію, що для послідовності цілих (задається масивом) знаходить найдовший відрізок, в якому всі сусідні числа відрізняються якнайменш на 3. Якщо таких відрізків кілька, то повернути перший (як індекс початку та довжину). Використовувати додаткові структури даних (у тому числі рядки) заборонено.

По завданням 2 та 3 оцінюється правильність та якість коду, а також економність обчислень.

%enditem
}
{\smallskip\smallskip \begin {scriptsize} Затверджено на засіданні кафедри теоретичної кібернетики, протокол \No 2  від   07.10.2024 р.\par\smallskip\smallskip\smallskip\smallskip
{\parbox {6.5 cm} {Зав. кафедри \dotfill Ю.В.~Крак}}\hspace {0.7cm} {\hbox to 6.5 cm { Екзаменатор \dotfill Т.О.~Карнаух}} \end{scriptsize}}
%end
\newpage
%begin
{{\begin {scriptsize}\par
{ \center  Київський національний університет імені Тараса Шевченка \par}
{\parbox {5 cm}{Освітня програма \hfill}{\parbox {7 cm}{Інформатика\hfill}}\parbox {6 cm}{\hfill Семестр 1}}\par
{\parbox {5 cm} {Навчальний предмет \hfill}\parbox {7 cm}{Програмування \hfill}\parbox {6cm}{\hfill}\par}\vspace{-15pt} \end{scriptsize}}{\center Екзаменаційний білет \No \textbf { 133 }\par}}
{%beginitem
1. Алгоритм двійкового пошуку та його складність.%enditem
\par
%beginitem
2. Написати програму, що запитує у користувача значення дійсних параметрів $x$ та $y$, обчислює для них значення виразу 
$$\frac{\sqrt x - 96}{(\tg x + 0.2) (y + 62)}$$ 
та повiдомляє введенi данi та результат обчислення.

%enditem
\par
%beginitem
3. Написати функцію, що для інтервалу $[a, b)$ знаходить найближчу пару чисел з цього інтервалу, що мають не менше трьох простих дільників та мають цифру 3 у своєму десятковому запису.

По завданням 2 та 3 оцінюється правильність та якість коду, а також економність обчислень.

%enditem
}
{\smallskip\smallskip \begin {scriptsize} Затверджено на засіданні кафедри теоретичної кібернетики, протокол \No 2  від   07.10.2024 р.\par\smallskip\smallskip\smallskip\smallskip
{\parbox {6.5 cm} {Зав. кафедри \dotfill Ю.В.~Крак}}\hspace {0.7cm} {\hbox to 6.5 cm { Екзаменатор \dotfill Т.О.~Карнаух}} \end{scriptsize}}
%end
\vspace{0.9cm}
%begin
{{\begin {scriptsize}\par
{ \center  Київський національний університет імені Тараса Шевченка \par}
{\parbox {5 cm}{Освітня програма \hfill}{\parbox {7 cm}{Інформатика\hfill}}\parbox {6 cm}{\hfill Семестр 1}}\par
{\parbox {5 cm} {Навчальний предмет \hfill}\parbox {7 cm}{Програмування \hfill}\parbox {6cm}{\hfill}\par}\vspace{-15pt} \end{scriptsize}}{\center Екзаменаційний білет \No \textbf { 134 }\par}}
{%beginitem
1. Оператори порівняння (relational), рівності (equality) та булеві оператори.%enditem
\par
%beginitem
2. Написати програму, що запитує у користувача значення дійсних параметрів $x$ та $y$, обчислює для них значення виразу 
$$\frac{\ x + 20}{(\sin x - 0.5) (y - 88)}$$ 
та повiдомляє введенi данi та результат обчислення.

%enditem
\par
%beginitem
3. Написати функцію, що в інтервалі $[a, b)$ знаходить ціле число, що ділиться на найбільший степінь числа 3. Якщо таких чисел декілька, повернути найбільше.

По завданням 2 та 3 оцінюється правильність та якість коду, а також економність обчислень.

%enditem
}
{\smallskip\smallskip \begin {scriptsize} Затверджено на засіданні кафедри теоретичної кібернетики, протокол \No 2  від   07.10.2024 р.\par\smallskip\smallskip\smallskip\smallskip
{\parbox {6.5 cm} {Зав. кафедри \dotfill Ю.В.~Крак}}\hspace {0.7cm} {\hbox to 6.5 cm { Екзаменатор \dotfill Т.О.~Карнаух}} \end{scriptsize}}
%end
\vspace{0.9cm}
%begin
{{\begin {scriptsize}\par
{ \center  Київський національний університет імені Тараса Шевченка \par}
{\parbox {5 cm}{Освітня програма \hfill}{\parbox {7 cm}{Інформатика\hfill}}\parbox {6 cm}{\hfill Семестр 1}}\par
{\parbox {5 cm} {Навчальний предмет \hfill}\parbox {7 cm}{Програмування \hfill}\parbox {6cm}{\hfill}\par}\vspace{-15pt} \end{scriptsize}}{\center Екзаменаційний білет \No \textbf { 135 }\par}}
{%beginitem
1. Рекурсивні означення, рекурсивні функції.%enditem
\par
%beginitem
2. Написати програму, що запитує у користувача значення дійсних параметрів $x$ та $y$, обчислює для них значення виразу 
$$\frac{\log_{10} x + 68}{(\ctg x - 0.4) (y + 86)}$$ 
та повiдомляє введенi данi та результат обчислення.

%enditem
\par
%beginitem
3. Написати функцію, що для послідовності цілих (задається масивом) знаходить найдовший відрізок, в якому всі сусідні числа різні. Якщо таких відрізків кілька, то повернути перший (як індекс початку та довжину). Використовувати додаткові структури даних (у тому числі рядки) заборонено.

По завданням 2 та 3 оцінюється правильність та якість коду, а також економність обчислень.

%enditem
}
{\smallskip\smallskip \begin {scriptsize} Затверджено на засіданні кафедри теоретичної кібернетики, протокол \No 2  від   07.10.2024 р.\par\smallskip\smallskip\smallskip\smallskip
{\parbox {6.5 cm} {Зав. кафедри \dotfill Ю.В.~Крак}}\hspace {0.7cm} {\hbox to 6.5 cm { Екзаменатор \dotfill Т.О.~Карнаух}} \end{scriptsize}}
%end
\newpage
%begin
{{\begin {scriptsize}\par
{ \center  Київський національний університет імені Тараса Шевченка \par}
{\parbox {5 cm}{Освітня програма \hfill}{\parbox {7 cm}{Інформатика\hfill}}\parbox {6 cm}{\hfill Семестр 1}}\par
{\parbox {5 cm} {Навчальний предмет \hfill}\parbox {7 cm}{Програмування \hfill}\parbox {6cm}{\hfill}\par}\vspace{-15pt} \end{scriptsize}}{\center Екзаменаційний білет \No \textbf { 136 }\par}}
{%beginitem
1. Математичний підхід до визначення складності обчислень.%enditem
\par
%beginitem
2. Написати програму, що запитує у користувача значення дійсних параметрів $x$ та $y$, обчислює для них значення виразу 
$$\frac{\pi x - 83}{(\arctg x + 0.5) (y - 71)}$$ 
та повiдомляє введенi данi та результат обчислення.

%enditem
\par
%beginitem
3. Написати функцію, що повертає найменше число з послідовності цілих (задається масивом), що має не більше трьох різних простих дільників. За відсутності має повертати None. Використовувати додаткові структури даних (у тому числі рядки) заборонено.

По завданням 2 та 3 оцінюється правильність та якість коду, а також економність обчислень.

%enditem
}
{\smallskip\smallskip \begin {scriptsize} Затверджено на засіданні кафедри теоретичної кібернетики, протокол \No 2  від   07.10.2024 р.\par\smallskip\smallskip\smallskip\smallskip
{\parbox {6.5 cm} {Зав. кафедри \dotfill Ю.В.~Крак}}\hspace {0.7cm} {\hbox to 6.5 cm { Екзаменатор \dotfill Т.О.~Карнаух}} \end{scriptsize}}
%end
\vspace{0.9cm}
%begin
{{\begin {scriptsize}\par
{ \center  Київський національний університет імені Тараса Шевченка \par}
{\parbox {5 cm}{Освітня програма \hfill}{\parbox {7 cm}{Інформатика\hfill}}\parbox {6 cm}{\hfill Семестр 1}}\par
{\parbox {5 cm} {Навчальний предмет \hfill}\parbox {7 cm}{Програмування \hfill}\parbox {6cm}{\hfill}\par}\vspace{-15pt} \end{scriptsize}}{\center Екзаменаційний білет \No \textbf { 137 }\par}}
{%beginitem
1. Цикли while та do-while, їх вигляд і семантика, умова продовження циклу.%enditem
\par
%beginitem
2. Написати програму, що запитує у користувача значення дійсних параметрів $x$ та $y$, обчислює для них значення виразу 
$$\frac{\log_{2} x - 27}{(\arccos x - 0.9) (y - 58)}$$ 
та повiдомляє введенi данi та результат обчислення.

%enditem
\par
%beginitem
3. Написати функцію, що для інтервалу $[a, b)$ знаходить найближчу пару чисел з цього інтервалу, що кожне з них має парну кількість різних простих дільників. Якщо таких пар кілька, то цікавить пара з найменшим добутком.

По завданням 2 та 3 оцінюється правильність та якість коду, а також економність обчислень.

%enditem
}
{\smallskip\smallskip \begin {scriptsize} Затверджено на засіданні кафедри теоретичної кібернетики, протокол \No 2  від   07.10.2024 р.\par\smallskip\smallskip\smallskip\smallskip
{\parbox {6.5 cm} {Зав. кафедри \dotfill Ю.В.~Крак}}\hspace {0.7cm} {\hbox to 6.5 cm { Екзаменатор \dotfill Т.О.~Карнаух}} \end{scriptsize}}
%end
\vspace{0.9cm}
%begin
{{\begin {scriptsize}\par
{ \center  Київський національний університет імені Тараса Шевченка \par}
{\parbox {5 cm}{Освітня програма \hfill}{\parbox {7 cm}{Інформатика\hfill}}\parbox {6 cm}{\hfill Семестр 1}}\par
{\parbox {5 cm} {Навчальний предмет \hfill}\parbox {7 cm}{Програмування \hfill}\parbox {6cm}{\hfill}\par}\vspace{-15pt} \end{scriptsize}}{\center Екзаменаційний білет \No \textbf { 138 }\par}}
{%beginitem
1. Змінні, їх властивості (у тому числі час існування, область дії оголошення імені) та операції над ними.%enditem
\par
%beginitem
2. Написати програму, що запитує у користувача значення дійсних параметрів $x$ та $y$, обчислює для них значення виразу 
$$\frac{\ln x + 30}{(\cos x + 0.8) (y + 80)}$$ 
та повiдомляє введенi данi та результат обчислення.

%enditem
\par
%beginitem
3. Написати функцію, що для інтервалу $[a, b)$ знаходить найближчу пару чисел з цього інтервалу, що кожне з них має парну кількість різних простих дільників. Якщо таких пар кілька, то цікавить пара з найбільшою сумою.

По завданням 2 та 3 оцінюється правильність та якість коду, а також економність обчислень.

%enditem
}
{\smallskip\smallskip \begin {scriptsize} Затверджено на засіданні кафедри теоретичної кібернетики, протокол \No 2  від   07.10.2024 р.\par\smallskip\smallskip\smallskip\smallskip
{\parbox {6.5 cm} {Зав. кафедри \dotfill Ю.В.~Крак}}\hspace {0.7cm} {\hbox to 6.5 cm { Екзаменатор \dotfill Т.О.~Карнаух}} \end{scriptsize}}
%end
\newpage
\end{document}
